\documentclass[11pt]{amsart}
\usepackage[T1]{fontenc}
\usepackage{lmodern,amsmath,amssymb,mathtools,microtype,booktabs}
\usepackage[margin=1.15in]{geometry}
\usepackage{xr-hyper}
\usepackage[hidelinks]{hyperref}
\newtheorem{theorem}{Theorem}[section]
\newtheorem{proposition}[theorem]{Proposition}
\newtheorem{lemma}[theorem]{Lemma}
\newtheorem{corollary}[theorem]{Corollary}
\theoremstyle{definition}\newtheorem{definition}[theorem]{Definition}
\theoremstyle{remark}\newtheorem{remark}[theorem]{Remark}
\newcommand{\R}{\mathbb R}
\newcommand{\Z}{\mathbb Z}
\newcommand{\E}{\mathbb E}
\newcommand{\Prob}{\mathbb P}
\newcommand{\one}{\mathbf 1}
\newcommand{\dd}{\,\mathrm d}
\DeclareMathOperator{\dist}{dist}
\DeclareMathOperator{\covol}{covol}
\DeclareMathOperator{\area}{area}
\DeclareMathOperator{\diam}{diam}
\numberwithin{equation}{section}

\externaldocument[M-][nocite]{main}[A2-v41-primary.pdf]
\title[Scalar collision laws: supplement]{Scalar collision laws and recognition of periodic dispersing billiards\\ Technical supplement}
\author{Qian Qi}
\date{October 4, 2026}
\subjclass[2020]{37D50, 52A10, 60J10, 62G05, 68Q32}
\keywords{Dispersing billiards, reciprocal collision laws, occupation measures,
translation rigidity, directional boundary measures, unknown launch density,
stationary minimax reconstruction}
\hypersetup{pdftitle={Scalar collision laws and recognition of periodic dispersing billiards: Technical supplement},pdfauthor={Qian Qi}}
\begin{document}
\raggedbottom
\begin{abstract}
This supplement contains complete proofs for the localized,
stationary, calibration, period-recognition and directional-germ
experiments associated with the article. It includes the single-law
short-flight inverse, its one-resolved-obstacle extension, and the
finite regularization theorem, together with the sharp stationary
benchmark for a fixed known uniform-disk law. Every theorem is stated
with the preparation model and geometric hypotheses used in its proof.
The primary article proves joint recovery from two fixed opposite
forward fields, a finite realization with the same commands, finite
probability recovery, and density and response estimates under a
variation bound. The present volume supplies the full interpolation,
period-locking and earlier statistical arguments used for comparison
and for the finite geometric construction.
\end{abstract}
\maketitle
\subsection*{Relation to the primary article}
The primary article, Theorem~\ref{M-thm:two-field-rigidity}, proves
joint rigidity from two fixed opposite forward fields. Its finite
geometric and probability conclusions are
Theorems~\ref{M-thm:two-field-finite-geometry} and
\ref{M-thm:two-field-finite-law}. The results in this supplement retain
their hypotheses and internal order. Section~\ref{M-sec:finite-interface}
of the primary article states the complete interface used by its finite
construction and points to the corresponding proofs here. The remaining
results develop the related localized, stationary and calibration
experiments, with their own observation models. In particular, the
directional-germ theorems use angular data and width hypotheses that
are not premises of the primary two-field inverse. This supplement
and the primary article constitute one submission; neither requires
consulting a manuscript revision history to read its theorem statements.
\section{The collision experiment}
\label{sec:setting}
We study joint geometric and probabilistic recovery from collision bits
when the actual launch position is unobserved and a solid start has the
same output as a free miss. The principal inverse identifies both the
obstacle configuration and a single unknown stationary launch density
from the directional first-order short-flight germ. It determines the
full response law and translation period group without assuming
periodicity, a second launch setting or a homothety relation.

We first fix the collision convention and the periodic priors used
only for uniform finite discrete conclusions. The broader exact class
is stated in Section~\ref{sec:single-law-rigidity}. The finite
single-law experiment follows in Section~\ref{sec:single-law-finite}.
The retained pooled, localized and calibration constructions supply
independent benchmarks and the auxiliary interpolation, period and
physical-packing arguments.

\subsection{The periodic class for finite reconstruction}
Fix $0<\beta\le1$ and put $s=6+\beta$. A table is the nonempty union
\begin{equation}\label{eq:presentation}
 \mathcal O=\bigcup_{i=1}^r(C_i+L\Z^2),\qquad 1\le r\le r_0.
\end{equation}
This presentation has exactly $r$ component orbits. It need not be
primitive. The constants in the following bounds are fixed and known:
\begin{equation}\label{eq:priors}
 \|L\|,\|L^{-1}\|\le L_0,\quad |\det L|\ge V_0>0,\quad
 \diam C_i\le D_0,\quad \dist(C,D)\ge d_0>0
\end{equation}
for distinct physical components. All components are strictly convex,
with curvature in $[\kappa_-,\kappa_+]\subset(0,\infty)$ and centered
support functions bounded by $M$ in $C^{6,\beta}(\R/2\pi\Z)$.
Reflections, repeated shapes and nonprimitive presentations are allowed.
Write $\mathcal B$ for this class, $\Pi=\{w:\mathcal O+w=\mathcal O\}$
for its full period group, $z_C$ for a component's Steiner point, and
$p_C$ for its centered support function. All norms on finite-dimensional
spaces are fixed once and for all.

A known positive nonperiod-patch margin $\eta$ will be required only
for uniform finite discrete decisions. Its precise definition is
\eqref{eq:patch-margin}. We denote the resulting subclass by
$\mathcal B_\eta$. The geometric loss is the maximum $C^2$ error of
matched component support functions in a protected laboratory patch,
together with the error of a matched primitive period basis and free
area. Representatives are chosen in a common protected patch. The
basis is matched to a true basis, not required to be a continuously
chosen reduced basis. Exact rational relations and the primitive orbit
count are stated separately from this real-valued loss.

\subsection{Collision bits and the retained localized subexperiment}
For a nominal displacement $a$, let $B_a(q)$ equal one if a free start
at $q$ has a first collision on the unreflected segment $[q,q+a]$.
It equals zero for a solid start and for a free miss. Both are attempted
preparations and enter the denominator. For pointwise statements solids
and swept sets are closed: a boundary solid start is zero, and a start
on a swept boundary but outside every solid is one. This fixes the
tangency convention; changing it on null sets has no effect on the means.
Fix an even smooth probability kernel $\kappa$ supported in the unit disk,
and set
\[
 f_{x,\sigma}(q)=\sigma^{-2}\kappa((q-x)/\sigma),\qquad
 u_\sigma=\kappa_\sigma*\one_\mathcal O.
\]
Choose $\sigma_0,t>0$ with $t+2\sigma_0<d_0$, and prescribe the compass
law
\begin{equation}\label{eq:compass}
 \rho=\tfrac14(\delta_{te_1}+\delta_{-te_1}
                         +\delta_{te_2}+\delta_{-te_2}).
\end{equation}
A forward command draws $q\sim f_{x,\sigma}$ and $a\sim\rho$
independently and attempts $(q,a)$. A reverse command draws that same
nominal joint law and attempts $(q+a,-a).$ Separate fresh preparations
are used; pairing the same physical particle is unnecessary. Only the
two pooled means are used, not means resolved by direction. The joint
reverse law, rather than only its marginal direction law, is prescribed.

In this localized bounded-error experiment, each actual preparation may be coupled to its nominal one with position,
duration and angle errors bounded by $\ell,\tau,\alpha$, respectively.
It remains a unit-speed straight flight until its first impact. The
bounds hold conditionally on past commands, and preparations at a given
command are fresh and conditionally independent. Section~\ref{sec:queries}
proves the bias bound $C(\ell+\tau+t\alpha)/\sigma$. Calibration is a
resource of the controller, not an estimate inferred from the bits.
The launch aperture is fixed by the prior and does not grow as accuracy
increases. It includes the translated reverse starts and the kernel
supports. No free point, obstacle label, impact location, impact angle,
collision time, primitive cell or area normalization is supplied.

\subsection{The directional germ and a single unknown launch law}
The principal inverse uses one stationary launch density $j$, with
unknown compact convex support $A$. The controller prescribes a nominal
center $x$, a direction $n_\theta$ and a length $t>0$, and observes only
the bit $B_{tn_\theta}(x+Z)$, where $Z$ has density $j$. The exact datum
is the family of raw forward means $F_{t,\theta}(x)$ for all centers,
directions and arbitrarily short positive lengths. Directions are
command labels. No launch realization or impact record is observed.

For locally finite separated $C^2$ strictly convex obstacles, the
first-order germ exists in $L^1_{\mathrm{loc}}$:
\begin{equation}\label{eq:intro-single-germ}
 q_\theta=\lim_{t\downarrow0}\frac{F_{t,\theta}}t.
\end{equation}
Let $c_C(\phi)$ and $r_C(\phi)$ be the boundary point with normal
$n_\phi$ and its radius of curvature. The collision-strip formula and
the planar angular identity give
\begin{equation}\label{eq:intro-single-normal}
 (\partial_\theta^2+1)q_\theta(x)
 =\sum_{C,\,\sigma=\pm1}r_C(\theta+\sigma\pi/2)
                     j(c_C(\theta+\sigma\pi/2)-x).
\end{equation}
The right side is a positive measure density. Its spatial components
resolve the antipodal pair left by angular inversion. When the
footprint diameter is smaller than the inter-obstacle separation and
every obstacle width, each component is exactly one translate of
$-A$. Its normalized density recovers $j$ in the centered footprint
frame. The union of its Steiner points over directions recovers the
entire obstacle boundary. Thus
Theorem~\ref{thm:single-law-rigidity} identifies the canonical triple
\begin{equation}\label{eq:intro-single-triple}
       (\mathcal O-s(A),\ A-s(A),\ j(\,\cdot+s(A))).
\end{equation}
It uses neither a second launch setting nor a homothety relation.
The complete observational fiber is the common translation of the
table and density, as stated in Corollary~\ref{cor:single-law-fiber}.
The first-order germ consequently determines every finite-length
forward response and the full translation period group.

Theorem~\ref{thm:single-resolved-obstacle} gives a broader exact
form: it suffices that $\Delta<d$ and one obstacle have diameter
larger than $\Delta$. That obstacle need not be identified in
advance, and no lower width bound is required for the others.
A minimum-area angular support component recovers the common density;
the odd part of the germ recovers occupation and hence the remaining
obstacles. The finite theorem below retains uniform separation of
all angular copies.

The exact datum in \eqref{eq:intro-single-germ} remains a field on the
whole plane. Its limit and angular derivatives are operations on exact
means. Section~\ref{sec:single-law-finite} gives a separate finite
experiment that estimates regularized versions of
\eqref{eq:intro-single-normal} from the original attempted bits.
Spatial and angular smoothing produce a positive comparison field;
integration by parts places all derivatives on known kernels. The
finite-length bias is bounded weakly, without a modulus of continuity
or an upper bound for the unknown density.
For $s=6+\beta$ and the boundary-mass exponent $\gamma$, its
explicit sufficient attempted-bit bound is
\begin{equation}\label{eq:intro-single-finite-cost}
 N_\nu\le C\nu^{-Q_\gamma}\log\frac{C}{\nu\delta},
 \qquad Q_\gamma=\frac{(3\gamma+27/2)s}{s-2}.
\end{equation}
This rate implements the joint unknown-footprint reconstruction;
its necessity is not asserted.
 The small probability of a
collision after spatial averaging gives a variance bound and an
explicit Bernstein sampling cost. A finite reconstruction of both
the footprint and obstacle boundaries follows under the quantitative
smoothness, boundary-mass and aperture assumptions stated there.

\subsection{The retained pooled experiment and its sharp power}
The reciprocal compass experiment uses a different, more restricted
command family. Its difference field $g=F-R$ determines occupation
through a stopped-Poisson inverse. The fixed-stencil formulation in
Theorem~\ref{thm:finite-stencil} is pointwise: its site count depends
on $(D+\Delta)/t$, and costs for several targets are added. It neither
supplies the directional derivatives in \eqref{eq:intro-single-normal}
nor makes a finite global periodicity assertion.

Section~\ref{sec:sharp-stationary} retains the sharp stationary power
for one fixed known uniform-disk law. Put
\[
                    q_0=\frac{3s/2+1}{s-2},\qquad s=6+\beta.
\]
On the same nondegenerate bounded physical class, at the same fixed
confidence and with worst-case expected attempts as cost,
Corollary~\ref{cor:stationary-minimax} gives
\begin{equation}\label{eq:intro-single-benchmark}
 c\nu^{-q_0}\le N^*_{\rm short}(\nu,\delta)
       \le N^*_{\rm pool}(\nu,\delta)
       \le C\nu^{-q_0}\log\frac{C}{\nu\delta}.
\end{equation}
The lower design permits adaptive directions and lengths tending to
zero. The upper design uses a fixed pooled compass with
$2t+2R_+<d_0$; its deterministic attempt cap also bounds expected
cost. One logarithmic factor remains. This fixed-known-law result is
not a minimax assertion for the new joint unknown-density inverse.

\subsection{Data, finite resources and organization}
Each attempted preparation is charged, including a solid start or a
miss. We write $N$ for attempted bits, $J$ for command-labelled batches
with revisits counted again, and $S$ for distinct nominal spatial
sites. Numerical description length, arithmetic and physical
positioning are separate resources. In particular, random centers
used to estimate a scalar integral are counted as well as boundary
search centers.

The new germ inverse uses raw forward means. In the retained
constructions, occupation and the area normalization use reciprocal
differences, whereas width and perimeter deficits use a raw forward
mean. The homothetic inverses continue to assume labelled homothetic
settings and recover their ratios and registrations. Their proofs
remain complete, but that apparatus relation is not a hypothesis of
the principal single-law theorem.

Exact period equality is a whole-field statement without a periodicity
prior. Uniform finite primitive-lattice decisions use the bounded
periodic class and its positive patch margin. A finite nonperiodic
configuration can instead be recovered in a protected aperture known
to contain it completely. These acquisition assumptions appear in
the finite statements; none asserts a finite test for crystallinity
of an arbitrary unknown infinite configuration.

Sections~\ref{sec:single-law-rigidity} and
\ref{sec:single-law-finite} contain the principal exact inverse and
its finite realization. Section~\ref{sec:sharp-stationary} provides
the fixed-known-law benchmark, and Section~\ref{sec:comparison}
places the methods relative to earlier work. The appendices contain
all auxiliary and retained constructions: global and finite-stencil
occupation, stationary boundary queries, homothetic calibration,
period recognition, physical packing and finite controls. The
directional-germ proof does not depend on the stopped-Poisson inverse.

\section{A single unknown launch law from directional short flights}
\label{sec:single-law-rigidity}
We now keep one stationary launch law and allow the direction and
length of a command to be prescribed.  The output of each attempted
preparation is still the bit $B_a$ of Section~\ref{sec:setting}; neither
the actual start nor a collision location or time is observed.  The
launch density and its support are unknown.  No relation between
different launch laws is assumed, because only one law is used.

Let $\mathcal O=\bigcup_{C\in\mathcal C}C$ be a nonempty locally
finite union of compact convex bodies with nonempty interiors,
diameters at most $D$, and mutual distances at least $d>0$.
In this section each boundary is $C^2$ with positive curvature, and
every width of every body is at least $w_*>0$.  Write
\[
 n_\phi=(\cos\phi,\sin\phi),\qquad
 \tau_\phi=(-\sin\phi,\cos\phi),
\]
and denote the point with outward normal $n_\phi$ and its radius of
curvature by $c_C(\phi)$ and $r_C(\phi)$, respectively.  Thus, for the
laboratory support function $h_C$,
\begin{equation}\label{eq:single-gauss-parametrization}
 c_C(\phi)=h_C(\phi)n_\phi+h_C'(\phi)\tau_\phi,
 \qquad r_C(\phi)=h_C(\phi)+h_C''(\phi)>0.
\end{equation}
The footprint $A$ is a compact convex body with nonempty interior,
and $j$ is an $L^1$ probability density supported on $A$ and positive
almost everywhere in its interior.  Their common position is
unrestricted.  The geometric separation assumption is
\begin{equation}\label{eq:single-footprint-separation}
 \diam A\le\Delta< b_*:=\min\{d,w_*\}.
\end{equation}
For example, an inball of radius $r_*$ in each obstacle permits
$w_*=2r_*$.  No periodicity, finite species list, footprint smoothness,
density modulus of continuity, or homothety assumption is imposed.

For $0<t<t_*$ and $\theta\in\R/2\pi\Z$, the observed mean is
\begin{equation}\label{eq:single-raw-response}
 F_{t,\theta}(x)=\int_A j(z)B_{t n_\theta}(x+z)\dd z.
\end{equation}
Here $t_*>0$ is any fixed upper length; limits below use $t<d$.
All solid starts and free misses contribute zero and remain in the
denominator.  The exact datum consists of these raw forward means at
all nominal centers and all prescribed directions for arbitrarily
short positive lengths.  This direction-resolved command family is
richer than a fixed pooled compass field.  It requires no second
setting and no prescribed reciprocal joint preparation.

\begin{lemma}[The short-flight boundary measure]
\label{lem:single-boundary-flux}
As $t\downarrow0$, the fields $t^{-1}F_{t,\theta}$ converge in
$L^1_{\mathrm{loc}}(\R^2)$, locally uniformly in $\theta$, to
\begin{align}
 q_\theta(x)
 &=\sum_{C\in\mathcal C}\int_{\partial C}
       j(y-x)(-n_C(y)\cdot n_\theta)_+\dd s(y)
       \label{eq:single-flux-boundary}\\
 &=\sum_{C\in\mathcal C}\int_0^{2\pi}
       r_C(\phi)j(c_C(\phi)-x)
                 (-\cos(\theta-\phi))_+\dd\phi.
       \label{eq:single-flux-gauss}
\end{align}
For a compactly supported $C^1$ test function $\psi$, there is the
quantitative weak estimate
\begin{equation}\label{eq:single-weak-flight-error}
 \left|\int\psi(x)\{t^{-1}F_{t,\theta}(x)-q_\theta(x)\}\dd x\right|
 \le\frac t2\|\nabla\psi\|_\infty
       \sum_{C\in\mathcal C_\psi}\mathcal W_\theta(C).
\end{equation}
Here $0<t\le t_0<\min\{d,t_*\}$,
$\mathcal W_\theta(C)$ is the width perpendicular to $n_\theta$,
and $\mathcal C_\psi$ consists of the finitely many bodies meeting
$\operatorname{supp}\psi+A+t_0\overline B$.
\end{lemma}
\begin{proof}
Use coordinates $y=w\tau_\theta+u n_\theta$.  For $w$ in the
projection interval $I_{C,\theta}$, the section of $C$ is an interval
$[\ell_{C,\theta}(w),u_{C,\theta}(w)]$.  Its incoming boundary
point is
$y_{C,\theta}(w)=w\tau_\theta+\ell_{C,\theta}(w)n_\theta$.
The free starts whose segment first meets this body form precisely
the strip
\[
 \{y_{C,\theta}(w)-s n_\theta:
               w\in I_{C,\theta},\ 0<s\le t\},
\]
up to planar null sets.  Since $t<d$, a segment cannot meet two
different bodies and this strip meets no other solid.  The strips
therefore give, almost everywhere in $x$,
\begin{equation}\label{eq:single-strip-disintegration}
 \frac{F_{t,\theta}(x)}t
 =\sum_C\int_{I_{C,\theta}}\frac1t\int_0^t
   j(y_{C,\theta}(w)-s n_\theta-x)\dd s\dd w.
\end{equation}
On a compact set of nominal centers only finitely many bodies occur,
uniformly in $\theta$ and $0<t\le t_0$.  Translation continuity in
$L^1(\R^2)$ shows that replacing $s$ by zero in
\eqref{eq:single-strip-disintegration} has local $L^1$ error at most
$ND\omega_j(t)$, where $N$ is this finite component count and
\[
 \omega_j(t)=\sup_{|v|\le t}\|j(\,\cdot+v)-j\|_{L^1(\R^2)}
                         \longrightarrow0.
\]
The projection identity
$\dd w=(-n_C\cdot n_\theta)_+\dd s$ on the incoming boundary
proves \eqref{eq:single-flux-boundary}.  The Gauss parametrization
and $\dd s=r_C(\phi)\dd\phi$ give
\eqref{eq:single-flux-gauss}.

To prove the weak estimate, integrate
\eqref{eq:single-strip-disintegration} against $\psi$ and change
variables from $x$ to $z=y_{C,\theta}(w)-s n_\theta-x$.
The two test values differ by at most
$s\|\nabla\psi\|_\infty$.  Integration against the probability
density $j$ preserves this bound, and averaging $s$ over $[0,t]$
gives $t/2$.  Finally,
$|I_{C,\theta}|=\mathcal W_\theta(C)$, and bodies outside
$\mathcal C_\psi$ make no contribution.
\end{proof}

\begin{lemma}[Angular resolution of the launch density]
\label{lem:single-angular-resolution}
The field $q$ is twice continuously differentiable in $\theta$ with
values in $L^1_{\mathrm{loc}}(\R^2)$.  Its angular derivative
\begin{equation}\label{eq:single-angular-field}
 S_\theta:=(\partial_\theta^2+1)q_\theta
   =\sum_{C\in\mathcal C}\sum_{\sigma\in\{-1,1\}}
       r_C(\theta+\sigma\pi/2)
       j(c_C(\theta+\sigma\pi/2)-\,\cdot\,)
\end{equation}
is nonnegative.  For every $\theta$, the connected components of
the support of the measure $S_\theta(x)\dd x$ are exactly
\begin{equation}\label{eq:single-angular-components}
 K_{C,\sigma,\theta}=c_C(\theta+\sigma\pi/2)-A.
\end{equation}
Different such components have distance at least $b_*-\Delta>0$.
\end{lemma}
\begin{proof}
For $k(\alpha)=(-\cos\alpha)_+$, direct differentiation on the
two open semicircles and the jumps of $k'$ give, as periodic
distributions,
\begin{equation}\label{eq:single-cosine-distribution}
             k''+k=\delta_{\pi/2}+\delta_{-\pi/2}.
\end{equation}
Apply this identity to \eqref{eq:single-flux-gauss}.  Local finiteness
justifies termwise distributional differentiation.  The right side
of \eqref{eq:single-angular-field} is continuous in $\theta$ in
$L^1$ on each compact set: each $r_C$ and $c_C$ is continuous,
translations act continuously on $L^1$, and only finitely many bodies
can contribute there.  The same is true of $q$ and its first
derivative, obtained by using the bounded piecewise continuous kernel
$k'$.  The distributional identity then proves the asserted $C^2$
regularity and specifies $S_\theta$ canonically for every angle.

The support of the positive measure with density
$j(c-x)$ is exactly $c-A$.  Positivity almost everywhere in
$\operatorname{int}A$ proves this assertion even if a chosen
representative of $j$ vanishes on a dense null set.  The weights in
\eqref{eq:single-angular-field} are positive.  Contact points on
distinct obstacles have distance at least $d$.  The two contact
points belonging to one obstacle have opposite normals, and
\[
 |c_C(\phi)-c_C(\phi+\pi)|
 \ge (c_C(\phi)-c_C(\phi+\pi))\cdot n_\phi
 =h_C(\phi)+h_C(\phi+\pi)\ge w_*.
\]
For two contact points $c,c'$ and any $a,a'\in A$,
$|(c-a)-(c'-a')|\ge|c-c'|-\Delta$.
Consequently the copies in \eqref{eq:single-angular-components}
have the stated positive separation.  They are connected compact
sets and form a locally finite family.  They are therefore precisely
the connected components of the measure support.
\end{proof}

Let $s(H)$ denote the Steiner point of a convex body $H$.  Its
translation covariance and $s(-H)=-s(H)$ will fix the common
translation of the unknown table and launch law.  Define
\begin{equation}\label{eq:single-canonical-parameters}
 A_0=A-s(A),\qquad j_0(z)=j(z+s(A)),\qquad
                      \mathcal O_0=\mathcal O-s(A).
\end{equation}

\begin{theorem}[Single-law directional-germ rigidity]
\label{thm:single-law-rigidity}
Under \eqref{eq:single-footprint-separation}, the directional
short-flight germ $\{q_\theta:\theta\in\R/2\pi\Z\}$ determines
the complete canonical triple $(\mathcal O_0,A_0,j_0)$, with equality
of densities understood almost everywhere.  In particular, the raw
mean family \eqref{eq:single-raw-response} determines this triple.

More explicitly, select any angle and any connected component $K$
of $\operatorname{supp}(S_\theta(x)\dd x)$.  Let
$S_{\theta,K}=\one_KS_\theta$ and
$m_K=\int_KS_\theta(x)\dd x$.  Then
\begin{equation}\label{eq:single-density-recovery}
 A_0=-(K-s(K)),\qquad
 j_0(z)=\frac{S_{\theta,K}(s(K)-z)}{m_K}.
\end{equation}
To recover the table without matching components across angles, put
\begin{equation}\label{eq:single-boundary-recovery}
 \Gamma=\bigcup_{\theta\in\R/2\pi\Z}
       \{s(K):K\text{ is a connected component of }
                        \operatorname{supp}(S_\theta(x)\dd x)\}.
\end{equation}
Then $\Gamma=\partial\mathcal O_0$; its connected components are
the individual obstacle boundaries, and
\begin{equation}\label{eq:single-table-recovery}
 \mathcal O_0=\bigcup_{\Lambda\in\operatorname{Comp}(\Gamma)}
                                      \operatorname{conv}\Lambda.
\end{equation}
\end{theorem}
\begin{proof}
By Lemma~\ref{lem:single-angular-resolution}, the selected component
has the form $K=c-A$, and on that component the field equals
$rj(c-x)$ for one strictly positive curvature radius $r$.
Consequently
\[
 s(K)=c-s(A),\qquad m_K=r\int_Aj=r.
\]
Reflection about $s(K)$ and division by this mass give
\eqref{eq:single-density-recovery}.  This proves that the result is
independent of the selected component and angle and recovers both
the support and the complete centered density.

Again by \eqref{eq:single-angular-components}, the set of component
Steiner points at angle $\theta$ is
\[
 \{c_C(\theta-\pi/2)-s(A),
               c_C(\theta+\pi/2)-s(A):C\in\mathcal C\}.
\]
As $\theta$ ranges over the circle, the Gauss parametrization ranges
over every boundary point of every obstacle.  Thus their union is
exactly $\bigcup_C\partial(C-s(A))$.
Each of these boundaries is connected, different boundaries have
distance at least $d$, and the family is locally finite.
Their union is closed and its connected components are precisely
the individual boundaries.  The same separation shows that this
union equals the boundary of the obstacle union.  A compact convex
body with nonempty interior is the convex hull of its boundary,
which proves \eqref{eq:single-table-recovery}.
\end{proof}

\begin{corollary}[The complete observational fiber]
\label{cor:single-law-fiber}
Two configurations and single stationary launch laws in the stated
class have the same directional germ if and only if there is
$h\in\R^2$ such that
\begin{equation}\label{eq:single-law-fiber}
 \mathcal O'=\mathcal O+h,\qquad A'=A+h,
                         \qquad j'(z)=j(z-h)\quad\text{a.e.}
\end{equation}
These conditions are also equivalent to equality of the entire
raw forward response for every nominal center and every finite
command displacement for which the bit is defined.
\end{corollary}
\begin{proof}
Equal germs give equal canonical triples by
Theorem~\ref{thm:single-law-rigidity}.  Taking
$h=s(A')-s(A)$ yields \eqref{eq:single-law-fiber}.
Conversely, under these relations couple a start $x+Z$ for the first
table with $x+Z+h$ for the second.  The two unreflected segments and
all obstacles differ by the same translation.  The bits therefore
agree for every displacement, including the solid-start convention.
Equality of the full raw family implies equality of its short-flight
germ by Lemma~\ref{lem:single-boundary-flux}.
\end{proof}

For the germ, translation equality is understood in
$L^1_{\mathrm{loc}}(\R^2)$, equivalently almost everywhere in the
nominal center for each angle.  For the continuous raw mean fields
it is pointwise.  The germ period group is therefore independent
of representatives of the density and flux fields.

\begin{corollary}[Germ periods and full response completion]
\label{cor:single-germ-periods}
The first-order directional germ determines the full forward
collision law at all finite command lengths.  Moreover,
\begin{equation}\label{eq:single-germ-periods}
 \operatorname{Per}(\{q_\theta\}_\theta)
 =\operatorname{Per}(\{F_{t,\theta}\}_{0<t<t_*,\theta})
 =\operatorname{Per}(\mathcal O).
\end{equation}
The common group is discrete.  It has rank two exactly when the
configuration admits a rank-two lattice of periods with finitely
many component orbits.  No periodicity hypothesis enters this
criterion.
\end{corollary}
\begin{proof}
The canonical triple supplies the defining integral of every
forward bit, and a common translation changes none of these means.
This proves response completion.  A period of the raw family is a
period of its germ.  A period of the germ preserves each measure
support in Lemma~\ref{lem:single-angular-resolution}, permutes its
connected components, and translates their Steiner points.
It therefore preserves $\Gamma$ and, by
\eqref{eq:single-table-recovery}, the obstacle union $\mathcal O_0$.
Every period of $\mathcal O$ is conversely a period of all raw
means by the collision-bit definition.
A nonzero translation period moves any compact obstacle to a
different one, and hence has length at least $d$.  The group is
discrete; the rank-two assertion follows from local finiteness in
a compact fundamental cell, as in
Theorem~\ref{thm:global-response-rigidity}.
\end{proof}

All these conclusions concern the exact direction-resolved active
germ on the plane.  The angular operation in
\eqref{eq:single-angular-field} is applied to response means and
is not an additional sensor output.  Its estimation from finitely
many attempted bits requires regularization and separate resource
bounds.  The proof above uses neither a reciprocal occupation inverse
nor a homothetic second footprint: angular differentiation separates
the two tangent contacts, and their common density identifies the
unknown launch law.  The small-footprint condition ensures that these
translated density copies can be separated spatially; no conclusion
without that condition is asserted here.

\subsection{One resolved obstacle suffices for exact recovery}
\label{sec:single-resolved-extension}
The uniform width restriction can be removed from all but one
obstacle.  This extension concerns exact recovery; the uniform
separation of every angular density copy remains the assumption for
the finite reconstruction below.  Equalities of germs, and their
translation periods, are understood in $L^1_{\mathrm{loc}}(\R^2)$
for each angle.  Thus they are independent of representatives of
the unknown density.

\begin{theorem}[Rigidity from one resolved obstacle]
\label{thm:single-resolved-obstacle}
Retain the single-law experiment, the locally finite convex
configuration, and the boundary regularity of
Section~\ref{sec:single-law-rigidity}, but replace
\eqref{eq:single-footprint-separation} by
\begin{equation}\label{eq:single-one-body-separation}
 \diam A\le\Delta<d,\qquad
       \diam C_*>\Delta\quad\hbox{for at least one }C_*\in\mathcal C.
\end{equation}
No lower width bound is required for the other obstacles, and the
identity or position of $C_*$ is not supplied.  The directional
germ determines $(\mathcal O_0,A_0,j_0)$ of
\eqref{eq:single-canonical-parameters}.  Its complete observational
fiber is \eqref{eq:single-law-fiber}; it determines the forward
collision law at every finite command length and the full obstacle
period group as in \eqref{eq:single-germ-periods}.

The footprint and density are recovered by
\eqref{eq:single-density-recovery} from any angular support component
of minimum area, where the minimum is taken over all components and
all angles.  This minimum is attained and equals $\area A$.
The occupation field is then the unique bounded continuous function
with $\inf v=0$ and distributional derivatives
\begin{equation}\label{eq:single-odd-occupation}
 \partial_{e_1}v=q_0-q_\pi,
 \qquad \partial_{e_2}v=q_{\pi/2}-q_{3\pi/2}.
\end{equation}
If $P$ ranges over the closures of the connected components of
$\{v>0\}$, the support functions of the recovered obstacles are
\begin{equation}\label{eq:single-occupation-cancellation}
                 h_{C_P^0}=h_P-h_{-A_0},
 \qquad \mathcal O_0=\bigcup_P C_P^0.
\end{equation}
\end{theorem}
\begin{proof}
The short-flight and angular identities
\eqref{eq:single-flux-boundary}--\eqref{eq:single-angular-field}
remain valid: their derivations do not use a comparison between
footprint size and obstacle width.  At each angle, different
obstacles contribute pairs of density supports whose mutual
distances are at least $d-\Delta>0$.  Within one pair the two
supports are distinct translates of $-A$; they may intersect.
Consequently each connected component of the angular measure
support is either one translate of $-A$ or the union of two
distinct translates.  Every component in the first case has area
$\area A$, while every component in the second has strictly larger
area.  Indeed, if two translates differ by $u\ne0$, their support
values in direction $u/|u|$ differ by $|u|$; the translate with
larger support contains a nonempty open cap outside the other.

Choose a diameter pair $p,q$ of $C_*$.  The maximality of
$|p-q|$ implies that $p$ and $q$ have outward normals
$n=(p-q)/|p-q|$ and $-n$, respectively.  For example,
$n\cdot(y-q)\le|y-q|\le|p-q|$ for every $y\in C_*$,
and the reverse endpoint is treated in the same way.
At the corresponding angular slice, $p-A$ and $q-A$ have
distance at least $|p-q|-\Delta>0$.  They are also separated
from every other obstacle's supports.  Both are therefore
individual support components of area $\area A$.
This proves that the global minimum is attained, and that each
minimizer is a single copy $K=c-A$.  On such a component
$S_\theta=rj(c-\,\cdot\,)$ for one $r>0$.
Its mass is $r$ and its Steiner point is $c-s(A)$, proving
\eqref{eq:single-density-recovery} exactly as before.

For completeness, define the physical occupation
\[
                 v(x)=\int_A j(z)\one_{\mathcal O}(x+z)\dd z.
\]
Since $(-a)_+-(a)_+=-a$, the flux identity gives
\begin{equation}\label{eq:single-odd-flux}
 q_\theta-q_{\theta+\pi}
   =-\sum_C\int_{\partial C}
          j(y-\,\cdot\,)\,n_C(y)\cdot n_\theta\dd s(y)
   =\partial_{n_\theta}v
\end{equation}
in distributions.  The last equality follows by testing against a
compactly supported smooth function and applying the divergence
theorem on each relevant body.  Fubini's theorem is applicable,
and only finitely many bodies meet the translated test support.
Translation continuity of $j$ in $L^1$ gives
\[
 \|v(\,\cdot+u)-v\|_\infty
                    \le\|j(\,\cdot-u)-j\|_{L^1}\longrightarrow0.
\]
Thus $v$ is bounded and continuous.  Positivity of $j$ almost
everywhere on $\operatorname{int}A$ and convexity give
\begin{equation}\label{eq:single-resolved-positivity}
        \{v>0\}=\bigcup_C\operatorname{int}(C+(-A)).
\end{equation}
Indeed, a positive integral is equivalent to a positive-area
intersection of $x+A$ with a body, and hence to an intersection of
their interiors.  The closures $C+(-A)$ have mutual gaps at least
$d-\Delta$.  A boundary point of any one of them lies in none of
these interiors and has zero occupation.  Therefore $\inf v=0$.
Two bounded continuous functions satisfying
\eqref{eq:single-odd-occupation} differ by a distribution with
zero gradient, hence by a constant.  This elementary assertion
also follows by convolution with smooth approximate identities:
every convolved difference has zero classical gradient and is
constant, and local uniform convergence gives the assertion for
the original difference.  The two zero-infimum normalizations
fix that constant.  Thus the germ determines $v$ uniquely.

The connected-component closures in
\eqref{eq:single-resolved-positivity} are exactly $P=C+(-A)$.
Writing $A=A_0+s(A)$ yields
$P=(C-s(A))+(-A_0)$.  Additivity of support functions proves
\eqref{eq:single-occupation-cancellation} and recovers every
obstacle, including those whose two angular copies overlap.
Equal germs therefore give equal canonical triples; conversely
a common translation couples every collision bit.  This proves
the claimed fiber and completion of the full forward law.

Finally, if $u$ is a period of the germ, then
$v(\,\cdot+u)$ has the same distributional gradient and the same
zero infimum as $v$.  Uniqueness gives $v(\,\cdot+u)=v$.
Translation by $u$ therefore permutes the sets $C+(-A)$;
cancelling the common footprint support proves that it permutes
the obstacles.  The converse follows from translation covariance
of the raw means.  The common period group is discrete, since a
nonzero period moves a compact component to another component at
distance at least $d$; the rank-two criterion follows as in
Corollary~\ref{cor:single-germ-periods}.
\end{proof}

\section{Finite reconstruction from one unknown launch law}
\label{sec:single-law-finite}
The angular identity also has a finite statistical realization.  Its
regularization must preserve positivity: a pointwise numerical second
derivative of a measured field would neither supply the necessary
support test nor admit a uniform bias estimate for arbitrary $L^1$
densities.  We instead differentiate a known angular test function and
average the spatial command centers with positive weights.  Both
averages are implemented by finite rational menus of ordinary forward
commands.

In this section $s=6+\beta$, and the obstacle class has the bounds of
\eqref{eq:priors}, including uniformly positive upper and lower
curvature radii.  All widths are at least a known $w_*>0$.  The one
unknown footprint and density satisfy
\begin{equation}\label{eq:single-finite-priors}
 \begin{gathered}
 \diam A\le\Delta<\min(d_0,w_*),\qquad
 \|p_A^\circ\|_{C^{6,\beta}}\le M_A,\\
 0<r_A^-\le p_A^\circ+(p_A^\circ)''\le r_A^+,
 \qquad j(z)\ge b_0\dist(z,\partial A)^\gamma
       \quad\hbox{for a.e. }z\in A,
 \end{gathered}
\end{equation}
where $b_0>0$ and $\gamma\ge0$ are fixed.  A uniform inball bound
could equivalently be added to these priors; the displayed lower
curvature radius already supplies one.  No upper bound, smoothness,
translation modulus, or evaluations of $j$ are available.  Put
\[
 a_0=s(A),\qquad A_0=A-a_0,\qquad
 \mathcal O_0=\mathcal O-a_0,
 \qquad \kappa_*:=\min(d_0,w_*)-\Delta>0.
\]
Here $s(A)$ denotes the Steiner point.  It is an unknown used to state
the identifiable output, not a supplied coordinate.  The experiment
uses only the means \eqref{eq:single-raw-response} and the individual
forward bits whose means they are.  All nominal centers and displacement
vectors selected below are rational.

\subsection{A positive weak response and its finite estimator}
Let $\chi(u)=c_\chi(1-u^2)^4$ for $|u|\le1$ and zero otherwise,
where $c_\chi$ normalizes its integral to one.  This constant is
rational.  Define
\begin{equation}\label{eq:single-positive-kernels}
 \kappa_r(u)=r^{-2}\chi(u_1/r)\chi(u_2/r),\qquad
 K_b(v)=b^{-1}\chi(v/b),\qquad \lambda_b=K_b''+K_b.
\end{equation}
The kernels have compact support and enough endpoint vanishing for
twice integrating by parts.  They may also be approximated by smooth
test functions in that identity.  Their bounds include
\begin{equation}\label{eq:single-kernel-bounds}
 \|\kappa_r\|_\infty\le Cr^{-2},\quad
 \|\nabla\kappa_r\|_\infty\le Cr^{-3},\quad
 \|\lambda_b\|_1\le Cb^{-2},\quad
 \|\lambda_b'\|_1\le Cb^{-3},\quad
 \int\lambda_b=1.
\end{equation}
Angles are integrated on a real interval around the chosen direction;
the physical responses are extended periodically.  Thus no choice of
a branch for an angular coordinate enters the construction.

Let $S_\theta$ be the positive angular measure in
\eqref{eq:single-angular-field}.  At a spatial target $x$ and angular
target $\alpha$ define the nonnegative function
\begin{align}
 U_{r,b}(x,\alpha)
  &=\int_{\R}\!K_b(v)\int_{\R^2}\!\kappa_r(u)
                      S_{\alpha+v}(x-u)\dd u\dd v
                       \label{eq:single-positive-response}\\
  &=\int_{\R}\!\lambda_b(v)\int_{\R^2}\!\kappa_r(u)
                      q_{\alpha+v}(x-u)\dd u\dd v.
                       \label{eq:single-weak-response}
\end{align}
The first formula defines a continuous nonnegative spatial function;
the second follows from Lemma~\ref{lem:single-angular-resolution}.
These are convolutions of locally finite measures.  In particular,
neither expression requires pointwise values of $j$.

\begin{lemma}[Finite sampling of the weak angular response]
\label{lem:single-weak-sampling}
Fix a bounded nominal aperture and all the stated priors.  For
$0<r,b,t<1$ sufficiently small, let $\ell$ be a spatial quadrature
mesh, $d_\theta$ an angular quadrature mesh, and $\zeta$ a bound on
the rounding of each displacement vector, with
\[
 \ell\le c\min(r,t),\qquad d_\theta\le cb,\qquad \zeta\le ct.
\]
There is a finite estimator $\widehat U(x,\alpha)$ using only forward
collision bits, whose deterministic bias relative to
$U_{r,b}(x,\alpha)$ is at most
\begin{equation}\label{eq:single-weak-bias}
 C\left\{tr^{-3}b^{-2}
       +\frac{\ell r^{-2}b^{-2}}t
       +d_\theta r^{-2}b^{-3}
       +\frac{\zeta r^{-2}b^{-2}}t\right\}.
\end{equation}
For any $a>0$ and $0<\varepsilon<1$, an additional statistical error
at most $a$ has probability at least $1-\varepsilon$ after
\begin{equation}\label{eq:single-weak-cost}
 n\le C\left(t^{-1}r^{-2}b^{-4}a^{-2}
                  +t^{-1}b^{-2}a^{-1}\right)
                    \log\frac{C}{\varepsilon}
\end{equation}
attempted bits, with an integer ceiling understood.  The constants
are uniform over the unknown density.  A fresh rational center and
vector selection is charged as a command occurrence, even when it
is chosen by randomizing a nominal integral.
\end{lemma}
\begin{proof}
First replace $q_\theta$ in \eqref{eq:single-weak-response} by
$t^{-1}F_{t,\theta}$.  Apply
\eqref{eq:single-weak-flight-error} to
$\psi(y)=\kappa_r(x-y)$ and integrate against $|\lambda_b|$.
Separation and the fixed diameter and aperture bounds give a uniform
bound on the number and widths of all components meeting these tests.
The first term in \eqref{eq:single-weak-bias} follows.  Integrating
over $j$ costs only its total mass one.  Thus an unobserved modulus
of continuity of $j$ has not been used in this replacement.

We give the finite quadrature details because rounding pointwise
values of an arbitrary density would not give this conclusion.  Put
\[
 w_u=\ell^2\kappa_r(u),\quad u\in\ell\Z^2,\qquad
 Z=\sum_u w_u,
 \qquad v_k=d_\theta\lambda_b(kd_\theta),\quad B=\sum_k|v_k|.
\]
Only the supports of the kernels are included.  For small fixed mesh
ratios, $1/2\le Z\le2$ and $B\le Cb^{-2}$.  For dyadic $r,b,\ell$
and $d_\theta$, all these weights are rational and are specified by
the fixed polynomials in \eqref{eq:single-positive-kernels}.

For a fixed launch displacement $z$ and fixed command $a$, the
function of the spatial variable being quadratured is a smooth weight
times the indicator of a collision strip.  For $|a|<d_0$, such a
strip is the difference between a convex swept body and its original
convex body.  In the relevant bounded region their total perimeter
is uniformly bounded.  The union of square mesh cells of side $\ell$
meeting these boundaries has area at most $C\ell$; this follows by
covering each convex boundary by its inner and outer parallel bands.
The weight variation on the remaining cells has integral at most
$C\ell r^{-1}$.  Since $\|\kappa_r\|_\infty\le Cr^{-2}$, the total
quadrature error is at most $C\ell r^{-2}$.  The estimate is uniform
in the translation $z$.  Integration against $j(z)\dd z$, division
by $t$ and summation against $|v_k|$ therefore give the second term
in \eqref{eq:single-weak-bias}.

For the angular quadrature write
\[
 A_t(\theta)=\int\kappa_r(u)F_{t,\theta}(x-u)\dd u.
\]
The swept strip has area at most $Ct$, and changing its displacement
vector from $t n_\theta$ to $t n_{\theta'}$ changes its area in
symmetric difference by at most $Ct|\theta-\theta'|$.  To see the
latter estimate, their convex swept bodies have Hausdorff distance
at most $t|n_\theta-n_{\theta'}|$ and uniformly bounded perimeter;
the unchanged solid is subtracted from both.  Consequently
\[
 0\le A_t(\theta)\le Ct r^{-2},\qquad
 |A_t(\theta)-A_t(\theta')|
                    \le Ct r^{-2}|\theta-\theta'|.
\]
These estimates remain uniform after integration over $j$.  The
one-dimensional quadrature error for $\lambda_b A_t$ is bounded by
its total variation times $d_\theta$.  Equation
\eqref{eq:single-kernel-bounds}, followed by division by $t$, gives
$Cd_\theta r^{-2}b^{-3}$.  Apply this argument to the continuous
spatial average; the spatial quadrature error at the finitely many
angular nodes has already been bounded separately.

Round $t n_{\alpha+kd_\theta}$ to a rational vector $a_k$ at distance
at most $\zeta$.  The same swept-body area estimate gives an error
$C\zeta r^{-2}$ in each continuous spatial mean.  The spatial
quadrature estimate is valid for the rounded vectors as well, so
this gives the final term of \eqref{eq:single-weak-bias}.

For each attempt select $u$ with probability $w_u/Z$, and select $k$
independently with probability $|v_k|/B$.  Attempt the ordinary
forward command at nominal center $x-u$ with displacement $a_k$, and
write its collision bit as $Y$.  The signed random variable
\begin{equation}\label{eq:single-weighted-observation}
                   W=\frac{ZB}{t}\operatorname{sgn}(v_k)Y
\end{equation}
has expectation exactly equal to the finite quadrature just described.
Zero weights can be omitted.  This selection does not modify the
stationary launch law: the random choices select known nominal
commands, each followed by a fresh draw from the same unknown $j$.

Uniformly at every angular node, the mean collision probability under
the spatial proposal is at most
$C(t+\ell+\zeta)r^{-2}\le Ct r^{-2}$.  It follows that
\begin{equation}\label{eq:single-weak-variance}
 |W|\le Ct^{-1}b^{-2},\qquad
 \operatorname{Var}(W)\le\E W^2\le Ct^{-1}r^{-2}b^{-4}.
\end{equation}
Bernstein's inequality for the average of independent copies of $W$
proves \eqref{eq:single-weak-cost}.  All observations, including zeros
from solid starts, enter this average.  The selection probabilities
and multipliers are known rational numbers; no response mean or
unknown density is an input to them.
\end{proof}

\begin{lemma}[Separated footprint copies from a finite threshold field]
\label{lem:single-finite-copies}
Let $0<e<e_0$ and choose $r,b$ comparable to sufficiently small fixed
multiples of $e$.  At each angular target $\alpha$ the function
$U_{r,b}(\cdot,\alpha)$ is supported in the union of the $Ce$-parallel
bodies of
\begin{equation}\label{eq:single-target-copies}
                 K_{C,\sigma}=c_C(\alpha+\sigma\pi/2)-A,
                         \qquad \sigma\in\{-1,1\}.
\end{equation}
These parallel bodies remain separated by at least $\kappa_*/2$.
There is a known $c_1>0$ such that
\begin{equation}\label{eq:single-interior-signal}
 U_{r,b}(x,\alpha)\ge c_1e^\gamma
            \quad\hbox{if }x\in(K_{C,\sigma})_{-e}.
\end{equation}
Suppose its values on a spatial mesh of width $c e$ in a protected
aperture are estimated with error at most $c_1e^\gamma/8$.
Thresholding, a fixed-distance graph, and convex hulls then recover
every complete protected copy in Hausdorff distance at most $Ce$,
with exactly one recovered cluster per copy.  Their Steiner points
estimate $c_C(\alpha+\sigma\pi/2)-a_0$ with error at most $Ce$.
\end{lemma}
\begin{proof}
Write $\phi=\alpha+\sigma\pi/2$.  The radius bound gives
$|c_C(\phi+v)-c_C(\phi)|\le r_C^+|v|$.  Each term in the first
formula \eqref{eq:single-positive-response} is therefore supported
in $K_{C,\sigma}+(r_C^+b+\sqrt2r)\overline B$.  Choose these two
widths to have sum at most $e/4$.  The exact separation of the copies
is at least $\kappa_*$, proving the support and gap assertions.

If $x\in(K_{C,\sigma})_{-e}$, then every displacement
$c_C(\phi+v)-x+u$ occurring in this term lies inside $A$ at depth
at least $3e/4$.  The lower bound in
\eqref{eq:single-finite-priors}, the lower obstacle curvature radius,
and the unit integrals of both positive kernels give
\eqref{eq:single-interior-signal}.  The other terms are nonnegative.
This argument integrates the almost-everywhere density inequality;
it requires no pointwise choice of its representative.

Retain grid nodes whose estimated value exceeds $c_1e^\gamma/4$.
Every node at depth $e$ is retained, and no node outside the
$e/4$-parallel bodies is retained.  Give two retained nodes an edge
when their distance is less than $\kappa_*/4$.  Distinct copies cannot
be joined for small $e$.  Within a copy, the nodes in a slightly deeper
erosion form a connected set for this fixed edge length: segments
joining points of the convex erosion can be followed by a grid chain
with error $O(e)$.  Every retained node lies within $Ce$ of that
erosion and hence of a retained deep node.  To justify this last
assertion uniformly up to the boundary, the footprint's lower
curvature radius implies that its inner parallel body at distance
$u<r_A^-$ has support $p_A-u$, so its Hausdorff distance from $A$
is $u$.  Alternatively its fixed inball gives the same bound by a
homothety towards an inball center.  Thus each complete copy yields
exactly one graph component.

The hull of its retained nodes lies in $K_{C,\sigma}+Ce\overline B$
and contains a grid approximation of $(K_{C,\sigma})_{-2e}$.
Its support consequently differs from that of $K_{C,\sigma}$ by
at most $Ce$.  Steiner centering is Lipschitz for support supremum
error, with constant two in the plane, proving the final assertion.
\end{proof}

\begin{theorem}[Finite forward reconstruction at one unknown setting]
\label{thm:single-law-finite}
Fix the priors \eqref{eq:single-finite-priors} and the bounded periodic
class $\mathcal B_\eta$.  At a single unknown stationary launch law,
a finite controller using only forward collision bits reconstructs
$A_0$ and a smooth strictly convex periodic table approximating
$\mathcal O_0$.  For all sufficiently small $\nu>0$ and
$0<\delta<1/4$, its footprint $C^2$ support error and the geometric
loss of the table are at most $C\nu$, and its primitive discrete data
are correct, with probability at least $1-\delta$.

Put
\begin{equation}\label{eq:single-finite-exponent}
                 Q_\gamma=\frac{(3\gamma+27/2)s}{s-2}.
\end{equation}
The controller has the deterministic attempted-bit bound
\begin{equation}\label{eq:single-finite-resources}
 N_\nu\le C\nu^{-Q_\gamma}\log\frac{C}{\nu\delta},
             \qquad S_\nu\le J_\nu\le N_\nu,
\end{equation}
where $J_\nu$ counts center-and-vector-labelled command occurrences
and $S_\nu$ counts distinct nominal sites.  The nominal centers remain in a
fixed prior-dependent aperture.  If a known coarse bound on $|a_0|$ is
also imposed, all actual launches and attempted flights lie in a
prior-bounded physical aperture; the bound does not supply the origin.
The actual starts lie in $W+A+O(r)\overline B$, and their flights
in its $t$-parallel body.  Its directional commands have lengths
of order
\begin{equation}\label{eq:single-finite-precision}
 \nu^{(\gamma+5)s/(s-2)},
 \qquad\hbox{and coordinate mesh of order }
             \nu^{(2\gamma+9)s/(s-2)}.
\end{equation}
The total command description is at most
$C J_\nu\log(C/(\nu\delta))$ bits, for fixed computable priors and
certified dyadic choices of these scales.  Numerical construction of
this finite menu and the geometric output is separate from the
observation count.  No sharpness of the exponent $Q_\gamma$ is asserted.
The fixed-known-disk minimax theorem concerns its own experiment and
retains its stated exponent.

The footprint translation $a_0$ is not estimated separately, as it
belongs to the common translation ambiguity of the experiment.  There
is one setting, no homothety control, no reciprocal preparation, and
no supplied density or footprint evaluations.  The theorem estimates
the geometry; it makes no strong-norm finite density estimate under
an unrestricted $L^1$ density class.
\end{theorem}
\begin{proof}
Choose a dyadic $e$ comparable to a sufficiently small multiple of
$\nu^{s/(s-2)}$.  In Lemma~\ref{lem:single-weak-sampling} take dyadic
scales comparable to
\begin{equation}\label{eq:single-finite-scales}
 r\asymp e,\quad b\asymp e,\quad
 t\asymp e^{\gamma+5},\quad
 \ell\asymp\zeta\asymp e^{2\gamma+9},\quad
 d_\theta\asymp e^{\gamma+5}.
\end{equation}
The constants in $r,b$ first meet the support reserves of
Lemma~\ref{lem:single-finite-copies}; the constant in $t$ and then
those in $\ell,d_\theta,\zeta$ make the four bias terms sum to at
most $c_1e^\gamma/16$.  These choices satisfy all mesh and length
restrictions.  With additional statistical error
$a=c_1e^\gamma/16$, \eqref{eq:single-weak-cost} becomes
\begin{equation}\label{eq:single-per-target-count}
                  n\le C e^{-(3\gamma+11)}
                                   \log(C/\varepsilon).
\end{equation}
Indeed the two Bernstein powers are $3\gamma+11$ and $2\gamma+7$,
and the first dominates.  The spatial boundary quadrature, rather
than an assumed continuity modulus for $j$, justifies the much finer
rational command mesh in \eqref{eq:single-finite-scales}.

We specify the aperture and discard rule before using these estimates.
Choose a fixed radius $R$ that contains the entire protected patch
needed by Lemma~\ref{lem:locking}, in the frame of $\mathcal O_0$.
A presentation lattice with the fixed bounds supplies representatives
in a prior-bounded box after any common translation, so no coordinate
of $a_0$ is needed for this choice.  Take a square $W$ whose interior
contains the ball of radius
$R+3D_0+5\Delta+10$ and its fixed response collars.  Sampling takes
place at the spatial target grid in $W$; the nominal centers used by
the quadrature extend this only by $\sqrt2r$.

At each angular target form the graph of Lemma~\ref{lem:single-finite-copies}.
Keep only graph components with a node at distance greater than
$2\Delta+1$ from $\partial W$.  Any such node is within $Ce$ of its
true copy, whose diameter is at most $\Delta$.  Its entire copy,
including the required $Ce$ margin, is therefore contained in $W$.
Consequently its hull and Steiner point satisfy that lemma; a
truncated exterior copy cannot be kept.  Conversely all copies of
obstacles meeting the protected region occur with the required
margin, at every angle.  This gives a finite complete-component
cutoff without an occupation query or a geometric seed supplied to
the controller.

Use angular targets in a rational mesh of maximum gap
$d=c\sqrt e$ around the circle.  There are at most $Ce^{-1/2}$ of
them, and at most $Ce^{-2}$ spatial targets at each.  Rational
approximations to $2\pi$ with a reserved margin give such a mesh;
the actual vectors at all quadrature nodes are rounded with the
certified error $\zeta$.  Apply a union bound over these $Ce^{-5/2}$
targets, setting their failure allowances to
$c\delta e^{5/2}$.  Equations \eqref{eq:single-per-target-count}
and \eqref{eq:single-finite-scales} give the attempted-bit bound in
\eqref{eq:single-finite-resources}.

On the resulting common event, let $\mathcal Z$ be the union of the
estimated Steiner points of retained copies.  Each point lies within
$Ce$ of $c_C(\phi)-a_0$ for some obstacle and the corresponding
sampled normal angle.  Join these points when their distance is less
than $d_0/3$.  Points coming from different obstacles cannot be
joined.  For a complete protected obstacle, consecutive angular
samples of its Gauss parametrization are at distance at most
$r_C^+d+Ce<d_0/3$, so all its samples belong to one component.
A cloud component with a point in the protected radius $R+D_0+1$
comes from an obstacle lying wholly in the larger aperture reserve.
All angular copies of that obstacle were retained.  Components from
outside this cutoff are not used as complete obstacles.  Every
obstacle needed for the period tests is present.

For each retained complete cloud component take its convex hull.
This hull estimates the corresponding $C-a_0$ in Hausdorff distance
$Ce$.  Here the angular mesh $\sqrt e$ suffices, despite the absence
of labels distinguishing the two tangent points.  For every normal
angle $\phi$ choose one sampled Gauss angle $\psi$ at distance at
most $d$.  Since $c_C'(u)=r_C(u)\tau_u$,
\begin{equation}\label{eq:single-gauss-support-deficit}
 0\le h_C(\phi)-c_C(\psi)\cdot n_\phi
    \le r_C^+(1-\cos|\phi-\psi|)
    \le\tfrac12r_C^+d^2.
\end{equation}
Thus the exact sampled boundary hull has support error $O(d^2)$;
the error in each measured center adds $Ce$.  The union of samples,
with no across-angle component labels supplied, gives this hull.

Take any complete copy hull from the angular stage.  Reflect its
Steiner-centered support to estimate $A_0$, also in support supremum
error $Ce$.  All true copy hulls are translates of $-A$ and therefore
give the same centered footprint.  Choosing any one of the finitely
many recovered complete copies creates no matching or asymmetry
assumption.

For completeness, these polygonal support estimates yield the claimed
$C^2$ output without treating an arbitrary polygon as a smooth body.
Convolve each support estimate with the signed approximation kernel
of Section~\ref{sec:stationary}, whose nonconstant moments through
degree six vanish, at scale $h=e^{1/s}$.  For a true $C^{6,\beta}$
support $p$ and its estimate $\widetilde p$ of supremum error $Ce$,
Taylor's formula and the $L^1$ derivative bounds of this fixed kernel
give
\[
 \|L_h*\widetilde p-p\|_\infty\le C(e+h^s),\qquad
 \|L_h*\widetilde p-p\|_{C^2}
                         \le C(eh^{-2}+h^{s-2})\le C\nu.
\]
The smoothed supports also have a uniform $C^3$ bound, since their
third derivative error is at most $C(eh^{-3}+h^{s-3})$.  Evaluate them
on an angular mesh of width $b_1$ comparable to a small multiple of
$\min(e^{1/3},\nu)$, to value error $O(b_1^3)$.  The polygonal
input and fixed kernel have explicit Lipschitz bounds for certifying
these finite integrals.  Quadratic local interpolation with the fixed
smooth partition used in Section~\ref{sec:stationary} then adds
supremum error $O(b_1^3)=O(e)$ and $C^2$ error $O(b_1)=O(\nu)$.
It gives a finite smooth support representation.  Steiner centering
changes the same bounds only by a fixed factor.  The known positive lower curvature
radii ensure $\widehat p+\widehat p''>0$ for small $\nu$, so all
outputs are supports of actual smooth strictly convex bodies.

Apply the retained period-locking and periodic assembly argument to
these original component estimates.  The margin transfers to the
frame $\mathcal O-a_0$ by
Lemma~\ref{lem:registration-period-margin}; it is invariant under a
common translation after reduction by presentation periods.  Matched
period vectors and free area have error $O(e)$, the discrete data
are exact on the common event, and the assembled periodic table has
the asserted geometric loss.  This final algebra uses the acquired
forward-data geometry and no reciprocal measurement.

Finally, each sampled center and displacement vector is specified on
the dyadic mesh $\ell\asymp\zeta$, within a fixed aperture and a
fixed displacement bound.  It therefore has
$O(\log(C/\nu))$ coordinate bits.  The kernels, menus and selection
probabilities have fixed polynomial formulas; certified dyadic scale
enclosures and finite rational arithmetic specify them.  Charging
every selected center-and-vector pair separately gives
$S_\nu\le J_\nu\le N_\nu$, with the description bound stated in
the theorem.  This accounting does not confuse an integral's target
with the many physical nominal sites used to estimate it.
\end{proof}

\begin{corollary}[A finite nonperiodic configuration]
\label{cor:single-finite-cloud}
Assume instead that the configuration has between one and a fixed
number $n_0$ of components, with the same geometric and density
priors, and that a known protected aperture contains all expanded
components $C-A$ together with the fixed collars used above.  The
same forward-only controller recovers $A_0$ and all components of
$\mathcal O_0$ in $C^2$ error at most $C\nu$, with probability at
least $1-\delta$ and the resource bounds
\eqref{eq:single-finite-resources}--\eqref{eq:single-finite-precision}.
Neither periodicity nor a nonperiod-patch margin is needed.
\end{corollary}
\begin{proof}
The protected aperture makes every relevant angular support copy
complete.  The two graph constructions and support smoothing in the
proof of Theorem~\ref{thm:single-law-finite} recover every component.
No period test or periodic assembly is performed.  All estimates and
sampling counts preceding that step remain unchanged.
\end{proof}

\section{The stationary minimax exponent}
\label{sec:sharp-stationary}
The stationary upper and lower powers can be matched. The geometric
point is a cancellation between the two curved sides of a collision
strip. Only boundary arcs for which one translated endpoint is in the
launch disk and the other is outside can contribute to the derivative
of a collision probability. Each such arc, and each swept facet, forces
a cubic lower bound for both the collision area and its complement.
This gives a Hellinger contraction of order $\epsilon^{3/2}$ uniformly
in the command length, including lengths tending to zero. The conditional
information bound and the physical packing then give the stationary
lower power. The launch disk remains fixed and its law may be disclosed
in full; the result does not concern a noise radius tending to zero.

\subsection{A uniform Hellinger bound for short commands}
The cancellation between the physical and swept boundaries remains
effective when the command length tends to zero. We establish the
geometric estimate needed to quantify this cancellation. Throughout this
subsection, $n(\theta)=(\cos\theta,\sin\theta)$ and
$\tau(\theta)=(-\sin\theta,\cos\theta)$.

\begin{lemma}[Effective collision boundary]
\label{lem:effective-collision-boundary}
There is a constant $R_0>0$ with the following property. Fix
$0<R_-\le R_+\le R_0$. For some $\zeta_*>0$, let $C$ be a smooth convex
body whose support function satisfies
\[
                         \|p_C-1\|_{C^2}\le\zeta_*.
\]
For $R\in[R_-,R_+]$, a disk $D=x+R\overline B$, a unit vector $e$,
and $t\ge0$, put
\[
 S=C+[-te,0],\qquad E=S\setminus C,\qquad
 A=\area(D\cap E),\qquad A_c=\area(D)-A.
\]
Write $b=p_Cn+p_C'\tau$, $r=p_C+p_C''$, and
$H_e=\{\theta:n(\theta)\cdot e<0\}$. Let $F_+$ and $F_-$ be the
two facets of $S$ with normals $e^\perp$ and $-e^\perp$, respectively.
Then
\begin{align}
 L_D(C;t,e)
 &: =\int_{H_e}r(\theta)
       \left|\one_D(b(\theta)-te)-\one_D(b(\theta))\right|\dd\theta
       +\mathcal H^1(F_+\cap D)+\mathcal H^1(F_-\cap D)
       \label{eq:effective-boundary-length}\\
 &\le K\min\{A,A_c\}^{1/3}.\label{eq:effective-boundary-area}
\end{align}
Here $K$ is independent of $x,e,t,R$, and $C$ in the stated class.
\end{lemma}
\begin{proof}
Choose $R_0$ sufficiently small and then $\zeta_*$ sufficiently small
relative to $R_-$. All curvature radii lie in $[1/2,2]$. Rotate the
coordinates so that $e=(1,0)$, and parametrize the entering boundary by
arclength $s$. The map
\[
               (s,q)\longmapsto b(s)-qe,\qquad 0<q<t,
\]
parametrizes $E$ injectively off its boundary. Its area Jacobian is
$a(s)=|n(s)\cdot e|$. This follows also by slicing horizontally:
if the left endpoint of $C$ at height $y$ is $l(y)$, the corresponding
collision interval is $[l(y)-t,l(y))$.

We first bound the number of boundary pieces in
\eqref{eq:effective-boundary-length}. The intersection of $\partial C$
with any disk of the stated radius is empty or a single short arc.
Indeed, $b(\theta)=n(\theta)+O(\zeta_*)$. If $b(\theta)\in D$, its
angle lies in an interval of length $O(R+\zeta_*)$ about the direction
of $x$. The whole candidate arc is at distance $O(R+\zeta_*)$ from
$x$. On that arc, with derivatives taken in arclength,
\[
 \frac{\dd^2}{\dd s^2}|b(s)-x|^2
       =2+2(b(s)-x)\cdot b''(s)\ge1.
\]
Its disk sublevel set is therefore an interval. Intersecting this short
arc with the entering hemisphere preserves this property. Applying the
same argument with center $x+te$ shows that the two disk-membership sets
for $b$ and $b-te$ are intervals. Their symmetric difference has at most
two interval components. Each facet contributes at most one chord.
Thus $L_D$ is the sum of at most four lengths.

Consider one active arc $I$, of arclength $\ell$. Its endpoint lying
in $D$ is denoted by $z(s)$. Either $z=b$ and the direction toward the
other endpoint is $-e$, or $z=b-te$ and that direction is $e$.
Write it as $\sigma e$, where $\sigma\in\{-1,1\}$; then
\[
                z(s)\in D,\qquad z(s)+\sigma te\notin D
                         \quad(s\in I),
\]
apart from endpoints of measure zero. Let $J$ be the middle half of
$I$. Its points are at arclength distance at least $\ell/4$ from
either endpoint of the entering hemisphere. The curvature bounds give
\begin{equation}\label{eq:active-entering-angle}
                         a(s)\ge c\ell\qquad(s\in J).
\end{equation}
For such $s$, set
\[
 v=z(s)-x,\qquad B=|v\cdot\tau(s)|,\qquad G=R^2-|v|^2.
\]
The squared-distance function $f(s)=|z(s)-x|^2$ satisfies $f''\ge1$
on $I$. Its values at $s\pm\ell/8$ do not exceed $R^2$, whereas
$f'(s)=2v\cdot\tau(s)$. The two lower Taylor estimates imply
\begin{equation}\label{eq:active-disk-clearance}
                              G\ge c\ell(B+\ell).
\end{equation}
Let $q(s)$ be the distance to the disk boundary along the ray
$z(s)+q\sigma e$. The other endpoint is outside $D$, so $q(s)\le t$.
The disk equation gives
\[
 q(s)=\frac{G}{\sqrt{(v\cdot e)^2+G}+\sigma v\cdot e}
       \ge\frac{G}{2|v\cdot e|+\sqrt G}.
\]
Since $|v\cdot e|\le Ra(s)+B$, the monotonicity of
$G/(c+\sqrt G)$ in $G$, together with
\eqref{eq:active-disk-clearance}, yields
\[
 q(s)\ge c\frac{\ell(B+\ell)}{a(s)+B+\ell},\qquad
                         a(s)q(s)\ge c\ell^2.
\]
For the last inequality use \eqref{eq:active-entering-angle}, considering
$B\le a(s)$ and $B\ge a(s)$. The ray up to $q(s)$ lies in the
collision strip in either of the two activity cases. Its injective
parametrization therefore gives
\begin{equation}\label{eq:active-arc-volume}
                A\ge\int_Ja(s)q(s)\dd s\ge c\ell^3.
\end{equation}
In particular, no division by $t$ is involved.

Next let $I=F_+\cap D$ have length $\ell$, and let $J$ again be its
middle half. The squared distance along the facet has second derivative
$2$. The preceding Taylor argument shows that the Euclidean disk
clearance on $J$ is at least $c\ell^2$. Moreover, $J$ is at distance
at least $\ell/4$ from both endpoints of the whole facet. At inward
depth $w$ below its supporting line, the horizontal displacement of
the entering boundary from its tangency point is at most $C\sqrt w$.
Indeed, the curvature-radius bounds make the vertical drop comparable
to the square of angular distance and the horizontal displacement
comparable to angular distance. For $0\le w\le c\ell^2$, choose $c$
so that this displacement is at most $\ell/8$. The horizontal collision
interval then covers $J$. A rectangle of length $\ell/2$ and inward
depth $c\ell^2$ lies in $D\cap E$, proving
\eqref{eq:active-arc-volume} for the facet chord. The bottom facet has
the identical argument with the inward direction reversed.

We also need the complementary estimate. On the middle half of any
active arc, \eqref{eq:active-disk-clearance} gives Euclidean clearance
\[
                     R-|v|=\frac{G}{R+|v|}\ge c\ell^2.
\]
If $z=b$, the normal strip of inward depth $c\ell^2$ lies in $C$.
If $z=b-te$, the normal strip of outward depth $c\ell^2$ lies outside
$S$, since $b-te$ is its support point with that normal. Both strips
lie in $D\setminus E$. Their normal-coordinate Jacobians are bounded
below uniformly. To justify inward injectivity explicitly, choose a
fixed $q_*>0$ smaller than every curvature radius. The function
$p_C-q_*$ is a support function, so $C=C_0+q_*\overline B$; inward
offsets of depth less than $q_*$ are disjoint parallel boundaries.
Outward injectivity follows from uniqueness of projection onto the
convex set $S$. These strips prove $A_c\ge c\ell^3$.
For a facet chord, its outward rectangle of the same depth lies outside
$S$ and inside $D$, proving the same estimate. Each of the at most four
effective lengths is consequently bounded by
$K\min\{A,A_c\}^{1/3}$. Their sum proves
\eqref{eq:effective-boundary-area}. When $t=0$, both sides are zero.
\end{proof}

\begin{lemma}[Uniform short-command Hellinger contraction]
\label{lem:all-short-hellinger}
Let $C_0,C_1$ belong to the class in
Lemma~\ref{lem:effective-collision-boundary}, and put
$\epsilon=\|p_{C_1}-p_{C_0}\|_\infty$. For the same $D,t,e$, let
\[
 P_i=\frac{\area\bigl(D\cap((C_i+[-te,0])\setminus C_i)\bigr)}
                 {\pi R^2},\qquad i\in\{0,1\}.
\]
Then, uniformly in the nominal center, direction, command length
$t\ge0$, and radius $R\in[R_-,R_+]$,
\begin{equation}\label{eq:all-short-hellinger}
 H^2(\operatorname{Ber}(P_0),\operatorname{Ber}(P_1))
 :=(\sqrt{P_0}-\sqrt{P_1})^2
    +(\sqrt{1-P_0}-\sqrt{1-P_1})^2
                         \le K\epsilon^{3/2}.
\end{equation}
\end{lemma}
\begin{proof}
Interpolate the support functions by
$p_\lambda=(1-\lambda)p_{C_0}+\lambda p_{C_1}$, $0\le\lambda\le1$,
and put $u=p_{C_1}-p_{C_0}$, $r_\lambda=p_\lambda+p_\lambda''$,
$b_\lambda=p_\lambda n+p_\lambda'\tau$. These bodies remain in the
same fixed neighborhood. Write $A(\lambda)$ for their collision areas
in $D$. The normal velocity of both the physical boundary and its
translate is $u(\theta)$. The facets have the corresponding constant
normal velocities. Differentiating the swept area minus the physical
area gives, for almost every $\lambda$,
\begin{align}
 A'(\lambda)
 &=\int_{H_e}u(\theta)r_\lambda(\theta)
     \bigl[\one_D(b_\lambda(\theta)-te)
                  -\one_D(b_\lambda(\theta))\bigr]\dd\theta
          \label{eq:collision-shape-derivative}\\
 &\quad+u(e^\perp)\mathcal H^1(F_{+,\lambda}\cap D)
       +u(-e^\perp)\mathcal H^1(F_{-,\lambda}\cap D).\nonumber
\end{align}
Here $u(\pm e^\perp)$ denotes evaluation at the corresponding normal.
For completeness, the parallel-area bound makes $A$ absolutely
continuous in $\lambda$. The moving curved boundaries have curvature
strictly smaller than that of the disk, and their disk intersections
are the finite arcs described above. Their intersections with
$\partial D$ therefore have arclength zero, as do those of the facets.
The moving-boundary strip formula, followed by dominated convergence,
justifies \eqref{eq:collision-shape-derivative}. The nonentering
curved boundaries cancel because their normal velocities and arclength
densities agree. The entering terms likewise cancel whenever both
corresponding points belong to $D$.

Since $\|u\|_\infty=\epsilon$, the preceding lemma implies, for
$P(\lambda)=A(\lambda)/(\pi R^2)$,
\begin{equation}\label{eq:collision-probability-differential}
 |P'(\lambda)|\le K\epsilon
                   \min\{P(\lambda),1-P(\lambda)\}^{1/3}.
\end{equation}
This inequality includes probabilities arbitrarily close to either
endpoint. Apply the ordinary absolutely continuous chain rule to
$(P+\delta)^{2/3}$ and $(1-P+\delta)^{2/3}$, where $\delta>0$.
Equation~\eqref{eq:collision-probability-differential} bounds each
derivative by $K\epsilon$. Integrate from $0$ to $1$ and let
$\delta\downarrow0$. Thus
\[
 |P_1^{2/3}-P_0^{2/3}|\le K\epsilon,\qquad
 |(1-P_1)^{2/3}-(1-P_0)^{2/3}|\le K\epsilon.
\]
For nonnegative $a,b$, $|a^{3/4}-b^{3/4}|\le|a-b|^{3/4}$.
Applying this inequality to the two preceding estimates, and then
squaring, proves \eqref{eq:all-short-hellinger}. At $t=0$ the
probabilities are both zero.
\end{proof}

The constants are uniform on the stated positive interval of launch
radii. Parameter-independent mixtures of directions preserve
\eqref{eq:all-short-hellinger} by joint convexity of squared Hellinger
distance. A reverse command only translates the nominal disk and
reverses the direction, both already uniform in the lemma.

\subsection{Binary information and physical reconstruction}
\begin{lemma}[Hellinger diameter bounds binary information]
\label{lem:hellinger-capacity}
Let $V$ have any distribution on a finite set and, conditional on
$V=v$, let $Y$ be Bernoulli with mean $p_v$. Put
$a=\min_vp_v$ and $b=\max_vp_v$. Then, in bits,
\begin{equation}\label{eq:hellinger-capacity}
 I(V;Y)\le\frac1{\ln2}
 H^2(\operatorname{Ber}(a),\operatorname{Ber}(b)).
\end{equation}
No positive lower bound on $p_v$ or $1-p_v$ is required. The same
assertion holds under every posterior distribution.
\end{lemma}
\begin{proof}
Put $q=(a+b)/2$. If $q$ is zero or one, the observation is constant.
Otherwise comparison with $Q=\operatorname{Ber}(q)$ gives
\[
 I(V;Y)\le\mathbb E D_2(\operatorname{Ber}(p_V)\Vert Q)
 \le\frac{(b-a)^2}{4\ln2\,q(1-q)}.
\]
Here $D_2$ denotes relative entropy with logarithm to base two.
The first inequality follows by subtracting the nonnegative divergence
of the mixture law from $Q$; the second follows from
$\ln u\le u-1$, which bounds relative entropy by chi-square
divergence. On the other hand,
\begin{align*}
 H^2(\operatorname{Ber}(a),\operatorname{Ber}(b))
 &=(b-a)^2\left\{
 \frac1{(\sqrt a+\sqrt b)^2}
 +\frac1{(\sqrt{1-a}+\sqrt{1-b})^2}\right\}\\
 &\ge\frac{(b-a)^2}{4q(1-q)}.
\end{align*}
Here $(\sqrt a+\sqrt b)^2\le2(a+b)=4q$, and the
complementary inequality is identical. Endpoint cases follow by
continuity. Nothing in the proof restricts the distribution of $V$.
\end{proof}

\begin{theorem}[The stationary lower power for arbitrary short commands]
\label{thm:sharp-stationary-lower}
For every $s=6+\beta$, $0<\beta\le1$, there are fixed physical
priors, a nondegenerate subclass $\mathcal P\subset\mathcal B_\eta$,
a fixed $t_{\max}>0$, and a fixed known uniform-disk launch law
with the following property. Every admissible adaptive controller using
fresh draws of that law, exact nominal centers, and collision bits from
arbitrary directions and lengths $|a|\le t_{\max}$, and achieving
centered component $C^2$ error at most $\nu$ with probability at least
$1-\delta$ uniformly on $\mathcal P$, satisfies
\begin{equation}\label{eq:sharp-stationary-lower}
 \sup_{\mathcal O\in\mathcal P}\mathbb E_{\mathcal O}\mathsf T
                  \ge c\nu^{-(3s/2+1)/(s-2)}
\end{equation}
for fixed $0<\delta<1/4$ and all sufficiently small $\nu$.

The same assertion holds if the controller chooses the known disk
radius in a fixed interval $0<R_-\le R\le R_+$. The interval and
the physical priors are chosen once, with
$t_{\max}+2R_+<d_0$ and with $R_+$ small enough for
Lemma~\ref{lem:all-short-hellinger}. The lattice, primitive orbit
count, species, component correspondences, and parameter-independent
laboratory frame may be disclosed. The actual launch points remain
unobserved. Continuous adaptive directions and arbitrarily small
positive command lengths are included.
\end{theorem}
\begin{proof}
Use the fixed lattice, fixed second obstacle, and support-bump family
of Lemmas~\ref{lem:packing} and \ref{lem:centered-packing}.
Choose the fixed near-circle neighborhood, launch-radius interval, and
strict separation slack first. For all small $h$ the packing lies in
that neighborhood and contains $M=2^m$ tables such that
\begin{align*}
 m&\ge c_0h^{-1},\qquad
 \|p^0_\omega-p^0_{\omega'}\|_{C^2}
                  \ge c_1h^{s-2}\quad(\omega\ne\omega'),\\
 \max_{\omega,\omega'}\|p_\omega-p_{\omega'}\|_\infty
                  &\le C_0h^s.
\end{align*}
The $C^2$ distance from the unit disk is $O(h^{s-2})$.
The laboratory frame and all disclosed data are the same for every
packing parameter; no unknown Steiner point is disclosed.

We specify the component reduction uniformly over the packing.
For a nominal disk $D_R(x)$ and a displacement $a$, every possible
flight lies in the capsule $D_R(x)+[0,a]$, whose diameter is at
most $2R_++t_{\max}$. Relative to one packing member, every other
member's indexed component is within Hausdorff distance $C_0h^s$.
Enlarge the capsule by this amount. For small $h$ its diameter is
still strictly less than $d_0$, so it meets at most one component
of the reference table. Consequently the original capsule can meet
only one common indexed component family across all parameters.
It may miss that family for some parameters. If that component is
fixed, the response law is parameter-independent. If it is a variable
component, Lemma~\ref{lem:all-short-hellinger}, after the common
lattice translation, applies to every pair. Thus every available
Bernoulli response family has Hellinger diameter at most
\begin{equation}\label{eq:sharp-per-attempt-hellinger}
                             C h^{3s/2}.
\end{equation}
All constants are uniform in the nominal center, direction, length,
and chosen launch radius. A reciprocal reverse command is the same
kind of command after translating its nominal center and reversing
its displacement. A parameter-independent mixture of directions
preserves the bound by joint convexity of squared Hellinger distance;
alternatively one may reveal the direction before discarding it.

Condition on the independent controller seed. At every nonterminal
history the posterior is supported on the same finite packing, and
the selected command is fixed by that history and seed.
Lemma~\ref{lem:hellinger-capacity} and
\eqref{eq:sharp-per-attempt-hellinger} bound its conditional
information by $C h^{3s/2}$. Pad the binary record with a fixed
symbol after stopping. The finite chain rule gives
\[
 I(V;Z_{\le n},U_0)
     \le C h^{3s/2}\sum_{k=1}^n\Prob\{\mathsf T\ge k\}
     = C h^{3s/2}\mathbb E\min(n,\mathsf T).
\]
The conditional-entropy convergence and monotone-convergence argument
of Proposition~\ref{prop:stopped-contraction} gives the same
inequality for the complete stopped record, with
$\mathbb E\mathsf T$ on the right. Infinite expected cost already
satisfies the theorem.

Take $h$ to be a sufficiently large fixed multiple of
$\nu^{1/(s-2)}$. The centered-shape accuracy sets are disjoint.
Fano's inequality in its stated stopped-transcript form yields
\[
 C h^{3s/2}\frac1M\sum_\omega\mathbb E_\omega\mathsf T
    \ge (1-\delta)m-h_2(\delta)\ge c h^{-1}.
\]
The supremum dominates the average, proving
\eqref{eq:sharp-stationary-lower}. The fixed-radius assertion is
the special case in which the controller uses one radius throughout.
\end{proof}

\subsection{Shrinking boundary layers and the attempted-bit logarithm}
The fixed-layer bisection of Theorem~\ref{thm:rare-stationary} uses the
final probability accuracy at every level. It remains a valid controller.
The following additional result uses a coarser layer while the radial
interval is wide. A safeguarded interval update permits the layer to
shrink without losing the true boundary from the interval.

\begin{proposition}[A shrinking-layer rare-collision controller]
\label{prop:shrinking-rare}
Under the hypotheses of Theorem~\ref{thm:rare-stationary}, its geometric
and primitive-data conclusions can be achieved with
\begin{align}
 J_\nu&\le C\nu^{-1/(s-2)}\log(C/\nu),
                                  \label{eq:shrinking-centers}\\
 N_\nu&\le C\nu^{-((\gamma+3/2)s+1)/(s-2)}
                                      \log(C/(\nu\delta)).
                                  \label{eq:shrinking-attempts}
\end{align}
The assertion includes both the calibrated one-scale experiment and
the two known scales with unknown footprint, with their respective
fixed compass-separation hypotheses. Nominal centers require only the
same finest precision $O(\nu^{s/(s-2)})$. The launch spread, direction
law and fixed aperture are unchanged.
\end{proposition}
\begin{proof}
Reserve failure probability $\delta/2$ for the fixed-accuracy acquisition
of Lemma~\ref{lem:rare-coarse-normals}. We first describe one radial
search on this successful coarse event. Translate its interior center
to the origin, and write
\[
 \overline B_r\subset P\subset\overline B_{R_1},
 \qquad R=R_P(u)
\]
for its known uniform radial bounds and unknown boundary radius in the
chosen unit direction $u$. Fix a dyadic constant $A\ge4R_1/r$.
For all sufficiently small $e>0$, the inball inclusions
\begin{equation}\label{eq:shrinking-radial-conversion}
 (1-\rho/r)P\subset P_{-\rho},\qquad
 P+\rho\overline B\subset(1+\rho/r)P
                      \quad(0<\rho<r)
\end{equation}
show the following. If a nominal point within distance $e/2$ of $au$
is used in Proposition~\ref{prop:rare-query} with physical layer $e$,
its label is correct whenever $|a-R|>Ae$, except on the conditional
failure event of that proposition. Indeed $a<R-Ae$ places $au$ in
$P_{-2e}$ by the first inclusion with $\rho=2e$, whereas
$a>R+Ae$ places it outside $P+2e\overline B$ by the second.
The rounding reserve preserves depth at least $e$ in the first case
and an exterior point in the second. Thus, on a successful test,
\begin{equation}\label{eq:shrinking-label-invariant}
 \text{label }1\ \Longrightarrow\ a\le R+Ae,
 \qquad
 \text{label }0\ \Longrightarrow\ a\ge R-Ae.
\end{equation}
The fixed tubular neighborhood of the original coarse bracket supplies
the same valid candidate normal for these rounded points.

Take an initial certified interval $[l_0,u_0]$ containing $R$, with
fixed width $w_0>0$. Coarse accuracy and numerical reserves may be
chosen so that this width is a common dyadic prior constant for all
radial searches, and so that $w_0/(4A)$ lies below all collar and
mass-lemma thresholds. Suppose that a current interval $[l_j,u_j]$
contains $R$ and has width $w_j$. At its midpoint $a_j$ query with
\begin{equation}\label{eq:shrinking-layer-schedule}
            e_j=\frac{w_j}{4A},\qquad q=\frac34.
\end{equation}
For label one retain $[a_j-Ae_j,u_j]$; for label zero retain
$[l_j,a_j+Ae_j]$. Equation~\eqref{eq:shrinking-label-invariant}
proves that the retained interval still contains $R$. It is a subinterval
of the preceding interval and has width exactly $q w_j$, whichever
label occurs. Consequently
\[
                     w_j=w_0q^j.
\]
This deterministic width schedule holds also on an erroneous record;
only the containment assertion uses successful labels.

Put $h\asymp\nu^{1/(s-2)}$ and choose a positive dyadic final radial
tolerance $E\asymp c h^s$. Let $K$ be the least integer with
$w_K\le2E$.
For all sufficiently small $E$, one has $K\ge1$ and
\begin{equation}\label{eq:shrinking-final-width}
                   2qE<w_K\le2E,
             \qquad K\le C\log(C/E).
\end{equation}
The final midpoint has radial error at most $E$. Let $M_h\le C h^{-1}$
be a deterministic bound on the total number of radial searches,
including both scales when present. At level $j=0,\ldots,K-1$ of
each search use conditional failure allowance
\begin{equation}\label{eq:shrinking-confidence}
 \varepsilon_j=\frac{\delta}{4M_h}\,2^{-(K-1-j)}.
\end{equation}
The sum for one search is less than $\delta/(2M_h)$, so all fine
tests succeed, conditional on coarse success, with probability at least
$1-\delta/2$. This allocation permits relatively larger failure
allowances at the more expensive fine levels. Locations, labels and
the resulting subintervals may depend on the full past; fresh
preparations give the required conditional bounds at every level.

Write $\kappa=\gamma+3/2$. By Proposition~\ref{prop:rare-query},
the attempts in one search are at most
\begin{align*}
 C\sum_{j=0}^{K-1} e_j^{-\kappa}
                                  \log(C/\varepsilon_j)
 &\le C\sum_{j=0}^{K-1} w_j^{-\kappa}
          \left\{\log(CM_h/\delta)+(K-1-j)\log2\right\}\\
 &\le C w_K^{-\kappa}\log(CM_h/\delta).
\end{align*}
For the last inequality put $k=K-j$ and use the convergent sums
$\sum_{k\ge1}q^{\kappa k}$ and
$\sum_{k\ge1}kq^{\kappa k}$. Equation~\eqref{eq:shrinking-final-width}
therefore bounds the total fine cost by
\[
                         C M_h E^{-\kappa}\log(CM_h/\delta).
\]
The coarse cost is $C\log(C/\delta)$ and is absorbed by this bound.
Substitution of $M_h\le C h^{-1}$ and $E\asymp c h^s$ proves
\eqref{eq:shrinking-attempts}. There are at most two centers per
level and $M_hK$ levels in all; the coarse stencil count is fixed.
This proves \eqref{eq:shrinking-centers}.

All bracket endpoints, layer widths and updates are dyadic when $A$
is chosen as a power of two. For an explicit implementation, choose
one dyadic target mesh so that each nominal point is computed within
$E/(8A)$ of its ideal radial point. Since
$w_{K-1}>2E$, this error is less than $e_j/2$ at every level.
For sufficiently small $E$ it also lies within the fixed bracket
neighborhood of Lemma~\ref{lem:rare-coarse-normals}. Add the exact
dyadic compass shift after rounding the target once. Thus the preceding
radial conversion and normal certification apply to every physical
command. Certified conservative integer batch sizes preserve the
probability and cost bounds. The finest coordinates and all batch
descriptions require only
$O(\log(C/\nu)+\log(C/\delta))$ digits.

Finally, the radial errors are $O(h^s)$ on the common successful event.
The interpolation, support subtraction, convexity and period-locking
steps of Theorem~\ref{thm:rare-stationary} apply without change. At
two known scales, include all searches in $M_h$, use the same uniform
constants, and apply its complete-component correspondence and common
footprint reconstruction. This gives the stated output conclusions
with total failure probability at most $\delta$.
\end{proof}


\begin{corollary}[The stationary minimax power]
\label{cor:stationary-minimax}
Fix a bounded class $\mathcal B_0\subset\mathcal B_\eta$ with the
priors of Theorem~\ref{thm:rare-stationary} and containing the preceding
packing, a confidence
$0<\delta<1/4$, and the same known uniform-disk laws. Let
$N^*_{\rm pool}(\nu,\delta)$ be the infimum of the worst-case
expected number of attempts among controllers with uniform geometric
error at most $\nu$ and success probability at least $1-\delta$,
using the fixed pooled compass and one fixed radius. Define
$N^*_{\rm short}(\nu,\delta)$ in the same way, allowing arbitrary
directions, lengths at most $t_{\max}$, and the fixed radius interval.
Choose the pooled compass length below $t_{\max}$ and with
$2t+2R_+<d_0$. Then, with
$q_0=(3s/2+1)/(s-2)$,
\begin{equation}\label{eq:stationary-sharp-bracket}
 c\nu^{-q_0}
 \le N^*_{\rm short}(\nu,\delta)
 \le N^*_{\rm pool}(\nu,\delta)
 \le C\nu^{-q_0}\log\frac{C}{\nu\delta}.
\end{equation}
In particular both observation designs have the same stationary
minimax power at fixed confidence. The statement leaves the
logarithmic factor and its confidence dependence unresolved.
\end{corollary}
\begin{proof}
Theorem~\ref{thm:sharp-stationary-lower} gives the first inequality:
its centered-shape loss is weaker than the geometric loss here.
The second is inclusion of the command classes. The rare-collision
upper construction, with the decreasing-resolution allocation of
Proposition~\ref{prop:shrinking-rare}, gives the final inequality.
That construction has a deterministic attempt cap, which is also a
bound for the expected-cost criterion used in this corollary.
Thus the lower and upper assertions refer to the same accuracy,
confidence, physical class, and expected-cost criterion.
\end{proof}

The density-independent upper theorem applies to a larger nuisance
class than the fixed disk experiment used for this converse. The
minimax conclusion concerns $\gamma=0$ and fixed positive disk
spreads with the stated geometry. It does not assert the necessity
of the upper exponent for every $\gamma>0$, nor does it make its
constants uniform when a permitted spread tends to zero.

\section{Information categories and relation to earlier inverses}
\label{sec:comparison}
The principal theorem recovers the obstacle configuration and the
complete stationary launch density from a directional short-flight
germ. Its finite version regularizes this identity and prices the
actual forward collision bits. The fixed-known-disk minimax result
provides a separate benchmark. The appendices retain the pooled
occupation and homothetic constructions, with their distinct data and
hypotheses. We describe the classical ingredients and the additional
identification steps for each experiment.

\subsection{Directional boundary measures and the unknown density}
First variation of translated-set probabilities has established
geometric antecedents. Galerne~\cite{Galerne} relates directional
covariogram derivatives to directional variation and perimeter, and
Kiderlen and Rataj~\cite{KiderlenRataj} study first variations of
dilation volume and contact distributions of stationary random sets.
Lemma~\ref{lem:single-boundary-flux} derives the spatially weighted
entering-boundary measure for the present attempted-collision law,
including a weak finite-length error that uses only the total mass
of the unknown density.

The planar angular differential inverse is also classical. In the
sine-transform normalization, Louis, Riplinger, Spiess and
Spodarev~\cite[Section 4]{LRSS} give $f=(g+g'')/2$ for the even
angular datum. Our one-sided kernel has the normalization
\eqref{eq:single-cosine-distribution}. It leaves the sum of two
antipodal normal contributions. Angular inversion alone therefore
does not identify a nonsymmetric obstacle or an unknown launch law.
The spatial variable supplies the additional information: the positive
angular field consists of translated density copies. Their component
supports and masses identify the common footprint and density; their
Steiner points trace the obstacle boundaries. The minimum-area
selection in Theorem~\ref{thm:single-resolved-obstacle} extends the
exact inverse to configurations with only one sufficiently large
obstacle. The proof classifies the full geometric and probabilistic
fiber and determines the complete collision law from its first-order
germ.

The finite inverse uses positive regularization and places derivatives
on known test functions, as in classical approximate inversion.
The collision-specific quantitative steps are the weak flight bias,
density-uniform spatial quadrature, the averaged rare-event variance,
and separated-copy acquisition. The resulting bound is an explicit
sufficient cost for unknown-law reconstruction. The sharp common-disk
power of Section~\ref{sec:sharp-stationary} is stated on its own
observation and nuisance class.

\subsection{Unknown probes and global response rigidity}
Geometric inference with an unknown probe has a close precedent in
scanned-probe microscopy. Villarrubia~\cite{Villarrubia} expresses
probe imaging through morphological dilation and develops blind
reconstruction from image features. Its consistent-probe envelope
is an important comparison for an unknown Minkowski summand.
In the retained homothetic construction, the extra collision flux
identifies footprint width, and two translated settings give exact
support-envelope cancellation. The resulting conclusion identifies the entire
geometric ambiguity, without asymmetric species, correspondence or a
periodic presentation.

The stopped Poisson representation and discrete obstacle iteration
use classical potential theory \cite{LawlerLimic}. Their role in
Theorem~\ref{thm:global-occupation} is to recover the positive set
without a supplied component or killing domain. Uniformly bounded
separated components give the stopping budget, and support
cancellation yields the exact equality of period groups. This
whole-field statement differs from the finite protected-patch
period tests; its finite-field locality bound explicitly records
how the observation region grows.

\subsection{Boundary estimation and informative rare responses}
Active boundary estimation is an essential comparison for the finite
theorem, not merely an analogy with passive billiard spectra.
Castro and Nowak~\cite{CN} study adaptive binary classification under
boundary-fragment smoothness and noise assumptions, with matching
orders for their excess-risk criterion. Their input selects a feature
point whose response is a class label. This directly available label
is not our collision bit: a solid start and a free miss have the same
output here. Neither their observation model nor their loss identifies
our finite periodic table. We use an independent physical packing and
binary-tree proof rather than importing a classification lower bound.

Locatelli, Carpentier and Kpotufe~\cite{LCK} give an adaptive strategy
for smooth decision boundaries without prior knowledge of smoothness
and noise parameters. Their membership-query setting and local
line-search/interpolation construction make the closest algorithmic
comparison. In contrast, our smoothness is known, our error is a $C^2$
geometry norm, and the boundary label must first be derived from
pooled collision data. We claim neither a new general
bisection principle nor an advance on adaptation to unknown smoothness.
The reductions in Propositions~\ref{prop:query} and
\ref{prop:rare-query} permit this active-estimation strategy in the
physical collision experiment, at their respective costs.

Yan, Chaudhuri and Javidi~\cite{YCJ} use informative abstention
responses in active learning. Under their assumptions the rarity of
an informative response can yield an inverse-probability cost. Their
observation includes an abstention symbol. Our sensor has only its
original collision bit. The additional argument here is that translated
pooled commands and a conjunction over outward compass candidates
realize the requisite one-sided event for each boundary query. Small
occupation alone would not justify such an improvement: the individual
means in a signed reciprocal difference need not be small. The exterior
zero and interior lower bound are therefore proved for the commands
that are actually sampled.

Weighted cap mass is also classical. Brunel~\cite{BrunelCaps} gives
the exponent $\gamma+(d+1)/2$ for rolling convex supports and
boundary-decaying densities, and translates cap assumptions into
Hausdorff rates for random convex hulls of observed spatial points.
In the plane the geometric exponent is $\gamma+3/2$. The present
upper bound must realize such information from hidden-launch bits.
The sharp lower bound additionally controls the actual swept-minus-solid
response: its effective-boundary cancellation is uniform for every
short direction and for lengths approaching zero. A cap probability
alone would not prove that all such commands have the required
Hellinger contraction.

\subsection{Localization and control precision}

The earlier scalar inverse reconstructed a uniformly accurate occupation
on a full two-dimensional grid. Its constructive upper bound was
$C\nu^{-3}\log(C/(\nu\delta))$ attempts on the same bounded
$C^{6,\beta}$ class, using $\sigma\asymp\nu^{3/2}$ and compensated
support smoothing. That theorem remains in the exact repository archive of A2 v31. The retained adaptive localized
procedure uses $O(h^{-1})$ angular lines, a fixed-depth dependency
stencil for each queried point, and the full $C^{6,\beta}$ regularity.
It has a different adaptive command design, not a sharper analysis of
the old uniform design. We have not proved a matching lower bound
for that old design or for uniform spatial preparations.

The resource improvement is accompanied by the explicit localization
and calibration cost in \eqref{eq:precision}. In that localized experiment the exact data are a pair
of functionals on controlled spatial inputs. The lower bound concerns
only their finite binary realizations; it does not assert a low-information
passive invariant. The angle/time/position tolerances shrink with
localization and are not learned from the collision outcomes. Apparatus
motion, precision in representing input locations, and arithmetic
cost are distinct resources. Goldberg and Kwek~\cite{GK} treat
query-coordinate precision as an explicit resource in learning convex
polytopes by exact membership queries. The digital bounds in
Theorems~\ref{thm:unk-scale-finite} and \ref{thm:registration-finite}
follow the same accounting principle
for a different class and observation: it includes the additional
integral-measurement centers as well as the fine angular boundary
queries. It does not convert a bit-description bound into a physical
metrology or travel bound.

For completeness the exact v31 manuscript is preserved unchanged
inside the repository archive of the reviewed v33 source. It retains the arbitrary centered-law and general
nontrivial bounded-law inverses, the different-zero-set comparison,
the full fixed-aperture proof, rational interval enclosures, the
uniform-grid construction, and all earlier moment, marked-law and
count-fiber results in their original nested volumes. The more general
randomized law is an exact functional theorem; the finite local query
proved here uses the four-atom compass law. No continuous transition
law is discretized without an error budget. Rational intervals remain
conditional on forcing and survival bounds, not certificates of the
apparatus or physical prior.

The finite periodic conclusion also retains the known positive patch
margin and the bounded periodic presentation. Exact rational period
relations are distinct from approximate Euclidean generators. Removing
the known margin gives only pointwise eventual recognition, as described
in Section~\ref{sec:periods}. The exact full-field period criterion in
Theorem~\ref{thm:global-response-rigidity} has no periodicity prior,
and Corollary~\ref{cor:registration-finite-cloud} treats a fully
captured finite configuration without a period margin. The direction-resolved forward germ of
Section~\ref{sec:single-law-rigidity} has prescribed nominal positions
and directions. Uniformly prepared count germs, analytic network
records, marked length spectra and passive trajectory laws are
different observations. The hypotheses of each inverse specify
which of these data are available.

\subsection{Stationary noise versus adversarial implementation}
The retained fixed-footprint construction belongs also to support estimation with
contaminated data. Brunel, Klusowski and Yang~\cite{BKY} observe
continuous random vectors drawn from an unknown convex body and then
contaminated by additive Gaussian noise. Their geometric estimator
avoids Fourier inversion and is analyzed in Hausdorff loss. Here no
random vector or direct inside/outside response is observed: the data
are attempted collision bits at controlled centers. The fixed compact
noise support, its boundary-mass lower bound and the reciprocal reduction
are essential to our polynomial upper bound. The Gaussian observation
and its tail-based analysis are not covered by this theorem, and its
rates are not being compared as rates in a common experiment.

Support addition, inner rolling disks and the Brunn--Minkowski inequality
are classical \cite{Schneider}. Their role in the retained occupation construction is to remove an
unknown stationary density without learning that density. The positivity set of
the recovered blur, not its detailed values, identifies the footprint
sum. The retained signed high-accuracy query uses the Green bound for exit from a fixed
rectangle; its constants can be large and its finite-depth work grows
logarithmically. The area normalization uses a finite adjoint of the
same killed operator. The direct collision-strip normalization instead
uses Cavalieri's principle and width additivity; only one integrated
pooled forward mean is required. We do not claim these convex-geometric
or random-walk
principles as new abstract theorems.

The localized and stationary upper bounds cannot be ranked by their
attempted-bit exponents alone. The localized model attains the smaller
sample power but requires shrinking spread and worst-case tolerance.
The stationary model keeps its random spread fixed and allows the
density to be unknown. Its uniform-disk minimax power is sharp by
Corollary~\ref{cor:stationary-minimax}; the comparison uses the same
physical class, confidence and expected-cost criterion, with a
stronger command class allowed for the lower bound. For the rare-query
upper, $2t+\max_i\diam A_i<d_0$ is displayed explicitly; the older
signed-query upper needs only $t+\max_i\diam A_i<d_0$.
The retained homothetic geometric inverse reconstructs the centered footprint,
ratios and relative setting translations without a supplied common
origin. Exact homothety, labelled settings, per-setting stationarity
and nominal positioning remain part of its observation model.
Quantitative boundary-mass, geometry, scale-gap and coarse offset
bounds enter its uniform finite version. The common-disk converse
does not extend automatically to every stationary density or to
vanishing launch spread, and digital command counts do not assign
manufacturing or metrology costs.

\appendix
\section{Global recovery of occupation and translation periods}
\label{sec:global-response}
The exact inverse admits a formulation without a periodic presentation.
Let $\mathcal G(D,d)$ consist of the nonempty, locally finite unions
$\mathcal O=\bigcup_{C\in\mathcal C}C$ of compact convex bodies with
nonempty interiors, diameters at most $D$, and distances between
distinct bodies at least $d>0$. The number of bodies and their shapes
need not be finite modulo translations. No boundary smoothness is
assumed in this section. Let $A$ be a compact convex body with nonempty
interior and diameter at most $\Delta$, and let $j$ be a probability
density supported on $A$ and positive almost everywhere in its
interior. The position of $A$ is unrestricted and its support function
need not be given. Fix the compass step so that
\begin{equation}\label{eq:global-step}
                         t+\Delta<d.
\end{equation}
Forward and reciprocal reverse preparations have the same meaning as
in Section~\ref{sec:setting}, with stationary displacement $Z\sim j$.
Write their exact mean fields as $F,R$, and put
\begin{equation}\label{eq:global-fields}
 v(x)=\int_A j(z)\one_{\mathcal O}(x+z)\dd z,\qquad
 g=F-R,\qquad
 Tf(x)=\frac14\sum_{e\in\{\pm e_1,\pm e_2\}}f(x+te).
\end{equation}
The short-segment identity gives $g=(T-I)v$. It uses no periodicity:
a segment of length $t<d$ cannot meet two different solids, and the
forward-minus-reverse bit is its endpoint occupation difference.
The positivity argument of Lemma~\ref{lem:noise-support} likewise gives
\begin{equation}\label{eq:global-positive-components}
 \{v>0\}=\bigcup_{C\in\mathcal C}\operatorname{int}P_C,
 \qquad P_C=C+(-A).
\end{equation}
These components have diameter at most $D+\Delta$ and mutual distance
at least $d-\Delta>t$. Also $0\le v\le1$. Translation continuity of
$j$ in $L^1$ implies uniform continuity of $v$, even when the family
of solids is infinite.

\begin{lemma}[Uniform exit from a positive component]
\label{lem:global-exit}
Let $X_k$ be the virtual compass walk with transition operator $T$,
started at $x\in\operatorname{int}P_C$, and let $\tau_C$ be its first
exit from that interior. Put
\begin{equation}\label{eq:global-exit-constants}
 H=\frac{(D+\Delta+t)^2}{t^2},\qquad b=\lceil2H\rceil.
\end{equation}
Then $v(X_{\tau_C})=0$ almost surely,
$\E_x\tau_C\le H$, and
\begin{equation}\label{eq:global-exit-tail}
                   \Prob_x(\tau_C>N)\le2^{-\lfloor N/b\rfloor}.
\end{equation}
The same estimates hold for every starting point and positive
component in the stated class. At a point where $v(x)=0$ set
$\tau_C=0$.
\end{lemma}
\begin{proof}
A step of length $t$ cannot pass from one positive component to
another. Hence the exit point has zero occupation, including when it
lies on the boundary. The process
$|X_k-x|^2-t^2k$ is a martingale. Before exit its distance from $x$ is
at most $D+\Delta$, and the exit step increases this bound by at most
$t$. Optional stopping at $n\wedge\tau_C$ therefore gives
$t^2\E_x(n\wedge\tau_C)\le(D+\Delta+t)^2$.
Monotone convergence proves the asserted mean bound and almost sure
exit. Markov's inequality bounds survival for $b$ steps by $1/2$.
At each surviving multiple of $b$, the same bound applies from the
current point in the same component. The Markov property proves
\eqref{eq:global-exit-tail}.
\end{proof}

\begin{theorem}[A global scalar-to-occupation inverse]
\label{thm:global-occupation}
Within the preceding class, $g$ determines $v$ uniquely, without
knowledge of the footprint or of the positive components. More
precisely, let $\mathcal T_H(x)$ be the stopping times of the virtual
compass walk, including zero, with $\E_x\tau\le H$. Then
\begin{equation}\label{eq:global-stopping-inverse}
 v(x)=\sup_{\tau\in\mathcal T_H(x)}
            \E_x\left[-\sum_{k=0}^{\tau-1}g(X_k)\right].
\end{equation}
In particular, two admissible occupation fields with forcing $g$ and
$\widetilde g$, possibly arising from different configurations,
footprints and densities satisfying the same bounds, obey
\begin{equation}\label{eq:global-inverse-stability}
                  \|v-\widetilde v\|_\infty
                         \le H\|g-\widetilde g\|_\infty.
\end{equation}
An explicit inverse is obtained from
\begin{equation}\label{eq:global-obstacle-iteration}
 v_0=0,\qquad v_{N+1}=\max\{0,Tv_N-g\}.
\end{equation}
It satisfies, on the whole plane,
\begin{equation}\label{eq:global-obstacle-error}
             0\le v-v_N\le2^{-\lfloor N/b\rfloor}.
\end{equation}
\end{theorem}
\begin{proof}
Since $(T-I)v=g$, the process
$v(X_n)-\sum_{k<n}g(X_k)$ is a martingale. For an integrable stopping
time, stopping first at $n\wedge\tau$ and then passing to the limit
is justified by boundedness of $v,g$ and integrability of $\tau$.
Thus
\[
 v(x)=\E_xv(X_\tau)
             -\E_x\sum_{k<\tau}g(X_k)
       \ge-\E_x\sum_{k<\tau}g(X_k).
\]
If $v(x)>0$, the exit time in Lemma~\ref{lem:global-exit} belongs to
$\mathcal T_H(x)$ and attains equality. If $v(x)=0$, the zero stopping
time does. This proves \eqref{eq:global-stopping-inverse}. The set of
stopping rules and the budget $H$ depend only on the stated bounds.
Each payoff changes by at most $H\|g-\widetilde g\|_\infty$ when
the forcing changes. Taking suprema proves
\eqref{eq:global-inverse-stability}, and also proves uniqueness.

Finite-horizon conditioning shows that \eqref{eq:global-obstacle-iteration}
is the value of the same payoff maximized over $\tau\le N$.
Consequently $0\le v_N\le v$. The particular rule
$\tau_C\wedge N$ gives
\[
 v(x)-v_N(x)\le\E_xv(X_{\tau_C\wedge N})
                     \le\Prob_x(\tau_C>N).
\]
The exit-tail estimate proves \eqref{eq:global-obstacle-error}.
This is the classical stopped Poisson argument, applied with the
unknown positive components discovered by the obstacle iteration;
the martingale and killed-walk identities are discussed in
\cite{LawlerLimic}.
\end{proof}

\begin{corollary}[Finite-field locality]
\label{cor:global-field-locality}
The value $v_N(x)$ depends only on the forcing at lattice points
reachable from $x$ in at most $N$ compass steps. If values
$\widetilde g$ at these points have absolute error at most $\epsilon$,
and $\widetilde v_N$ is computed by the same finite iteration, then
\begin{equation}\label{eq:global-field-locality}
 |v(x)-\widetilde v_N(x)|
             \le2^{-\lfloor N/b\rfloor}+N\epsilon.
\end{equation}
\end{corollary}
\begin{proof}
Reachability follows inductively from the transition rule. Averaging
and the positive-part map are nonexpansive in the supremum norm, so
the error in the $N$th iterate is at most $N\epsilon$. Combine this
with \eqref{eq:global-obstacle-error}.
\end{proof}

For a field or collection of fields $f$, let
$\operatorname{Per}(f)=\{w:f(x+w)=f(x)\text{ for all }x\}$, with
equality required in every component of a collection.

\begin{theorem}[Response rigidity for the full period group]
\label{thm:global-response-rigidity}
For every configuration and stationary law above,
\begin{equation}\label{eq:global-period-identity}
 \operatorname{Per}(F,R)=\operatorname{Per}(g)
      =\operatorname{Per}(v)=\operatorname{Per}(\mathcal O).
\end{equation}
Thus one exact reciprocal difference field determines the full
translation period group even when the footprint and its density are
unknown. This group is discrete. In particular, the configuration has
a rank-two lattice of periods, with finitely many component orbits,
if and only if the scalar difference field has two independent periods.
No periodicity or bounded-presentation hypothesis is needed for this
equivalence.
\end{theorem}
\begin{proof}
Every period of $\mathcal O$ is a period of both collision mean
fields, by translation covariance of the experiment. Every common
period of $F,R$ is a period of $g$. The inverse in
Theorem~\ref{thm:global-occupation} commutes with translations; hence
every period of $g$ is a period of $v$.

A period $w$ of $v$ preserves its positive set. The separated closures
of the positive components are precisely the bodies $P_C$, so
translation by $w$ permutes them. If $P_C+w=P_{C'}$, equality of
laboratory support functions gives, for every unit vector $u$,
\[
 p_C^{\rm lab}(u)+p_{-A}^{\rm lab}(u)+w\cdot u
       =p_{C'}^{\rm lab}(u)+p_{-A}^{\rm lab}(u).
\]
Cancellation yields $C+w=C'$. Applying the inverse translation proves
that $w$ is a period of $\mathcal O$, completing the equality chain.

A nonzero period cannot fix one compact component under translation.
It therefore takes that component to a distinct component, and its
length is at least $d$. The period group is consequently discrete.
If it contains two independent vectors, it is a rank-two lattice;
local finiteness and bounded component diameters imply that only
finitely many component orbits meet a compact fundamental cell.
Conversely such a periodic configuration has two independent periods,
which are inherited by $g$ through \eqref{eq:global-period-identity}.
\end{proof}

The exact datum here is a scalar response field on the plane. The
locality estimate specifies its finite-region approximation; it does
not replace a statistical count of the attempted bits used to estimate
those field values. Uniform finite period decisions from a fixed
protected aperture are a separate conclusion and retain the bounded
presentation and positive patch margin specified in
Section~\ref{sec:setting}. The present theorem identifies crystallinity
from the exact active response without assuming crystallinity of the
underlying configuration.

\section{A finite-stencil inverse and its sampling law}
\label{sec:finite-stencil}
The global stopping formula can be evaluated on a fixed finite set of
nominal centers. In particular, neither its spatial aperture nor its
forcing precision needs to grow with an iteration horizon. The finite
inverse is an optimal stopping linear program on a known killed walk.
Its killing boundary is computational; no assertion that the physical
occupation vanishes on that boundary is made.

Throughout this section the configuration and launch law satisfy the
hypotheses of Section~\ref{sec:global-response}. Put
\begin{equation}\label{eq:stencil-constants}
 K=\max\{1,\lceil(D+\Delta)/t\rceil\},\quad
 S=\{-K,\ldots,K\}^2,\quad M=(2K+1)^2,\quad
 H_S=K^2+(K+1)^2.
\end{equation}
For a target $x$, the nominal stencil is $x+tS$. Let $Q$ be the
substochastic compass matrix on $S$: a step leaving $S$ is killed and
its mass is not redistributed. Write $o=(0,0)$, let $e_o$ be its unit
coordinate vector, and set
\begin{equation}\label{eq:stencil-green}
 G=(I-Q)^{-1},\qquad h=G\one,\qquad
 \mathcal M_o=\{m\ge0:(I-Q^{\mathsf T})m\le e_o\}.
\end{equation}
Existence of $G$ and compactness of $\mathcal M_o$ are included in the
next lemma. For a real vector $f$ on $S$ define
\begin{equation}\label{eq:stencil-functional}
 \mathcal V(f)=\max_{m\in\mathcal M_o}(-f^{\mathsf T}m).
\end{equation}
The matrix, feasible set, and all their coefficients are independent
of the obstacles, the launch density, and the data.

\begin{lemma}[The killed stopping polytope]
\label{lem:stencil-polytope}
The killed compass walk exits $S$ almost surely, and
$0<h_z\le H_S$ for every $z\in S$. The set $\mathcal M_o$ is a
nonempty compact polytope. Its elements are precisely the expected
continuation-visit vectors of randomized stopping rules that stop no
later than exit from $S$. Moreover,
\begin{equation}\label{eq:stencil-budget}
 0\le m_z\le G_{oz},\qquad \sum_zm_z\le h_o\le H_S.
\end{equation}
For every real $f$ one has the attained dual formula
\begin{equation}\label{eq:stencil-dual}
 \mathcal V(f)=
 \min\{u_o:u\ge0,\ (I-Q)u\ge-f\}.
\end{equation}
Equivalently, $\mathcal V(f)$ is the value at $o$ of the unique
solution $u=\max\{0,Qu-f\}$.
\end{lemma}
\begin{proof}
Use lattice units and let $\tau_S$ be the first exit time. The process
$|X_n|^2-n$ is a martingale. At $n\wedge\tau_S$ the squared norm is
at most $K^2+(K+1)^2$. Bounded stopping, followed by monotone
convergence, gives
$\E_z\tau_S\le H_S-|z|^2\le H_S$. Thus exit is almost sure,
$Q^n\to0$, and the nonnegative Green series $\sum_{n\ge0}Q^n$
converges to $G$. Its row sums are the expected exit times.

If $m_z$ counts expected continuations at $z$ and $s_z$ counts
expected stops there, conservation of visits gives
\[
                 m+s=e_o+Q^{\mathsf T}m,\qquad m,s\ge0.
\]
Conversely, given these equations, continue at $z$ with probability
$m_z/(m_z+s_z)$, interpreted as zero when the denominator vanishes.
The resulting continuation transition matrix is bounded entrywise by
$Q$ and is therefore transient. Its visit equations have a unique
solution, namely $m+s$, so its continuation and stopping masses are
exactly $m$ and $s$. This proves the asserted realization of every
feasible vector; history-dependent rules satisfy the same conservation
identities.

Multiplication by the nonnegative matrix $G^{\mathsf T}$ yields
\[
             m=G^{\mathsf T}(e_o-s)\le G^{\mathsf T}e_o.
\]
Summation proves \eqref{eq:stencil-budget}, which also proves
compactness. The zero vector is feasible. The dual feasible set is
nonempty because $u=\|f\|_\infty h$ is feasible. Finite-dimensional
linear programming duality gives \eqref{eq:stencil-dual} and attainment.
The visit-measure formulation is the usual total-cost stopping
formulation; see \cite{Altman}. The argument above states its complete
finite-state realization in the present experiment.

For completeness, finite-horizon stopping gives the increasing
iteration $u_0=0$, $u_{n+1}=\max(0,Qu_n-f)$. Its values are bounded
by $\|f\|_\infty h$ and converge to the stopping value. The Bellman
map satisfies the componentwise estimate
$|B(a)-B(b)|\le Q|a-b|$. Hence any two fixed points differ in
absolute value by at most $Q^n$ times their original difference.
Transience proves uniqueness. Its value at $o$ equals
\eqref{eq:stencil-functional} by the visit representation.
\end{proof}

\begin{theorem}[Exact fixed-stencil occupation recovery]
\label{thm:finite-stencil}
For every target $x\in\R^2$, put $f_x(z)=g(x+tz)$, $z\in S$.
Then
\begin{equation}\label{eq:finite-stencil-inverse}
                         v(x)=\mathcal V(f_x).
\end{equation}
Thus precisely the same finite-dimensional functional recovers the
occupation at every point from a prior-bounded stencil. The positive
set, footprint, density, and containing component are not inputs.
No periodicity or boundary regularity is required.

The functional is defined also for data that are not the response of
any physical experiment. For arbitrary real vectors $f,\widetilde f$,
\begin{equation}\label{eq:stencil-stability}
 |\mathcal V(f)-\mathcal V(\widetilde f)|
 \le\sum_{z\in S}G_{oz}|f_z-\widetilde f_z|
 \le h_o\|f-\widetilde f\|_\infty.
\end{equation}
In particular, arbitrary measurement errors incur no factor equal to
an iteration horizon.
\end{theorem}
\begin{proof}
By the preceding lemma, $\mathcal V(f_x)$ is the supremum of
$\E_x[-\sum_{k<\tau}g(X_k)]$ over rules stopping by exit from
$x+tS$. Every such time has finite expectation. The martingale
identity used in Theorem~\ref{thm:global-occupation} gives
\[
 -\E_x\sum_{k<\tau}g(X_k)
                  =v(x)-\E_xv(X_\tau)\le v(x).
\]
If $v(x)=0$, stopping immediately attains equality. Otherwise $x$
belongs to one expanded component $P_C$. Every point of $P_C$ lies
within distance $D+\Delta\le Kt$ of $x$. Consequently the first exit
from its interior occurs no later than exit from $x+tS$ along the
compass walk. At that first exit occupation is zero, since one step
cannot reach another positive component. This admissible stopping rule
attains $v(x)$, proving \eqref{eq:finite-stencil-inverse}.

Other components may meet the edge of the stencil, and occupation
there may be positive. The inequality for every stopping rule and
the equality for the containing-component exit are precisely why no
physical zero boundary condition is needed.

For every feasible $m$, \eqref{eq:stencil-budget} bounds the change
of its objective by $\sum_zG_{oz}|f_z-\widetilde f_z|$. Taking maxima
in both directions proves \eqref{eq:stencil-stability}.
\end{proof}

\begin{corollary}[Local scalar stability]
\label{cor:stencil-local-stability}
Two admissible experiments with the same stated geometric bounds,
but possibly different configurations, footprints, and densities,
satisfy
\begin{equation}\label{eq:stencil-local-stability}
 |v(x)-\widetilde v(x)|
       \le h_o\max_{z\in S}|g(x+tz)-\widetilde g(x+tz)|.
\end{equation}
For a fixed experiment and every vector $a$,
\begin{equation}\label{eq:stencil-period-defect}
 |v(x+a)-v(x)|
       \le h_o\max_{z\in S}|g(x+a+tz)-g(x+tz)|.
\end{equation}
These are local estimates on fixed nominal apertures, not whole-field
error assumptions.
\end{corollary}
\begin{proof}
Apply the same functional at the two targets or experiments and use
\eqref{eq:stencil-stability}. The stencil matrix does not change under
translation.
\end{proof}

\subsection{Finite observations and numerical certificates}
The inverse gives both a finite sample law and a directly checkable
interval for every numerical output. If
$|f_z-\widehat f_z|\le\xi_z$ with $\xi_z\ge0$, then
\begin{equation}\label{eq:stencil-certified-interval}
 \mathcal V(\widehat f+\xi)\le\mathcal V(f)
                      \le\mathcal V(\widehat f-\xi).
\end{equation}
This follows because $m\ge0$ in every objective. Any feasible primal
vector supplies a lower bound and any feasible dual vector supplies
an upper bound in \eqref{eq:stencil-dual}. Thus finite rational
feasibility checks and a primal--dual gap certify the computed interval;
an unverified numerical optimizer status is unnecessary.

\begin{theorem}[A fixed-aperture finite occupation experiment]
\label{thm:stencil-sampling}
Fix $0<\varepsilon<1$ and $0<\delta<1$. At any nominal target $x$,
there is a procedure using at most $M$ distinct nominal sites and
\begin{equation}\label{eq:stencil-attempt-count}
 N\le 2M\left\lceil
       8H_S^2\varepsilon^{-2}\log\frac{4M}{\delta}
                         \right\rceil
\end{equation}
attempted collision bits, whose output has error at most
$\varepsilon$ with probability at least $1-\delta$. All sites lie
in $x+[-Kt,Kt]^2$, independently of the target accuracy. The result
is uniform over the unknown stationary density and convex footprint
in the exact class, without a quantitative boundary-mass assumption.
It remains valid conditionally at a target chosen from previous data.

There are $2M$ command-labelled batches, one forward and one reverse
per site; every batch repetition is included in $N$. Numerical
arithmetic is separate from these observation counts. If $x$ and $t$
are rational, all sites and both linear programs have rational data;
no progressively finer set of nominal positions is needed for this
occupation query.
\end{theorem}
\begin{proof}
At each site use $n$ fresh attempts for each mean, where $n$ is the
ceiling in \eqref{eq:stencil-attempt-count}, and let $\widehat f$
be the empirical forward-minus-reverse vector. Hoeffding's inequality
and a union bound give
\[
 \Prob\left\{\max_z|\widehat f_z-f_z|>
                      \frac{\varepsilon}{2H_S}\right\}
 \le4M\exp\left(-\frac{n\varepsilon^2}{8H_S^2}\right)
 \le\delta.
\]
Indeed it suffices to estimate each of the $2M$ means to
$\varepsilon/(4H_S)$. On this event
\eqref{eq:stencil-stability} bounds the error of
$\mathcal V(\widehat f)$ by $\varepsilon/2$. Compute its value to
certified error at most $\varepsilon/2$ using the finite rational
primal and dual programs, then project the result onto $[0,1]$.
Projection cannot increase the error because $v(x)\in[0,1]$.

Every recorded bit, including a solid start or a free miss, is an
attempt in the displayed count. Fresh preparations give exactly the
same bounds conditionally on previous observations. The stochastic
walk and stopping rule are computational objects and require no
additional physical flights. The rational coefficient and batch-count
bit lengths grow as $O(\log(1/\varepsilon)+\log\log(4M/\delta))$
after the rational target, step and fixed priors are specified.
\end{proof}

\begin{remark}[Aperture, interpretation, and comparison]
\label{rem:stencil-scope}
A bound on the absolute position of the footprint is needed only to
translate the nominal aperture into a fixed physical launch aperture:
forward and reverse starts lie in
$x+[-Kt,Kt]^2+A+t\overline B$. The number $M$ may be large; its
explicit dependence on the ratio $(D+\Delta)/t$ is not suppressed in
\eqref{eq:stencil-constants}. For a list of adaptive targets, allocate
conditional failure probabilities and add their batch costs. This
pointwise occupation theorem does not turn finitely many values into
a whole-field periodicity certificate.

The iteration of Lemma~\ref{lem:stopped} and
Corollary~\ref{cor:global-field-locality} remains valid. The present
finite program removes the horizon-dependent amplification for noisy,
not necessarily physically admissible, forcing. It does not replace
the one-sided local rare event of Proposition~\ref{prop:rare-query}:
estimating a signed occupation value and detecting a physical rare
collision are different statistical operations. In particular the
sharp stationary boundary power of
Corollary~\ref{cor:stationary-minimax} is unchanged.
\end{remark}

\section{Reconstruction with a fixed, unknown stationary launch law}
\label{sec:stationary}
The calibration converse quantifies over implementation maps that may
change with the table, the command and the nominal draw. Stationary
additive noise is a different experiment. We now give a constructive
inverse for this experiment at a fixed, nonzero launch spread. The
probability density need not be known, symmetric or mean zero. Its
support, and a lower bound on its mass near the support boundary, are
calibrated apparatus data.

\subsection{A common launch law and its geometric support}
Fix a known strictly convex body $K\subset\R^2$, containing the origin
in its interior, whose support function is $C^{6,\beta}$ and whose
curvature radii have positive lower and finite upper bounds. Assume
\begin{equation}\label{eq:fixed-footprint}
                       t+\diam K<d_0.
\end{equation}
This strict inequality, the bounds on $K$, and its position in the
laboratory frame are fixed independently of the target accuracy.
Let $\gamma\ge0$ and $b_0>0$ be fixed, with the following density class
nonempty:
\begin{equation}\label{eq:jitter-class}
 \mathcal J=\left\{j\in L^1(K):\ j\ge0,\quad\int_Kj=1,\quad
 j(z)\ge b_0\dist(z,\partial K)^\gamma
                  \text{ for almost every }z\in\operatorname{int}K\right\}.
\end{equation}
For $\gamma=0$ the lower bound is $b_0$ throughout the interior. No
upper bound, derivative bound or evaluation oracle for $j$ is needed.
The law $j$ is fixed during the experiment and is independent of the
commands, their compass direction and the table. The guarantee below
is uniform over this unknown nuisance law.

At nominal center $x$, draw fresh independent $Z\sim j$ and $a\sim\rho$.
The forward preparation attempts $(x+Z,a)$; its reciprocal reverse
attempts $(x+Z+a,-a)$ with the same prescribed joint distribution and
fresh draws. Only the pooled collision bit is returned. In particular,
$Z$, the actual start, free-start status and direction-resolved means
are not observed. Durations, directions and nominal centers are exact
in this stochastic model; certified rounding of the latter is handled
below. There is no additional arbitrary implementation error.

Set
\begin{equation}\label{eq:blurred-occupation}
 v_j(x)=\int_K j(z)\one_{\mathcal O}(x+z)\dd z,\qquad
 P_C=C+(-K).
\end{equation}
The first expression is a reflected convolution when $j$ is not even.
The same reciprocal identity gives $g_j=(T-I)v_j$. Every component
$P_C$ has diameter at most $D_0+\diam K$, and distinct such components
are at distance at least $d_0-\diam K>t$. Indeed, for $c\in C,d\in D$
and $z,w\in K$, one has
$|(c-z)-(d-w)|\ge |c-d|-|z-w|$.

\begin{lemma}[Support recovery without density deconvolution]
\label{lem:noise-support}
For every $j\in\mathcal J$,
\begin{equation}\label{eq:noise-support}
 \{v_j>0\}=\bigcup_C\operatorname{int}P_C,\qquad
             p_C^{\rm lab}=p_{P_C}^{\rm lab}-p_{-K}^{\rm lab}.
\end{equation}
Consequently the exact reciprocal mean functions on a sufficiently
large fixed protected aperture determine all protected component bodies,
without knowing $j$. With exact zero tests the full period group is then
recognized as in Lemma~\ref{lem:recognition}, without a supplied positive
nonperiod margin.
\end{lemma}
\begin{proof}
The intersection $C\cap(x+K)$ has nonempty interior exactly when
$x\in\operatorname{int}(C+(-K))$. One way to see the interior assertion
is to apply the open linear map $(c,z)\mapsto c-z$ to
$\operatorname{int}C\times\operatorname{int}K$; its image is the interior
of the Minkowski difference. On that intersection $j$ is positive
almost everywhere after translation, so its integral is positive.
Outside the interior, the intersection has area zero. The separation
above makes the positive components distinct. Support functions add
under Minkowski addition, which proves the second formula.

The proof of Lemma~\ref{lem:stopped} uses only $0\le v\le1$,
separated positive components of bounded diameter, and $g=(T-I)v$.
It does not require a smooth or known convolution kernel. Applied to
$v_j$, its killed iterates recover $v_j$ uniformly on the protected
region from its exact forcing. This proves the local assertion. The
period statement applies to the recovered original bodies, not to an
assumed lattice of the noise law. These support and Minkowski facts are
classical convex geometry; see \cite{Schneider}.
\end{proof}

\subsection{A uniform lower bound near the blurred boundary}
A quantitative version of positivity is needed for a finite experiment.
For a convex body $A$, write $A_{-r}=\{x:x+r\overline B\subset A\}$.
Whenever $r$ is smaller than its uniform interior rolling radius,
$p_{A_{-r}}=p_A-r$. The identity follows by subtracting $r$ from the
support function: its curvature measure remains nonnegative, and the
resulting Minkowski summand is exactly the erosion.

\begin{lemma}[Overlap and boundary mass]
\label{lem:cap-mass}
There are $c,e_0>0$, depending only on the physical class, $K$, $b_0$
and $\gamma$, such that for $0<e<e_0$ and every protected component,
\begin{equation}\label{eq:cap-mass}
 x\in(P_C)_{-e}\quad\Longrightarrow\quad
              v_j(x)\ge c e^{\gamma+3/2}.
\end{equation}
On the complement of $\bigcup_C\operatorname{int}P_C$ the value is zero.
Both assertions are uniform over $j\in\mathcal J$.
\end{lemma}
\begin{proof}
Choose uniform interior rolling radii $r_C,r_K$ and put $r_P=r_C+r_K$.
Reduce $e_0$ so that $e_0<\min(r_C,r_K,r_P,1)/64$; after the
collar depth $d_*$ below is fixed, require also $e_0<d_*/8$.
Thus every erosion by $e/4$ used below is strictly below the relevant
rolling radius, and $K_{-e/4}$ retains a uniform positive rolling radius.
All subsequent support identities are therefore invoked only within their
valid erosion range.

First consider the unweighted overlap
$w(x)=\area(C\cap(x+K))$. Write $P=C+(-K)$.
At $x_0\in\partial P$ with outer normal $n$, the uniquely determined
support points satisfy $x_0=y-z$, where $y\in\partial C$ has normal
$n$ and $z\in\partial K$ has normal $-n$. Let $r_C,r_K$ be fixed
positive interior rolling radii, smaller if necessary than their
uniform bounds. At $x=x_0-dn$, the intersection contains the lens of
the two disks with centers $y-r_Cn$ and $y+(r_K-d)n$ and radii
$r_C,r_K$. Their centers are separated by $r_C+r_K-d$.

Here is an explicit lower bound for the lens. In tangent-normal
coordinates about $y$, choose a fixed sufficiently small $c_1>0$.
For $|u|\le c_1\sqrt d$, the top of the first disk is
$-r_C+\sqrt{r_C^2-u^2}\ge-d/4$, whereas the bottom of the second is
$r_K-d-\sqrt{r_K^2-u^2}\le-3d/4$. For small $d$ the rectangle
\[
                  |u|\le c_1\sqrt d,\qquad
                         -3d/4\le y_n\le-d/4
\]
lies in both disks. Thus $w(x)\ge c_2d^{3/2}$ throughout a fixed
inner collar of $P$. An interior point at sufficiently small distance
$d$ from $\partial P$ has the displayed form at a nearest boundary
point, so this covers the whole collar, not only prescribed normals.

The bound extends to deeper points. The function $w^{1/2}$ is concave
on $P$: the Minkowski combination of the intersections at $x_1,x_2$
is contained in the intersection at their corresponding combination,
and the planar Brunn--Minkowski inequality gives concavity. Fix a small
collar depth $d_*$. On $\partial P_{-d_*}$ the preceding bound is
$c_2d_*^{3/2}$. A chord through any point of $P_{-d_*}$ has endpoints
on this boundary. Concavity bounds $w$ there by the same lower bound.
It follows that
\begin{equation}\label{eq:unweighted-cap}
 w(x)\ge c_3\min\{\dist(x,\partial P),d_*\}^{3/2},\qquad x\in P.
\end{equation}
The constants are uniform for the small erosions of $K$ used next:
their support norms and rolling radii remain uniformly controlled.

For $x\in P_{-e}$ put $K'=K_{-e/4}$. By the support identities,
$C+(-K')=P_{-e/4}$. Hence $x$ has distance at least $3e/4$ from
the boundary of this smaller sum. Applying \eqref{eq:unweighted-cap}
to $K'$ gives
$\area(C\cap(x+K'))\ge c_4e^{3/2}$. On $K'$ the density lower
bound is $b_0(e/4)^\gamma$. Restricting the defining integral to
this intersection proves \eqref{eq:cap-mass}. For $\gamma=0$ the
same argument works, although the erosion is unnecessary.
\end{proof}


\begin{proposition}[A density-independent inverse modulus]
\label{prop:jitter-modulus}
Consider two physical experiments $(\mathcal O_0,j_0)$ and
$(\mathcal O_1,j_1)$ with $j_0,j_1\in\mathcal J$ and the same calibrated
footprint. Let $g_0,g_1$ be their reciprocal pooled-mean differences and
put $\Delta=\|g_0-g_1\|_{W,\infty}$. On a common protected patch,
for sufficiently small $\Delta$, their complete components admit a
bijection satisfying
\begin{align}
 \max_C d_H(C_0,C_1)&\le C\Delta^{1/(\gamma+3/2)},
                                      \label{eq:jitter-Hausdorff}\\
 \max_C\|p_{C_0}^{\rm lab}-p_{C_1}^{\rm lab}\|_{C^2}
    &\le C\Delta^{(s-2)/(s(\gamma+3/2))}.
                                      \label{eq:jitter-modulus}
\end{align}
The densities need not agree. On $\mathcal B_\eta$, the primitive
orbit count and rational period relations in matched generating lists
also agree below the known
locking thresholds, and matched real period bases and free areas
have error bounded by the right side of \eqref{eq:jitter-Hausdorff}.
\end{proposition}
\begin{proof}
The comparison by stopping at the zero set of the first occupation in
Section~\ref{sec:queries}, and then interchanging the two occupations,
gives
$\|v_{j_0}-v_{j_1}\|_{U,\infty}\le H_K\Delta$.
This uses the bounded separated positive components $P_C$ and not
equality of the two densities. Put
$e=(2H_K\Delta/c)^{1/(\gamma+3/2)}$, with $c$ from
Lemma~\ref{lem:cap-mass}. For a point of $P_{C_0}$ choose a point of
$(P_{C_0})_{-e}$ at distance at most $e$; the erosion identity supplies
one. Its first occupation is at least $2H_K\Delta$, hence its second
occupation is positive. The two unions of expanded components are
therefore at Hausdorff distance at most $e$, by symmetry. Protective
collars include the test points for all components under comparison.

For $e$ below a fixed fraction of the expanded-component separation,
each connected expanded component lies in the $e$-neighborhood of a
unique component of the other union. The neighborhoods are pairwise
disjoint. Applying the same argument in reverse gives a bijection of
complete components. In support coordinates their distances are at
most $e$. Subtract the same function $p_{-K}$ to obtain
\eqref{eq:jitter-Hausdorff} for the original bodies.

For completeness, on a uniformly bounded $C^{6,\beta}$ class the
interpolation estimate
\[
 \|f\|_{C^2}\le C\|f\|_\infty^{(s-2)/s},\qquad
                              \|f\|_\infty\le1,
\]
follows by convolution with an explicitly signed approximation
kernel. Starting from an even compact smooth probability kernel $\kappa$,
the linear combination of its dilates at scales $1,2,3,4$ with
coefficients $(8/5,-4/5,8/35,-1/35)$ has total mass one and vanishing
nonconstant moments through degree six. It is a signed numerical
approximation kernel, not a launch probability density. At smoothing scale $b$, its second
derivative is bounded by $C b^{-2}\|f\|_\infty$, and the derivative
remainder is $C b^{s-2}$. Choosing $b=\|f\|_\infty^{1/s}$ proves
the estimate; the first and zeroth derivatives have weaker bounds.
Apply this to the differences of matched support functions to prove
\eqref{eq:jitter-modulus}. Lemma~\ref{lem:locking} gives the final
assertion, with a matched independent-pair basis rather than a
continuously selected reduced lattice basis.
\end{proof}

\subsection{High-accuracy local queries at fixed launch spread}
The required query accuracy now tends to zero; a fixed $1/8$ error
would not locate the boundary of the positive set. A refinement of the
killed-aperture estimate avoids charging a spurious factor of the
iteration depth to every probability measurement.

\begin{lemma}[Aperture Green bound for perturbed iterates]
\label{lem:aperture-green}
Choose the fixed observation rectangle $W$ with collars protecting the
complete $P_C$ containing all targets in $U$. There are constants
$H_K,H_W,L_K$ such that the exact killed iterates for $g_j$ have error
at most $2^{-\lfloor n/L_K\rfloor}$ on $U$. If the forcing at each
state has error at most $\xi$ and each update error is at most $\zeta$,
the clipped approximate iterate satisfies
\begin{equation}\label{eq:jitter-green-error}
 \|\widehat V_n-v_j\|_{U,\infty}
       \le2^{-\lfloor n/L_K\rfloor}+H_W(\xi+\zeta).
\end{equation}
At any requested target $x$, all these iterates use only the states
$(x+t\Z^2)\cap W$, whose number is bounded by a fixed prior constant,
independently of $n$ and the target accuracy.
\end{lemma}
\begin{proof}
Take $H_K=(D_0+\diam K+t)^2/t^2$ and
$L_K=\max(1,\lceil2H_K\rceil)$. The stopped positive-component
argument proves the first assertion. For the walk killed on leaving
$W$, put $\tau_W=\inf\{n\ge0:X_n\notin W\}$ and take
$H_W=(\diam W+t)^2/t^2$. The bounded-stopping martingale calculation
gives $\sup_{x\in W}\E_x\tau_W\le H_W$.

Let $E_n$ be the pointwise absolute difference of the noisy and exact
killed iterates on the lattice through $x$. Averaging, positive part
and clipping to $[0,1]$ give
$E_{n+1}\le T^WE_n+\xi+\zeta$, with both errors zero outside $W$.
Iterating from zero yields
\[
 E_n(x)\le(\xi+\zeta)
       \sum_{k=0}^{n-1}(T^W)^k\one_W(x)
   =(\xi+\zeta)\E_x\min(n,\tau_W)\le H_W(\xi+\zeta).
\]
This bound uses exit from the entire fixed rectangle; it does not assume
that the approximate iterate knows the positive set. Combining it with
the exact error proves \eqref{eq:jitter-green-error}. A bounded rectangle
meets at most a fixed number of points of any translate of $t\Z^2$.
Transitions inside that state set are exact additions by the compass
step; transitions leaving it are killed, not wrapped or renormalized.
\end{proof}

\begin{proposition}[An indeterminate-layer query under fixed jitter]
\label{prop:jitter-query}
For each target $x\in U$, accuracy scale $0<e<e_0$ and conditional
failure allowance $0<\varepsilon<1/4$, a finite reciprocal experiment
returns the correct membership label of $\bigcup_CP_C$ whenever the
target lies outside its inner boundary layer of width $e$. In that
layer the label is unrestricted. It uses at most a fixed number of
command centers and
\begin{equation}\label{eq:jitter-query-cost}
                    C e^{-(2\gamma+3)}\log(C/\varepsilon)
\end{equation}
attempted collision bits. Its launch law has the same fixed support
$K$ at every query and every accuracy. No sample of the launch position
or of an occupation indicator is returned.
\end{proposition}
\begin{proof}
Let $b_e=c e^{\gamma+3/2}$ be the lower bound in
Lemma~\ref{lem:cap-mass}, decreasing $c$ so $b_e\le1$.
Estimate $v_j(x)$ to error at most $b_e/8$. In
Lemma~\ref{lem:aperture-green}, choose
$n=O(\log(C/b_e))$ with the first error at most $b_e/32$, and
estimate each pooled mean to error $b_e/(128H_W)$. Forcing is the
difference of two such means. Reserve update error $b_e/(64H_W)$.
These allocations give error less than $b_e/8$. There are only a
prior-bounded number of means, and their sample averages have Hoeffding
tails. A union bound gives the stated cost. Fresh preparations justify
the same assertion conditional on all previous data.

Threshold at $b_e/2$. Outside the positive components the true value is
zero, and at depth at least $e$ it is at least $b_e$. The test is therefore
correct at both kinds of points. Only the layer between them is
unrestricted. Rounding a target once by $O(e)$, followed by exact
addition of the dyadic compass step, enlarges this layer by $O(e)$;
it does not change the probability accuracy requirement. The finite
state iterates may be evaluated by rational arithmetic on the hit
counts or with the stated update reserve.
\end{proof}

\subsection{Uniform finite reconstruction and the role of calibration}
\begin{theorem}[Unknown stationary jitter on a fixed calibrated support]
\label{thm:stationary-jitter}
Fix the bounds of $\mathcal B_\eta$, a footprint $K$ satisfying
\eqref{eq:fixed-footprint}, and the nonempty density class
\eqref{eq:jitter-class}. There are constants $C,c,\nu_0>0$ and a finite
adaptive controller, independent of the unknown density $j$, such that
for $0<\nu<\nu_0$ and $0<\delta<1/4$,
\begin{equation}\label{eq:jitter-uniform-quantifier}
 \inf_{\mathcal O\in\mathcal B_\eta}\ \inf_{j\in\mathcal J}
 \Prob_{\mathcal O,j}\{\text{geometric loss}\le C\nu
              \text{ and the primitive discrete data are correct}\}
                         \ge1-\delta.
\end{equation}
The discrete data and the laboratory loss have the meanings fixed in
Section~\ref{sec:setting}. The controller uses one fixed launch support
$K$, the reciprocal compass law, and only the pooled collision bits.
It satisfies
\begin{align}
 J_\nu&\le C\nu^{-1/(s-2)}\log(C/\nu),\label{eq:jitter-centers}\\
 N_\nu&\le C\nu^{-((2\gamma+3)s+1)/(s-2)}
                  \log(C/\nu)\log(C/(\nu\delta)).\label{eq:jitter-attempts}
\end{align}
Every attempt, including a solid start or a miss, is counted. The
aperture is fixed by the priors. No localization scale of the random
start shrinks with $\nu$. For densities bounded below on $K$
($\gamma=0$), the power in \eqref{eq:jitter-attempts} is
$(3s+1)/(s-2)$.

Nominal centers have dyadic precision $O(\nu^{s/(s-2)})$ and require
$O(\log(C/\nu)+\log(C/\delta))$ digits per batch, including the
repetition count. Certified enclosures of the fixed numerical priors, including $b_0$
and $\gamma$, and certified support and derivative evaluations for $K$
are assumed when a numerical output is requested.
These are description and algorithmic assertions, not bounds on
manufacturing the launch law or calibrating its support.
\end{theorem}
\begin{proof}
The expanded bodies $P_C$ are strictly convex with uniformly bounded
$C^{6,\beta}$ support functions, positive curvature bounds, bounded
diameters and positive separation. Their curvature radii are the sums
of those of $C$ and $-K$. Choose a fixed enlarged protected aperture
for them and all coarse search brackets.

A fixed-scale version of Proposition~\ref{prop:jitter-query} gives
coarse hulls and interior centers for the expanded bodies. The proof of
Lemma~\ref{lem:coarse} applies verbatim to an indeterminate-layer label:
its deep grid nodes are retained and no node outside the expanded body
is retained; allowing an additional layer only increases the fixed
error reserve. The complete-component cutoff excludes fragments and
protects every required component. The isolated radial brackets and
positive inball therefore follow with changed prior constants.

Use $m$ angular directions, $h=2\pi/m\asymp\nu^{1/(s-2)}$, and
query at layer width $e=c h^s$. The relaxed bisection invariant of
Lemma~\ref{lem:bisection} needs only correctness outside a boundary
layer, not an exactly unbiased label or a known density. It returns
radial values of $P_C$ with error $O(e)$ in $O(\log(C/e))$ queries
per direction. Lemma~\ref{lem:radial} then gives smooth convex bodies
$\widehat P_C$ whose laboratory supports have $C^2$ error $O(h^{s-2})$
and supremum error $O(h^s)$.

Subtract the known footprint in support coordinates:
\begin{equation}\label{eq:estimated-subtraction}
                 \widehat p_C= p_{\widehat P_C}^{\rm lab}-p_{-K}^{\rm lab}.
\end{equation}
The true function $p_C+p_C''$ is bounded below by $1/\kappa_+$.
For small $\nu$, the $C^2$ error ensures
$\widehat p_C+\widehat p_C''>0$. Thus \eqref{eq:estimated-subtraction}
is the support function of an actual strictly convex body, not a formal
subtraction of sets. Its $C^2$ error is $O(h^{s-2})$ and its Hausdorff
error is $O(h^s)$. Subtracting a first harmonic for Steiner centering
preserves these orders. If a smooth output rather than a $C^{6,\beta}$
output is desired, use the finite smoothing described below.

Apply Lemma~\ref{lem:locking} to these original component bodies.
Their errors eventually lie below the known patch and bounded-denominator
thresholds, so the primitive orbit count and exact rational relations
are recovered. Real basis vectors and free area remain estimates, with
error $O(h^s)$. Positive separation and curvature yield a physical
periodic output exactly as in the proof of Theorem~\ref{thm:adaptive}.
The test is on the original bodies after subtraction; it is not an
assumption that an arbitrary noisy difference of sets preserves periods.

There are at most $Q_h=C(1+h^{-1}\log(C/h))$ adaptive local queries,
including coarse acquisition. Allocate conditional failure probability
$\delta/Q_h$ to each. Each query has a prior-bounded number of command
centers, even though its killed iteration depth grows logarithmically.
Equation~\eqref{eq:jitter-query-cost} yields total cost
\[
 C h^{-1}e^{-(2\gamma+3)}\log(C/h)\log(C/(h\delta)).
\]
Substituting $e=c h^s$ and $h\asymp\nu^{1/(s-2)}$ proves both counts.
Conditional confidence and the union bound prove
\eqref{eq:jitter-uniform-quantifier}. No property of $j$ other than its
support, lower bound and stationarity was used.

For numerical realization choose a dyadic compass step satisfying the
strict gap and one dyadic target mesh comparable to $e$. The state set
is the intersection of its translated exact compass lattice with a
fixed dyadic rectangle. Rounded radial bisection has the same mesh-floor
bound as in Theorem~\ref{thm:digital}. Sample counts here grow
polynomially in $1/\nu$, but their binary descriptions still require
only logarithmically many digits. The thresholds and count bounds can
be replaced by conservative dyadic enclosures of the fixed constants;
the receiver never needs to evaluate the unknown $j$.

For a finite smooth output, evaluate \eqref{eq:estimated-subtraction}
on an angular grid of width $b\asymp c\min(e^{1/3},\nu)$, to value
accuracy $O(b^3)$. The reconstructed supports have a uniform $C^3$
bound. Quadratic local interpolation patched with the fixed smooth
partition therefore has supremum error $O(b^3)$ and $C^2$ error
$O(b)$: Taylor remainders and value errors contribute the same orders
after up to two derivatives. These are respectively $O(e)$ and
$O(\nu)$. All support evaluations are computed from the finite radial
representation and the known footprint, not obtained through additional
collision observations. The curvature margin persists after reducing
constants. This final step is a disclosed numerical overhead, not part
of the attempted-bit count.
\end{proof}

\begin{corollary}[No fixed-jitter resolution floor in this experiment]
\label{cor:no-jitter-floor}
A fixed nondegenerate footprint satisfying \eqref{eq:fixed-footprint}
permits arbitrarily accurate reconstruction from finite collision bits,
uniformly over \eqref{eq:jitter-class}. In particular, stationary
uniform disk noise of any fixed radius $R$ with $t+2R<d_0$ has no
positive $C^2$ resolution floor in this model.
\end{corollary}
\begin{proof}
Keep $K$ and $j$ fixed and let $\nu$ decrease in
Theorem~\ref{thm:stationary-jitter}. Its sample bound is finite at every
positive accuracy. For a uniform disk, $\gamma=0$ and
$b_0=1/(\pi R^2)$ are admissible. The full density was not otherwise
needed by the construction.
\end{proof}

There is no contradiction with Theorem~\ref{thm:calibration-floor}.
The latter permits different, command- and draw-dependent
implementations in its indistinguishable worlds. Here all commands use
one stationary translate of a density with a calibrated support and
uniform boundary mass. The noise radius is not a worst-case actuator
error budget. An unknown displacement of the entire footprint, an
uncontrolled change of its support with the command, arbitrary errors
in duration or angle, and unbounded noise distributions are not included
in the new theorem. Nominal positioning and support calibration remain
separate apparatus requirements. Nor does the theorem claim a sharp
sample exponent: it proves that repetition can replace narrowing the
random launch spread in this stochastic experiment, not that every
kind of calibration can be replaced by repetition.

\section{Joint recovery of the table and an unknown launch footprint}
\label{sec:unknown-footprint}
The one-scale stationary theorem calibrates the whole convex support of
the launch law. That support can instead be an unknown of the inverse
problem when the same launch mechanism is operated at two known,
nonzero scales. The extra control is a dilation about a fixed laboratory
origin, not a position measurement. Both spreads stay fixed as the
requested accuracy tends to zero. All one-scale results remain valid
without this additional control.

\subsection{The two-scale experiment}
Fix numbers $0<\lambda_1<\lambda_2$, known to the controller. Let
$\mathcal K$ be a class of unknown strictly convex footprints with
\begin{equation}\label{eq:unknown-K-prior}
 r_K\overline B\subset K\subset R_K\overline B,\qquad
 \|p_K^{\rm lab}\|_{C^{6,\beta}}\le M_K,\qquad
 0<\rho_K^-\le p_K^{\rm lab}+(p_K^{\rm lab})''\le\rho_K^+.
\end{equation}
The constants are fixed and known, but neither $p_K$ nor its Steiner
point is supplied. The scale gap $\lambda_2-\lambda_1$ is fixed and
positive. Require
\begin{equation}\label{eq:unknown-K-separation}
                      t+2\lambda_2R_K<d_0.
\end{equation}
For each $K$, an unknown stationary density $j$ has integral one on
$K$ and satisfies
\begin{equation}\label{eq:unknown-K-density}
 j(z)\ge b_0\dist(z,\partial K)^\gamma\quad\text{a.e. in }K,
               \qquad b_0>0,\quad \gamma\ge0.
\end{equation}
Take the admissible family of pairs $(K,j)$ nonempty. An attempt at
scale $\lambda_i$ and nominal center $x$ uses $x+\lambda_i Z$,
where $Z$ is a fresh draw from $j$. The reverse attempt uses
$x+\lambda_i Z+a$ and $-a$. The compass direction is pooled as
before. Only a collision bit is returned; no realization of $Z$, impact
position, free-start status or direction-resolved outcome is supplied.
The same unknown $j$ is used throughout. Its scaled density obeys
\begin{equation}\label{eq:scaled-density-bound}
 j_i(z)=\lambda_i^{-2}j(z/\lambda_i)
 \ge b_0\lambda_i^{-(\gamma+2)}
                    \dist(z,\partial(\lambda_iK))^\gamma.
\end{equation}
Nominal centers, durations, directions and the two scale factors are
exact controls, apart from the numerical rounding specified below.
There is no additional arbitrary displacement error.

Write $v_i(x)=\int_Kj(z)\one_{\mathcal O}(x+\lambda_i z)\dd z$
and $g_i=(T-I)v_i$. The reciprocal identity and the stopped comparison
recover $v_i$ from $g_i$ on protected regions. Lemma~\ref{lem:noise-support}
gives their complete positive components as
\begin{equation}\label{eq:two-scale-components}
                 P_{i,C}=C+(-\lambda_iK),\qquad i=1,2.
\end{equation}
The controller knows uniform geometric bounds for these components,
although it does not know their support functions. Their radii of
curvature are sums, their diameters are bounded by $D_0+2\lambda_2R_K$,
and their mutual separation is at least $d_0-2\lambda_2R_K>t$.
These are all the geometric bounds needed before footprint subtraction
in the proof of Theorem~\ref{thm:stationary-jitter}.

\begin{lemma}[Matching scales without component labels]
\label{lem:scale-matching}
Put $m_K=(\lambda_2-\lambda_1)r_K$ and
$d_K=d_0-2\lambda_2R_K$. The small-scale component belonging to
$C$ satisfies
\begin{equation}\label{eq:scale-nesting}
                 P_{1,C}+m_K\overline B\subset P_{2,C}.
\end{equation}
Consequently the two collections in \eqref{eq:two-scale-components}
determine their common component correspondence without obstacle labels,
asymmetry or distinct-shape assumptions.

More precisely, suppose complete convex estimates at each scale have
Hausdorff error at most $e$, under their true within-scale correspondences.
If $e<c\min(m_K,d_K)$ for a fixed sufficiently small $c$, the Steiner
point of each estimated small-scale component lies in exactly its
corresponding estimated large-scale component. This rule recovers the
same correspondence from the estimates.
\end{lemma}
\begin{proof}
Convexity and $0\in K$ give
$\lambda_2K=\lambda_1K+(\lambda_2-\lambda_1)K$.
The latter summand contains $m_K\overline B$, proving
\eqref{eq:scale-nesting}. In particular every point of $P_{1,C}$
is at distance at least $m_K$ from the complement of $P_{2,C}$.
The distinct large-scale components are $d_K$-separated, so containment
pairs each small component with a unique large one. Every large one
contains its own small component, proving bijectivity.

The Steiner point lies in its convex body, and its change under
Hausdorff error $e$ is at most $2e$, by \eqref{eq:centering-map}.
Thus the estimated small-scale Steiner point has distance at least
$m_K-2e$ from the complement of the true large-scale component.
For compact convex bodies at Hausdorff distance $e$, support inequalities
give $A_{-e}\subset\widehat A\subset A+e\overline B$. Hence the
point belongs to the corresponding large estimate when $3e<m_K$.
It belongs to no other large estimate when $3e<d_K$. These inequalities
supply a common positive margin for certified containment tests.
No unstable choice of a nearest congruent shape is involved.
\end{proof}

All collections in this lemma consist of complete components. In a
bounded experiment take the protective collars large enough to contain
both dilations of every component needed for the period tests. The
coarse complete-component cutoff of Lemma~\ref{lem:coarse}, with the
larger diameter bound above, discards exterior fragments. Components
needed in the protected patch therefore occur at both scales; unmatched
exterior fragments are not used to infer a missing component or a period.
The enlarged aperture still depends only on the fixed priors.

\begin{theorem}[Recovery under calibrated homothety and an inverse modulus]
\label{thm:self-exact}
The exact pair of reciprocal mean functions $(g_1,g_2)$ in the
two-scale experiment determines the unknown footprint in its laboratory
frame and every protected original component. For a matched pair,
\begin{align}
 p_{-K}^{\rm lab}
   &=\frac{p_{P_{2,C}}^{\rm lab}-p_{P_{1,C}}^{\rm lab}}
                   {\lambda_2-\lambda_1},\label{eq:self-footprint}\\
 p_C^{\rm lab}
   &=\frac{\lambda_2p_{P_{1,C}}^{\rm lab}
                  -\lambda_1p_{P_{2,C}}^{\rm lab}}
                   {\lambda_2-\lambda_1}.\label{eq:self-body}
\end{align}
With exact zero tests the full period group is also determined, under
the bounded periodic presentation and without a positive patch margin.

For two admissible experiments with possibly different $K,j$ and tables,
put $\Delta=\max_i\|g_i-\widetilde g_i\|_{W,\infty}$.
For small $\Delta$, the protected components admit a correspondence with
\begin{align}
 d_H(K,\widetilde K)+\max_Cd_H(C,\widetilde C)
       &\le C\Delta^{1/(\gamma+3/2)},\label{eq:self-modulus-0}\\
 \|p_K^{\rm lab}-p_{\widetilde K}^{\rm lab}\|_{C^2}
  +\max_C\|p_C^{\rm lab}-p_{\widetilde C}^{\rm lab}\|_{C^2}
       &\le C\Delta^{(s-2)/(s(\gamma+3/2))}.
                                                   \label{eq:self-modulus-2}
\end{align}
The constants include the reciprocal of the fixed scale gap. On
$\mathcal B_\eta$, primitive discrete data lock below the retained
thresholds, and matched real period bases and free areas obey the
Hausdorff order in \eqref{eq:self-modulus-0}.
\end{theorem}
\begin{proof}
The stopped inverse first identifies the two positive-component
collections; Lemma~\ref{lem:scale-matching} identifies their pairing.
Minkowski support addition gives
$p_{P_{i,C}}^{\rm lab}=p_C^{\rm lab}+\lambda_i p_{-K}^{\rm lab}$.
Solving these two linear equations proves
\eqref{eq:self-footprint}--\eqref{eq:self-body}. The answer for $K$
is the same from every pair because the experiment uses one common
footprint. Reflection of a support argument by $\pi$ recovers $p_K$
from $p_{-K}$. This also recovers a nonzero Steiner point of $K$;
no mean-zero convention or hidden centering of the launch law was used.
Apply Lemma~\ref{lem:recognition} to the recovered original bodies.

For the modulus, the stopped comparison of Section~\ref{sec:queries}
uses only uniformly bounded, separated positive components. Thus it
compares $v_i$ and $\widetilde v_i$ even for different footprints and
different densities. Its supremum error is at most $H_*\Delta$.
The proof of Lemma~\ref{lem:cap-mass} is uniform over
\eqref{eq:unknown-K-prior}--\eqref{eq:scaled-density-bound}, so an
inner depth $e=C\Delta^{1/(\gamma+3/2)}$ has occupation greater
than $2H_*\Delta$. The erosion argument gives Hausdorff distance
$O(e)$ between the positive components at each scale, with a bijection
of complete protected components. Uniform $C^{6,\beta}$ support bounds
and the signed-kernel interpolation estimate give $C^2$ error
$O(e^{(s-2)/s})$ at each scale. This step does not subtract a supposedly
common unknown footprint.

The nesting margin makes the two within-scale correspondences consistent
with the cross-scale pairing. Indeed the small-scale Steiner point and
its counterpart remain in the corresponding large-scale bodies below
the same margin; a different pairing would put one point in two
separated large components. Now subtract the two explicit linear systems.
Their inverse has norm bounded in terms of
$(\lambda_2-\lambda_1)^{-1}$, which proves both estimates, for $K$
as well as the original bodies. Finally apply Lemma~\ref{lem:locking}
to those original bodies, after subtraction, not to the expanded union.
\end{proof}

\begin{theorem}[Finite two-scale reconstruction without a footprint oracle]
\label{thm:self-finite}
Fix the bounds of $\mathcal B_\eta$, \eqref{eq:unknown-K-prior},
the strict margin \eqref{eq:unknown-K-separation}, the two scale factors,
and $b_0,\gamma$. One finite adaptive controller, independent of $K$
and $j$, reconstructs a strictly convex footprint $\widehat K$ and a
smooth strictly convex periodic table. For all sufficiently small
$\nu>0$ and $0<\delta<1/4$, uniformly over admissible $\mathcal O,K,j$,
its geometric loss and footprint laboratory $C^2$ support error are at
most $C\nu$, and its primitive discrete data are correct, with
probability at least $1-\delta$. Its costs satisfy
\begin{align}
 J_\nu&\le C\nu^{-1/(s-2)}\log(C/\nu),\label{eq:self-centers}\\
 N_\nu&\le C\nu^{-((2\gamma+3)s+1)/(s-2)}
                \log(C/\nu)\log(C/(\nu\delta)).\label{eq:self-attempts}
\end{align}
Both random launch supports are fixed independently of $\nu$.
No footprint support or derivative evaluations and no density evaluations
are inputs. Nominal positioning is refined to
$O(\nu^{s/(s-2)})$; the exact scale factors and their common dilation
origin, the footprint envelope and regularity bounds, and the boundary-mass
constants are calibrated controls or priors. The theorem does not charge
their physical manufacture or calibration to $N_\nu$.
\end{theorem}
\begin{proof}
Run the construction of Theorem~\ref{thm:stationary-jitter} separately
at the two scales, stopping each construction before its footprint
subtraction. Every earlier step uses only occupation queries, the
component bounds and the lower boundary-mass constants; it never evaluates
$p_K$. Equation~\eqref{eq:scaled-density-bound} supplies known lower
constants at both scales. The aperture Green bound, coarse acquisition,
indeterminate-layer bisection and radial interpolation therefore give
complete convex estimates $\widehat P_{i,C}$ with support supremum
error $O(e)$ and $C^2$ error $O(h^{s-2})$, where
$h\asymp\nu^{1/(s-2)}$ and $e=c h^s$.

Use Lemma~\ref{lem:scale-matching} to pair the two finite collections.
The estimates have a common error event of probability at least
$1-\delta$ after splitting confidence between scales and queries.
On this event every pairing needed in the protected patch is correct.
For the finitely many complete pairs take the average of
\[
 q_C=\frac{p_{\widehat P_{2,C}}^{\rm lab}
                 -p_{\widehat P_{1,C}}^{\rm lab}}{\lambda_2-\lambda_1}
\]
and call it $\widehat q$. It estimates $p_{-K}^{\rm lab}$ in supremum
norm by $O(e)$ and in $C^2$ by $O(\nu)$; averaging cannot enlarge
the maximum error. Put
$\widehat p_C=p_{\widehat P_{1,C}}^{\rm lab}-\lambda_1\widehat q$.
These functions have the same two error orders for $p_C^{\rm lab}$.
Using one average makes the output footprint common to all components.
Uniform positive lower radii of curvature of $K$ and $C$ imply
$\widehat q+\widehat q''>0$ and
$\widehat p_C+\widehat p_C''>0$ for small $\nu$. Hence all these
estimates represent actual strictly convex bodies, not arbitrary formal
support differences. The known inball of $K$ persists after a fixed
reduction of its radius. The finite final smoothing of
Section~\ref{sec:stationary}, applied also to $\widehat q$, preserves
the stated orders and strict curvature.

Apply the complete-component period tests and rational locking of Lemma~\ref{lem:locking}
to the original-body estimates. All margins are met after reducing
$\nu_0$, giving a physical periodic output with the same primitive orbit
count and exact rational relations. Matched real basis vectors and free
area have error $O(e)$ and in particular $O(\nu)$. The footprint itself
is not subjected to period tests.

There are two families of radial queries. Put
$Q_h=C(1+h^{-1}\log(C/h))$. Each query uses at most
\[
                   C e^{-(2\gamma+3)}\log(CQ_h/\delta)
\]
attempts, by the fixed-aperture estimate. This proves
\eqref{eq:self-centers}--\eqref{eq:self-attempts}. Separate fresh
attempts at both scales justify conditional confidence allocation; no
coupling of the physical noise draws is assumed. Exact dyadic compass
addition, rounded radial bisection and certified arithmetic use the same
error reserves as before. The numerical constants require only the
finite priors in the theorem, not an oracle for the unknown footprint.
\end{proof}

\begin{remark}[What has, and has not, been calibrated]
\label{rem:self-contract}
This theorem replaces knowledge of an entire function $p_K$ by two
known homothetic settings and uniform geometric and mass bounds. It
does not learn the nuisance density. Exact homothety is substantive:
independently drifting unknown footprints are not covered. If the two
settings have an additional common unscaled laboratory offset $b$, the
formulas instead recover the translated table $\mathcal O-b$ and the
same $K$. Absolute table location then needs an external reference.
This follows directly from
$C-b+(-\lambda_iK)$ in \eqref{eq:two-scale-components}, and is not
a new noise floor for centered shape or for periods.
\end{remark}

\section{Boundary tests from rare pooled collisions}
\label{sec:rare-stationary}
The preceding stationary construction estimates a difference of two
collision probabilities to the accuracy of a small boundary mass.
Near an exposed boundary one can instead arrange that an exterior
test has a zero collision probability. An interior test then needs only
one observed collision. The compass direction remains unobserved.
To cover the directions at which two compass vertices maximize the same
support functional, we test a finite set of candidate vertices and take
the conjunction of their outcomes.

The compass length in this section is chosen once from the fixed
priors. Let $D_*$ be a known upper bound for the diameter of every
available launch footprint, and require
\begin{equation}\label{eq:rare-margin}
                         2t+D_*<d_0.
\end{equation}
For the calibrated one-scale experiment one may take $D_*=\diam K$;
for the two known scales of Section~\ref{sec:unknown-footprint}, take
$D_*=2\lambda_2R_K$. The extra factor of two accounts for both the
shift of a nominal center and the attempted flight. This sufficient
design condition is stronger than the one used for the earlier
reciprocal estimates. Those estimates remain valid under their stated
conditions. Whenever the footprint diameter bound is strictly below
$d_0$, a positive dyadic $t$ satisfying \eqref{eq:rare-margin} can be
chosen independently of the requested accuracy. No new output or
direction control is required.

\subsection{Coarse normal information}
The expanded components $P=C+(-K)$ have a uniform positive lower
curvature radius, a uniform support regularity bound, and a positive
separation. Fix a collar width $d_c>0$ smaller than their lower rolling
radius and such that
\begin{equation}\label{eq:rare-collar}
            d_c<t/8,\qquad D_*+2t+d_c<d_0.
\end{equation}
Reducing $d_c$ below a fixed fraction of these bounds will be harmless.
For every $y$ in this collar of a complete expanded component, there
are a boundary point $y_0$, its outward unit normal $n$, and a signed
depth $d\in[-d_c,d_c]$ such that
\begin{equation}\label{eq:rare-signed-depth}
              y=y_0-dn,\qquad n\cdot y=p_P^{\rm lab}(n)-d.
\end{equation}
Here $d>0$ inside $P$. The representation is unique in the chosen
collar: the support function $p_P-d$ still has a positive curvature
radius for every inward offset in this range, and describes the inner
parallel boundary. Exterior metric projections are unique by convexity.

\begin{lemma}[Fixed-accuracy normals for radial brackets]
\label{lem:rare-coarse-normals}
At any prescribed confidence $1-\varepsilon_c$, a prior-dependent
finite number of the earlier reciprocal queries supplies an interior
center for each protected expanded component and isolated radial
brackets containing its boundary in every radial direction. Each bracket
has a tubular neighborhood of one fixed positive prior-controlled radius
lying in the collar \eqref{eq:rare-signed-depth}. For each radial
direction the same finite data compute a rational vector $\widehat n$
satisfying
\begin{equation}\label{eq:rare-normal-error}
                         |\widehat n-n|\le1/32
\end{equation}
for the projection normal of every point in that tubular neighborhood.
The number
of attempts is at most $C\log(C/\varepsilon_c)$. For unknown
footprints the construction uses only their uniform priors.
\end{lemma}
\begin{proof}
Apply the fixed-layer version of Proposition~\ref{prop:jitter-query}
and the coarse acquisition of Lemma~\ref{lem:coarse} to the expanded
components, obtaining interior centers and isolated radial intervals.
With respect to each center, take a fixed angular mesh of width $h_c$
and refine its radial values to error $e_c\le c h_c^s$ by the same
reciprocal queries. All these parameters are fixed from the priors.
Lemma~\ref{lem:radial} yields a finite smooth radial function
$\widetilde R$ with
\[
 \|\widetilde R-R_P\|_\infty\le C h_c^s=:E_c,
 \qquad
 \|\widetilde R-R_P\|_{C^1}\le C h_c^{s-1}.
\]
Choose $h_c$ small enough that both error bounds meet the fixed
reserves below. In radial direction $u(\varphi)$, use a certified
interval containing
$[\widetilde R(\varphi)-2E_c,\widetilde R(\varphi)+2E_c]$.
Its endpoints enclose the true radial boundary, and its width can be
made arbitrarily small at this fixed stage. It remains in the original
isolated interval after reducing $h_c$ and its numerical reserve.

The normal to a positive radial graph is
\[
 \frac{R_P(\varphi)u(\varphi)
                 -R_P'(\varphi)u^\perp(\varphi)}
        {\sqrt{R_P(\varphi)^2+R_P'(\varphi)^2}}.
\]
The inball and regularity bounds make this expression uniformly
Lipschitz in $R_P,R_P'$. Substitution of $\widetilde R$ therefore
approximates the true radial boundary normal to any prescribed fixed
accuracy. The projection normals throughout the bracket have the same
property. To verify this last assertion quantitatively, let $r_*>0$
be a uniform lower rolling radius of $P$. Write $P=Q+r_*\overline B$,
using the support identity from Section~\ref{sec:stationary}. If
$x_i\in\partial P$ has outer normal $n_i$, the monotonicity of the
support-point relation for $Q$ gives
\[
 (x_1-x_2)\cdot(n_1-n_2)\ge r_*|n_1-n_2|^2,
 \qquad |n_1-n_2|\le |x_1-x_2|/r_*.
\]
The projection of any collar point lies within twice its distance
from the true radial boundary point. Hence on a tubular neighborhood
of radius $a_c$ about the bracket its normal differs from the radial
boundary normal by $O(E_c+a_c)$. Choose $h_c$ and then a fixed
$a_c>0$ so small that the whole neighborhood lies in the chosen collar
and that this error and the radial approximation error sum to less
than $1/64$. Compute a rational normal approximation with another
error less than $1/64$. This proves \eqref{eq:rare-normal-error}
uniformly on that neighborhood; its radius is independent of all
subsequent accuracy parameters.

Only a fixed number of radial values, bisections and reciprocal command
centers was used. All requested probability accuracies are fixed, so
their conditional confidence allocation costs
$C\log(C/\varepsilon_c)$ attempts. Every step precedes footprint
subtraction and depends only on the uniform bounds for the expanded
bodies. It consequently applies at either unknown-footprint scale.
\end{proof}

\subsection{A test using only positive collision records}
Put $E=\{e_1,-e_1,e_2,-e_2\}$ and, for the rational vector supplied
by Lemma~\ref{lem:rare-coarse-normals}, form the finite set
\begin{equation}\label{eq:rare-candidates}
 V(\widehat n)=
 \left\{v\in E:\ \widehat n\cdot v
       \ge\max_{w\in E}\widehat n\cdot w-\tfrac14\right\}.
\end{equation}
All these comparisons are rational. This set contains every maximizer
of $n\cdot v$ over $E$: the two normal errors cost at most $1/16$.
Moreover, every candidate satisfies
\begin{equation}\label{eq:rare-positive-normal}
 n\cdot v\ge\frac1{\sqrt2}-\frac14-\frac1{16}>\frac38.
\end{equation}
The threshold in \eqref{eq:rare-candidates} is positive, so the set
contains at most two compass vertices. In particular the construction
does not require a unique maximizing direction.

\begin{proposition}[A rare-collision boundary query]
\label{prop:rare-query}
Assume \eqref{eq:rare-margin} and condition on the successful coarse
event of Lemma~\ref{lem:rare-coarse-normals}. For a target $y$ in one
of its bracket neighborhoods, a scale $0<e<d_c$ and a conditional failure
allowance $0<\varepsilon<1/4$, pooled forward commands at at most two
shifted centers return the correct membership label whenever $y$ is
outside the expanded component or is inside it at distance at least
$e$ from its boundary. The label in the intervening inner layer is
unrestricted. The query uses at most
\begin{equation}\label{eq:rare-query-cost}
                     C e^{-(\gamma+3/2)}\log(C/\varepsilon)
\end{equation}
attempted collision bits. An exterior target is labelled zero with
probability one. No launch position, free-start indicator or compass
direction is observed.
\end{proposition}
\begin{proof}
For every $v\in V(\widehat n)$ make fresh pooled forward preparations
at $x_v=y+tv$. Mark this candidate positive if its batch contains at
least one collision. Return one precisely when every candidate is
marked positive. The random direction in each preparation still has
the prescribed uniform law on $tE$.

First exclude all other obstacles from these preparations. At the
projection in \eqref{eq:rare-signed-depth}, support addition writes
$y_0=c_0-z_0$, where $c_0\in C$ and $z_0\in K$ are the support
points with normals $n$ and $-n$. Every point of every attempted
segment has the form
\[
 c_0+(z-z_0)+tv+utw-dn,
       \qquad z\in K,\quad w\in E,\quad 0\le u\le1.
\]
Its distance from $c_0$ is at most $D_*+2t+d_c<d_0$. The physical
separation therefore excludes every component other than $C$, for
all launch offsets and all compass directions.

Suppose that $y$ is exterior, so $d<0$, and choose a true maximizing
vertex $v_*\in V(\widehat n)$. For every $z\in K$, $w\in E$ and
$0\le u\le1$, support addition gives
\begin{align*}
 n\cdot(y+tv_*+z+utw)
 &\ge p_C^{\rm lab}(n)-d+t(n\cdot v_*+u n\cdot w)\\
 &>p_C^{\rm lab}(n).
\end{align*}
The second line uses symmetry of $E$ and maximality of $v_*$, which
imply $n\cdot v_*+u n\cdot w\ge0$. Thus every segment for this
candidate misses $C$, and the preceding separation excludes all other
collisions. Its collision probability is zero. Its mark, and therefore
the conjunction, is deterministically zero. Possible boundary tangencies
are irrelevant to the asserted exterior and positive-depth labels.

Now suppose that $y$ has inner depth $d\ge e$. By
Lemma~\ref{lem:cap-mass},
\[
        \Prob\{y+Z\in C\}=v_j(y)
                    \ge c e^{\gamma+3/2}=:b_e.
\]
Decrease the fixed constant $c$ if necessary so that $b_e\le1$.
For every candidate $v$ and every $z\in K$,
\[
 n\cdot(y+tv+z)
       \ge p_C^{\rm lab}(n)-d+t n\cdot v
       >p_C^{\rm lab}(n),
\]
by \eqref{eq:rare-collar} and \eqref{eq:rare-positive-normal}.
All starts are consequently free. On the independent events
$y+Z\in C$ and $a=-tv$, the attempted segment ends in $C$ and
collides. The pooled success probability for each candidate is at
least $b_e/4$; the selected direction need never be revealed.

Use
\[
            M_e=\left\lceil\frac4{b_e}
                                     \log\frac2\varepsilon\right\rceil
\]
fresh preparations for each candidate. Its probability of having no
collision is at most $\exp(-M_eb_e/4)\le\varepsilon/2$.
A union bound over at most two candidates proves correctness at
positive depth and \eqref{eq:rare-query-cost}. Fresh preparations give
the same bound conditional on every preceding adaptive history.
No upper bound or regularity assumption on the density was used.
\end{proof}

\subsection{The improved stationary reconstruction rate}
\begin{theorem}[Stationary reconstruction by rare collisions]
\label{thm:rare-stationary}
Fix the priors of Theorem~\ref{thm:stationary-jitter} and choose a
fixed compass length satisfying $2t+\diam K<d_0$. There is a finite
adaptive controller, independent of the stationary density, whose
geometric loss is at most $C\nu$ and whose primitive discrete data
are correct with probability at least $1-\delta$, uniformly over the
same physical and density classes. For sufficiently small $\nu>0$
and $0<\delta<1/4$, its costs satisfy
\begin{align}
 J_\nu&\le C\nu^{-1/(s-2)}\log(C/\nu),
                                      \label{eq:rare-centers}\\
 N_\nu&\le C\nu^{-((\gamma+3/2)s+1)/(s-2)}
              \log(C/\nu)\log(C/(\nu\delta)).
                                      \label{eq:rare-attempts}
\end{align}
Every preparation is counted. The launch footprint and compass length
stay fixed, and all nominal centers and physical launches lie in a
prior-dependent fixed aperture. Dyadic nominal positioning to
$O(\nu^{s/(s-2)})$ suffices.

For the two known scales and unknown footprint of
Theorem~\ref{thm:self-finite}, the same conclusions and costs hold
under $2t+2\lambda_2R_K<d_0$, including laboratory $C^2$ footprint
error at most $C\nu$. No footprint-function or density evaluations
are then required. The exact scales, common dilation origin and fixed
geometric and boundary-mass priors retain their stated roles.
\end{theorem}
\begin{proof}
Perform the fixed-accuracy construction of
Lemma~\ref{lem:rare-coarse-normals}, reserving failure probability
$\delta/2$. Put $h=2\pi/m\asymp\nu^{1/(s-2)}$ and $e=c h^s$.
At each of the $m$ equally spaced radial angles, use
Proposition~\ref{prop:rare-query} for the midpoint tests of the
corresponding coarse bracket. Its rational candidate vector is fixed
for that radial direction and remains valid throughout the bracket.
The relaxed bisection invariant of Lemma~\ref{lem:bisection} applies
to the inner indeterminate layer as well. The inball converts a normal
depth of order $e$ to a radial uncertainty of order $e$. Thus
$O(\log(C/e))$ tests per direction give radial values with error
$O(e)$, without a monotonicity assumption on erroneous layer labels.

Lemma~\ref{lem:radial} gives expanded bodies $\widehat P_C$ with
support error $O(h^s)$ in supremum norm and $O(h^{s-2})$ in $C^2$.
For a known footprint, subtract $p_{-K}^{\rm lab}$ and apply the
curvature, final smoothing and period-locking arguments of
Theorem~\ref{thm:stationary-jitter}. These deterministic steps give
the stated physical output and retain the primitive conclusions.

There are at most $Q_h=C(1+h^{-1}\log(C/h))$ fine queries. Give
each one conditional failure allowance $\delta/(2Q_h)$. On the
successful coarse event, Proposition~\ref{prop:rare-query} and a
union bound give simultaneous valid labels with probability at least
$1-\delta/2$. The attempts are bounded by
\[
 C\log(C/\delta)
   +C h^{-1}e^{-(\gamma+3/2)}
                   \log(C/h)\log(C/(h\delta)).
\]
Substituting $e=c h^s$ proves \eqref{eq:rare-attempts}. At most two
shifted command centers replace each fine target; coarse reciprocal
stencils have fixed size. This proves \eqref{eq:rare-centers}.
An enlargement of the earlier fixed aperture by the fixed shifts
protects all these commands and their launch supports.

For numerical commands, round each target once to a dyadic grid of
mesh $O(e)$ and then add the exact dyadic vector $tv$. For sufficiently
small $\nu$, every rounded target lies in the fixed tubular neighborhood
of its radial bracket. The same candidate vector therefore satisfies
\eqref{eq:rare-normal-error}, even though rounding may move the target
off its radial line. Rounding changes distances to the expanded component
by $O(e)$ and hence only enlarges the indeterminate layer by that
amount. It does not turn a claimed exterior test at the rounded point
into a probability estimate. Conservative dyadic lower bounds for
$b_e$ and integer batch counts give the same costs. Their descriptions,
and those of the centers, have
$O(\log(C/\nu)+\log(C/\delta))$ digits per batch.

In the two-scale experiment apply this construction separately to
$P_{i,C}=C+(-\lambda_iK)$, using
\eqref{eq:scaled-density-bound} and the common diameter envelope
$2\lambda_2R_K$. The stronger separation assumption gives
\eqref{eq:rare-margin} at both scales. Match complete components by
Lemma~\ref{lem:scale-matching}, and use the averaged linear support
inverse and convexity verification of Theorem~\ref{thm:self-finite}.
Both scales have the same exponents and uniform prior-controlled
constants. Splitting the confidence budget and adding their costs
therefore preserves \eqref{eq:rare-centers}--\eqref{eq:rare-attempts}.
\end{proof}

For $\gamma=0$, the new upper power is $(3s/2+1)/(s-2)$.
The lower bound of Theorem~\ref{thm:stationary-lower}, specialized to
a fixed sufficiently small uniform-disk footprint and the present
compass choice, consequently gives, on a fixed physical class containing
its packing,
\begin{equation}\label{eq:rare-stationary-bracket}
 c\nu^{-(s+1)/(s-2)}
       \le N^*_{\rm stat}(\nu)
       \le C\nu^{-(3s/2+1)/(s-2)}\log^2(C/\nu).
\end{equation}
The difference between the displayed exponents is now
$s/(2(s-2))$. These bounds do not assert an exact minimax rate.
The improvement comes from a rare positive event whose absence has
an exact geometric explanation on the exterior side; it does not
follow from increasing the precision of reciprocal mean cancellation.

\section{Recovery of an uncalibrated homothety ratio}
\label{sec:unknown-scale}
Integrating a pooled forward collision mean eliminates the unknown
launch density and measures a sum of widths of the original obstacle.
Together with the two positivity supports, this scalar identifies the
ratio of two homothetic settings. An independent normalization uses
the area of a recovered occupation component and requires a weaker
separation condition.

Write $A$ for the physical footprint at the first setting, and let
$Z$ have an unknown probability density $j$ on $A$. The two settings
use the displacements $Z$ and $rZ$, where
\begin{equation}\label{eq:unk-scale-prior}
             1+g_*\le r\le R_*,\qquad g_*>0.
\end{equation}
Only these bounds, not the value of $r$, are supplied. The common
homothety origin is fixed in the nominal laboratory coordinates.
The footprint satisfies \eqref{eq:unknown-K-prior}, with $A$ in place
of $K$, and the density satisfies \eqref{eq:unknown-K-density}.
Assume $t+2R_*R_K<d_0$. Thus the two expanded-component collections
have a known positive separation greater than $t$, and their
cross-setting nesting margin is at least $g_*r_K$.

If the apparatus is originally described by unknown factors
$\lambda_1,\lambda_2$ acting on a parameter footprint $K$, then
$A=\lambda_1K$ and $r=\lambda_2/\lambda_1$. This is an operational
normalization: multiplying $K$ by a positive constant and dividing
both factors by that constant leaves the physical experiment unchanged.
The quantities recovered below are the physical first footprint and
the ratio, without assigning an additional unit to $K$.

All support functions in this section use laboratory coordinates.
For a unit vector $u$, $p_H(u)$ denotes $\sup_{x\in H}x\cdot u$,
the same support function expressed on the unit circle.

\subsection{An integrated collision width}
For a convex body $H$ write
\[
 \mathcal W(H)=p_H(e_1)+p_H(-e_1)+p_H(e_2)+p_H(-e_2),
\]
the sum of its two coordinate widths. This functional is invariant
under translation and additive under Minkowski addition. Let $A$ be
the first launch footprint, let $j$ be its probability density, and
put $P_{1,C}=C+(-A)$. Denote the first-setting pooled forward mean by
$F_1$.

\begin{lemma}[An isolated pooled flux measures width]
\label{lem:unk-scale-flux}
Suppose $2t+\diam A<d_0$. After sufficiently accurate fixed coarse
acquisition of a complete $P_{1,C}$, one can construct a bounded
domain $E$, consisting of finitely many squares of a fixed dyadic
mesh, which contains $P_{1,C}+t\overline B$ and no forward-response
support belonging to another component. Then
\begin{equation}\label{eq:unk-scale-flux}
 \Phi_C:=\int_E F_1(x)\dd x
                      =\frac t2\mathcal W(C).
\end{equation}
The domain and its area are computed from the coarse observations and
the fixed priors. Neither $j$ nor an occupation integral is an input.
\end{lemma}
\begin{proof}
For a displacement $a=tv$, $|v|=1$, set
$S_a(C)=C+[-a,0]$. A line parallel to $v$ meets $C$ in an interval,
a point, or the empty set. On every nonempty interior slice,
$S_a(C)\setminus C$ is an interval of length $t$. The transverse
projection of $C$ has length
$p_C(v^\perp)+p_C(-v^\perp)$. Cavalieri's formula therefore gives
\begin{equation}\label{eq:unk-scale-strip-area}
 \int_{\mathbb R^2}
       (\one_{S_a(C)}-\one_C)
   =t\{p_C(v^\perp)+p_C(-v^\perp)\}.
\end{equation}
Endpoint and tangency conventions affect only null sets. This is also
the segment case of the classical mixed-area formula \cite{Schneider}.

Since $t<d_0$, one segment of length $t$ cannot meet two distinct
physical components, and a start in one component cannot reach another.
Consequently the collision indicator decomposes pointwise as
\[
 B_a^{\mathcal O}
       =\sum_C(\one_{S_a(C)}-\one_C).
\]
The sum is locally finite and its summands are nonnegative. The
contribution $F_{1,C}$ obtained by averaging a single summand against
$j$ and the compass law is supported in
\[
 \bigcup_{v\in\{\pm e_1,\pm e_2\}}
       \{P_{1,C}+[-tv,0]\}
                    \subset P_{1,C}+t\overline B.
\]
Distinct expanded components are separated by at least
$d_0-\diam A$: offsets from the same footprint differ by at most its
diameter. Their displayed enlarged envelopes are therefore separated
by at least $d_0-\diam A-2t>0$. A sufficiently accurate coarse hull,
a fixed enlargement, and a sufficiently small fixed dyadic square mesh
give a domain $E$ containing the complete envelope of $C$ while
remaining disjoint from every other envelope. All error reserves use
fixed fractions of this positive separation. Thus $F_1=F_{1,C}$ on
$E$, and $F_{1,C}=0$ outside $E$.

Tonelli's theorem, translation invariance of area, and $\int_Aj=1$
now show that $\int_EF_1$ is the average of the four strip areas in
\eqref{eq:unk-scale-strip-area}. Each coordinate width occurs twice,
which proves \eqref{eq:unk-scale-flux}. No density differentiability,
upper bound, or symmetry is used.
\end{proof}

\begin{theorem}[Calibration by an integrated collision mean]
\label{thm:unk-scale-flux}
Suppose $2t+2R_*R_K<d_0$. The exact observed functions
$(g_1,g_2,F_1)$ determine the unknown ratio, the first physical
footprint in laboratory coordinates, and every protected original
component. In particular the two observed reciprocal mean pairs
suffice. Put $Q=-A$. For a matched pair $P_1=C+Q$, $P_2=C+rQ$, put
$p_D=p_{P_2}-p_{P_1}$ and use the flux $\Phi_C$ of
Lemma~\ref{lem:unk-scale-flux}. Then
\begin{equation}\label{eq:unk-scale-flux-inverse}
 a=\frac{\mathcal W(P_1)-2\Phi_C/t}{\mathcal W(D)},
 \qquad r=1+a^{-1},\qquad
 p_Q=a p_D,\qquad p_C=p_{P_1}-a p_D.
\end{equation}
The inverse is uniformly stable: expanded-support errors $e$ in
supremum norm and $u$ in $C^2$, together with flux error $m$, give
ratio and Hausdorff errors $O(e+m)$ and laboratory $C^2$ errors
$O(u+m)$. These assertions hold for sufficiently small $u+m$,
with constants depending only on the fixed priors.
\end{theorem}
\begin{proof}
The stopped inverse recovers the complete positivity supports from
$g_1,g_2$, and their fixed nesting margin gives their pairing.
Since $D=(r-1)Q$, one has $P_1=C+aD$ with $a=(r-1)^{-1}$.
Width additivity and \eqref{eq:unk-scale-flux} yield the first formula
in \eqref{eq:unk-scale-flux-inverse}; the remaining formulas follow
by support addition. The denominator is bounded below by $4g_*r_K$.
Also $a$ belongs to the fixed interval
$[1/(R_*-1),1/g_*]$. Consequently errors $e$ in the support values
and $m$ in the flux change $a$ and $r$ by at most $C(e+m)$.
The reconstructed supports then have supremum error $O(e+m)$ and
$C^2$ error $O(u+m)$. Their uniform positive curvature margins
persist at sufficiently small errors, so they represent physical
strictly convex bodies. Reflection recovers $A$ from $Q=-A$.
No first harmonic is discarded, and the original components retain
their laboratory locations.
\end{proof}

\begin{proposition}[A finite integrated collision measurement]
\label{prop:unk-scale-flux-measure}
Assume $2t+2R_*R_K<d_0$. After a prior-dependent coarse acquisition
of one complete first-setting component $P_1$, its integrated
forward mean $\Phi_C$ can be estimated to error $m$ with conditional
confidence $1-\varepsilon$ using
\begin{equation}\label{eq:unk-scale-flux-cost}
                       C m^{-2}\log(C/\varepsilon)
\end{equation}
attempted pooled collision bits. The same order bounds the number
of nominal centers. A dyadic nominal mesh of width $c m$ suffices,
uniformly over the possibly unbounded launch density.
\end{proposition}
\begin{proof}
Let $d_*>2t$ be the lower separation of the first-setting expanded
components. Choose a fixed $b<(d_*-2t)/8$. A sufficiently accurate
coarse hull determines a domain $E$, which is a finite union of
fixed dyadic squares, such that
\[
 P_1+t\overline B\subset E
       \subset P_1+(t+3b)\overline B.
\]
Choose all squares meeting an enlargement of the coarse hull;
reducing its fixed error and the fixed square size gives these
inclusions with positive margins. No other component's forward
response support meets $E$, since each such support is contained
in $P'_1+t\overline B$ and $d_*-2t-3b>0$.
Consequently $F_1=F_{1,C}$ on $E$ and
$\int_E F_1=\Phi_C$.

For $X$ uniform on $E$, one forward bit $Y$ at $X$ gives the bounded
variable $|E|Y$, whose expectation is $\Phi_C$. Independent fresh
repetitions and Hoeffding's inequality give the asserted count.
This uses the pooled forward bit itself and does not require a
direction-resolved mean.

To give the commands finite descriptions, subdivide the fixed
squares by a dyadic mesh $\zeta$ and sample their fine midpoints
uniformly. The expectation changes by $O(\zeta)$ without any
pointwise regularity assumption on $j$. Indeed, condition first
on the launch displacement and a compass direction. The hit
indicator is the difference between a convex segment sweep and
the original convex body. Intersect both with the fixed squares
of $E$. Their total boundary length and their number are uniformly
bounded. The midpoint sum differs from the integral only on fine
squares meeting these boundaries, whose total area is
$O(\zeta)$ by the convex boundary-tube estimate. The bound is
uniform in the launch displacement; averaging over its probability
law and over the four directions preserves it. Take
$\zeta\le c m$ and reserve half the error for this bias. The
confidence assertion is conditional on the coarse acquisition
and uses only subsequent fresh preparations.
\end{proof}

\subsection{Uniform finite reconstruction}
\begin{theorem}[Finite recovery without scale values]
\label{thm:unk-scale-finite}
Fix the numerical bounds of $\mathcal B_\eta$, the footprint and
boundary-mass bounds and \eqref{eq:unk-scale-prior}. Choose the fixed
dyadic compass step to satisfy $2t+2R_*R_K<d_0$. A finite adaptive
controller, supplied with the
two setting labels but neither their ratio nor a footprint or density
oracle, recovers the ratio, the physical first footprint, and a
smooth strictly convex periodic table. For all sufficiently small
$\nu>0$ and $0<\delta<1/4$, its ratio error, footprint laboratory
$C^2$ support error, and geometric loss are at most $C\nu$, and its
primitive discrete data are correct, with probability at least
$1-\delta$ uniformly over the admissible experiments.

Put
\[
                    q_\gamma=
                    \frac{(\gamma+3/2)s+1}{s-2}.
\]
The attempted-bit and nominal-center counts satisfy
\begin{align}
 N_\nu&\le
 C\nu^{-q_\gamma}\log(C/\nu)\log(C/(\nu\delta))
       +C\nu^{-2}\log(C/\delta),\label{eq:unk-scale-attempts}\\
 J_\nu&\le
 C\nu^{-1/(s-2)}\log(C/\nu)
       +C\nu^{-2}\log(C/\delta).\label{eq:unk-scale-centers}
\end{align}
Since $6<s\le7$ and $\gamma\ge0$, one has $q_\gamma>2$, so the
second term of \eqref{eq:unk-scale-attempts} is absorbed by the first.
Both physical spreads stay fixed as $\nu$ tends to zero. The finest
nominal mesh has width $c\nu^{s/(s-2)}$. The list of nominal centers,
setting labels and repetition counts has binary description length
at most
\begin{equation}\label{eq:unk-scale-control-description}
 C\left(\nu^{-1/(s-2)}\log(C/\nu)
              +\nu^{-2}\log(C/\delta)\right)
                         \log(C/(\nu\delta)).
\end{equation}
This is a bound on the digital control description. The fixed
homothety mechanism, its common origin, the known prior bounds,
and the physical realization of the nominal mesh remain apparatus
assumptions.
\end{theorem}
\begin{proof}
Run the rare-event boundary construction of
Proposition~\ref{prop:rare-query} and
Theorem~\ref{thm:rare-stationary} at both settings, stopping before
footprint subtraction. No step uses the scale value: the known
bounds on $r$ give the requisite component, curvature, separation,
and boundary-mass constants. With
$h\asymp\nu^{1/(s-2)}$ and $e=c h^s$, the two collections have
support errors $O(e)$ in supremum norm and $O(\nu)$ in $C^2$.
The fixed nesting margin pairs all complete components needed in
the protected patch. Split the confidence allowance between these
constructions and an independent integrated collision measurement.

Use one complete first-setting component and
Proposition~\ref{prop:unk-scale-flux-measure} with $m=c\nu$ to
estimate $\Phi_C$. Evaluate the four support values giving each
width sum, and apply \eqref{eq:unk-scale-flux-inverse}. The fixed
positive denominator and the stability statement of
Theorem~\ref{thm:unk-scale-flux} give $O(\nu)$ error in $r$ and
in the physical footprint. Subtract this one footprint from
every first-setting expanded support. The resulting supports have
$C^2$ and Hausdorff error $O(\nu)$, and retain strictly positive
curvature. The finite smoothing and period-locking constructions
apply after reducing $\nu_0$. Thus the output is a physical periodic
table with the required primitive data, real period coordinates,
and free area.

The two geometric reconstructions use
$C h^{-1}\log(C/h)$ local queries, each costing at most
$C e^{-(\gamma+3/2)}\log(C/(h\delta))$ attempts. The integrated
measurement adds $C\nu^{-2}\log(C/\delta)$ attempts and at most
that many centers. This proves both counts. Fresh preparations
give the required conditional estimates and a union bound gives
the stated confidence.

Choose the common fine dyadic mesh comparable to $e$ and dividing
the fixed coarse mesh in the integrated collision construction. Since $e\le c\nu$,
it also meets that construction's bias reserve. Every nominal point
lies in a fixed aperture and therefore needs $O(\log(C/\nu))$
digits. Every repetition count is polynomial in $1/\nu$ times a
logarithmic confidence factor and needs at most
$O(\log(C/(\nu\delta)))$ digits. Encoding the finite command list
therefore gives \eqref{eq:unk-scale-control-description}.
\end{proof}

\subsection{An independent area normalization}
For a component $C$, put $Q=-A$ and
\[
 P_1=C+Q,\qquad P_2=C+rQ,
 \qquad v_{1,C}(x)=\int_Aj(z)\one_C(x+z)\dd z.
\]
The subscript $C$ denotes a single complete positive component. The
stopped inverse recovers this restriction of the occupation without
knowing $j$. Fubini's theorem gives the additional invariant
\begin{equation}\label{eq:unk-scale-mass}
 M_C:=\int_{P_1}v_{1,C}(x)\dd x
     =\int_Aj(z)|C-z|\dd z=|C|.
\end{equation}
In particular, the area on the right is measured by the collision
experiment; it is not an input concerning the obstacle.

For planar convex bodies $E,F$, let $V(E,F)$ be their mixed area,
with $V(E,E)=|E|$. In support coordinates,
\begin{equation}\label{eq:unk-scale-mixed-area}
 |E|=\frac12\int_0^{2\pi}(p_E^2-(p_E')^2)\dd\theta,
 \qquad
 V(E,F)=\frac12\int_0^{2\pi}(p_Ep_F-p_E'p_F')\dd\theta.
\end{equation}
These expressions are unchanged by translation. Support addition and
the quadratic area formula are the usual mixed-area identities
\cite{Schneider}.

\begin{theorem}[An unknown scale ratio is identifiable]
\label{thm:unk-scale-exact}
Under the initial condition $t+2R_*R_K<d_0$, the exact reciprocal
differences $(g_1,g_2)$ alone determine $r$, the physical footprint
$A$ in laboratory coordinates,
and every protected original component. The full period group is then
determined by the retained exact recognition theorem.

More explicitly, match a complete pair $P_1,P_2$ by nesting and let
$D$ be the convex body with support
\begin{equation}\label{eq:unk-scale-difference}
                         p_D=p_{P_2}-p_{P_1}.
\end{equation}
Set
\[
 a_0=|P_1|,\qquad b_0^*=V(P_1,D),\qquad c_0=|D|,
 \qquad \mathcal D=(b_0^*)^2-c_0(a_0-M_C).
\]
Then $c_0>0$, $\mathcal D>0$, and
\begin{align}
 a&=\frac{b_0^*-\sqrt{\mathcal D}}{c_0}
    =\frac{a_0-M_C}{b_0^*+\sqrt{\mathcal D}},
 &r&=1+a^{-1},\label{eq:unk-scale-root}\\
 p_Q&=a p_D,
 &p_C&=p_{P_1}-a p_D.\label{eq:unk-scale-inverse}
\end{align}
All support functions in these formulas are in the laboratory frame.
No centering convention for $A$ or for its density is imposed.
\end{theorem}
\begin{proof}
Lemma~\ref{lem:scale-matching} uses only the lower nesting margin and
the separation, so it supplies the pairing without the value of $r$.
For the true pair, $D=(r-1)Q$ and $a=(r-1)^{-1}$. In particular $D$
is a full-dimensional strictly convex body containing a fixed inball.
The equation $P_1=C+aD$ and the mixed-area identity give
\[
           M_C=a_0-2a b_0^*+a^2c_0.
\]
Moreover,
\begin{equation}\label{eq:unk-scale-root-margin}
 b_0^*-ac_0=V(C,D)>0,
 \qquad \mathcal D=V(C,D)^2.
\end{equation}
Thus the true $a$ is precisely the smaller root in
\eqref{eq:unk-scale-root}. Any other positive value $a'$ for which
$p_{P_1}-a'p_D$ is the support of a convex body of area $M_C$ would
also satisfy $b_0^*-a'c_0>0$, and hence would be the same smaller
root. This proves uniqueness among physical candidates.
Formula~\eqref{eq:unk-scale-inverse} now recovers $C$ and $Q$;
reflection recovers $A$. The answer for $r$ and $A$ is common to all
matched pairs. The recognition theorem applies to the recovered
original components.
\end{proof}

\begin{proposition}[Stability of the area-normalized inverse]
\label{prop:unk-scale-stability}
Under the fixed priors above, suppose the matched expanded supports
at both settings are known with supremum error $e$ and $C^2$ error $u$, and
$M_C$ is known with error $m$. For sufficiently small $u+m$, the
smaller-root construction is well defined and gives
\begin{align}
 |\widehat r-r|+d_H(\widehat A,A)
                 +\max_Cd_H(\widehat C,C)&\le C(e+m),
                                      \label{eq:unk-scale-stability-0}\\
 \|p_{\widehat A}-p_A\|_{C^2}
       +\max_C\|p_{\widehat C}-p_C\|_{C^2}&\le C(u+m).
                                      \label{eq:unk-scale-stability-2}
\end{align}
The constants depend on the stated bounds, including $g_*$, but not
on the unknown ratio, footprint, or density.
\end{proposition}
\begin{proof}
The true $D$ has uniformly positive curvature radius. Hence the
estimated support difference remains a strictly convex support for
small $u$. On bounded convex classes area and mixed area are Lipschitz
in Hausdorff distance, so their estimated coefficients have error
$O(e)$. For example, area follows from the two parallel-body
inclusions, and mixed area follows by polarization using $|E+F|$.

The inball of $D$ and the lower perimeter bound for $C$ imply
$V(C,D)\ge c>0$. Equation~\eqref{eq:unk-scale-root-margin} therefore
keeps the discriminant uniformly away from zero. The two equivalent
root formulas have bounded derivatives throughout a fixed
neighborhood of the admissible data. Also
$1/(R_*-1)\le a\le1/g_*$. Consequently
$|\widehat a-a|+|\widehat r-r|\le C(e+m)$.
Multiplication by $p_D$, whose $C^2$ norm is uniformly bounded in the
protected frame, proves both estimates. The positive curvature
margins of the original bodies and $A$ persist, so the recovered
supports represent actual convex bodies. A single pair may be used
for $a$ and $A$, after which the same footprint is subtracted from all
other first-setting components.
\end{proof}

\subsection{A finite measurement of one component area}
The area normalization can be measured without a fine deterministic
grid of occupation queries. The following construction uses a fixed
coarse domain and a bounded weight computable before the area
measurements begin.

\begin{proposition}[Integrated occupation from pooled bits]
\label{prop:unk-scale-mass}
After a prior-dependent coarse acquisition of one complete expanded
component $P_1$, its mass $M_C$ can be estimated to error $m$ with
conditional confidence $1-\varepsilon$ using
\begin{equation}\label{eq:unk-scale-mass-cost}
                       C m^{-2}\log(C/\varepsilon)
\end{equation}
attempted pooled collision bits at the first setting. The same order
bounds the number of nominal centers. A dyadic nominal mesh of width
$c m$ suffices. All centers lie in one fixed protected aperture, and
the assertion is uniform over the possibly unbounded density $j$.
\end{proposition}
\begin{proof}
Let $d_*>t$ be the known lower separation of the expanded components.
Choose a fixed positive $b<(d_*-t)/8$. From a sufficiently accurate
coarse hull construct a domain $E$, which is a finite union of squares
of a fixed dyadic mesh $d$, such that
\begin{equation}\label{eq:unk-scale-domain}
 P_1\subset E,\qquad
 \dist(P_1,E^c)>b,\qquad E\subset P_1+3b\overline B.
\end{equation}
The mesh may be chosen to divide the dyadic compass step $t$.
Taking all mesh squares meeting a fixed small enlargement of the
coarse hull gives these inclusions, after reducing its fixed error
and $d$. In particular no other positive component meets
$E+t\overline B$. Boundaries of squares may be assigned consistently
to one neighboring square; they have measure zero.

On $E$ the first-setting forcing therefore satisfies
\[
             g_1=(T^E-I)v_{1,C}.
\]
Let $w_E=(I-T^E)^{-1}\one_E$. The killed-walk exit bound gives
$0\le w_E\le H_E$, where $H_E$ is a fixed prior constant. Because
$d$ divides $t$, the state graph of the killed walk is identical for
all starting points in the interior of any one mesh square. Thus
$w_E$ is constant on each of the finitely many squares. Its values
are the solution of a finite rational linear system with coefficients
$0,1,1/4$, and are available without any occupation observations.
The inverse exists since the walk exits the bounded domain almost
surely and has uniformly bounded mean exit time.

The compass kernel is symmetric. Integrating the killed operator on
$E$ and interchanging its finitely many translations gives
\begin{equation}\label{eq:unk-scale-duality}
 \int_E w_E g_1
   =\int_E v_{1,C}(T^E-I)w_E
   =-\int_Ev_{1,C}=-M_C.
\end{equation}
If $X$ is uniform on $E$, take one fresh forward bit $Y^+$ and one
fresh reciprocal reverse bit $Y^-$ at $X$. The variable
\[
            Z_E=-|E|w_E(X)(Y^+-Y^-)
\]
has expectation $M_C$ and absolute value at most $|E|H_E$. Independent
repetitions and Hoeffding's inequality prove
\eqref{eq:unk-scale-mass-cost} for continuous nominal sampling.
This calculation does not assume that an occupation bit is observed.

For a finite command description choose a dyadic mesh $\zeta$ dividing
$d$ and sample $X_\zeta$ uniformly from the fine square midpoints in
$E$. This changes the expectation by $O(\zeta)$, uniformly over $j$.
Here is a direct bound which does not require an upper bound or a
modulus of continuity for the density. Before averaging in the
launch displacement, each forward or reverse bit is a finite
difference of indicators of convex obstacles and their segment
sweeps, translated by the displacement. Intersect these sets with
the fixed squares of $E$, on which $w_E$ is constant. Their total
boundary length and number are uniformly bounded. A midpoint sum
and its integral agree on every fine square which does not meet
one of these boundaries. The total area of the remaining squares is
$O(\zeta)$, by the boundary-tube estimate for convex sets. The bound
is uniform in the launch displacement, so integration against the
probability density preserves it. The same argument applies to the
four pooled directions and the reciprocal shifts. Taking
$\zeta\le c m$ reserves at most half the permitted error for this
quadrature bias. Hoeffding's inequality supplies the other half.
The procedure and its confidence bound hold conditional on the
successful coarse acquisition, using fresh independent preparations.
\end{proof}

The normalization argument uses the first setting's density only
through its integral and the quantitative boundary lower bound.
It therefore remains valid if the two settings have different
stationary densities on $A$ and $rA$, provided their normalized
boundary-mass bounds are uniform. Their supports must still be
exactly homothetic. Finally, an unknown common unscaled displacement
$b$ in the commands replaces $C$ by $C-b$ throughout, leaving $r$
and the physical footprint unchanged. Translating the table and
$b$ by the same vector leaves every collision bit unchanged. The
absolute location in that extended experiment requires a laboratory
reference.

\section{Collision rigidity with unregistered launch laws}
\label{sec:unregistered-laws}
The common origin of the homothetic settings can also be removed.
An integrated collision mean first measures the width of the launch
footprint itself. A translation-covariant functional of the complete
collection of centered component supports then separates that footprint
from the obstacles. Neither step requires corresponding obstacles to
be identified across settings. The resulting exact inverse applies to
arbitrary separated convex configurations, before periodicity or a
finite list of component species is assumed.

Let $\mathcal O$ be a nonempty locally finite union of compact convex
bodies with nonempty interiors, diameter at most $D_0$, and mutual
distance at least $d_0$. At setting $i\in\{1,\ldots,m\}$, an unknown
stationary probability density $j_i$ is positive almost everywhere in
the interior of its convex support $A_i$. Suppose the supports are
translated positive homothetic copies of one another. Normalize them as
\begin{equation}\label{eq:registration-model}
 A_i=b_i+\rho_i A_0,\qquad s(A_0)=0,\qquad
 \rho_1=1,\qquad \rho_i>0,
\end{equation}
where $s$ denotes the Steiner point. Thus $A_0=A_1-s(A_1)$ is the
physical first footprint after centering, and $b_i=s(A_i)$. None of
$A_0,b_i,\rho_i$ is supplied. At least two scale factors are distinct.
The densities at different settings need not be rescalings of one
common density. Assume
\begin{equation}\label{eq:registration-separation}
                    2t+\max_i\diam A_i<d_0.
\end{equation}
The observations are the pooled forward and reciprocal reverse means
$F_i,R_i$ at all nominal centers, and $g_i=F_i-R_i$. No regularity
beyond convexity, and no quantitative boundary-mass lower bound, is
required for the exact results of this section.

Write $p_H$ for a laboratory support function and
$p_H^\circ(u)=p_H(u)-s(H)\cdot u$ for its centered version,
$u\in\mathbb S^1$. The centered support functions of all obstacle
components are uniformly bounded and uniformly Lipschitz. The same
is true of their expanded components at any fixed setting.

\subsection{A footprint width without a matching component}
By Theorem~\ref{thm:global-occupation}, $g_i$ determines the occupation
and its complete positive-component collection
\[
               \mathcal P_i=\{P_{i,C}=C+(-A_i):C\in\mathcal C\},
\]
where $\mathcal C$ is the collection of original components.
For any $P\in\mathcal P_i$, choose a bounded domain $E_P$ containing
$P+t\overline B$ and meeting no other component's forward-response
support. Such a domain exists by \eqref{eq:registration-separation}.
Let $\mathcal W(H)$ denote the sum of the two coordinate widths and put
\begin{equation}\label{eq:registration-deficit}
 d_i(P)=\mathcal W(P)-\frac2t\int_{E_P}F_i(x)\dd x.
\end{equation}

\begin{lemma}[The footprint deficit]
\label{lem:registration-deficit}
The quantity in \eqref{eq:registration-deficit} is independent of
$P\in\mathcal P_i$ and of the admissible domain $E_P$. Its common
value is
\begin{equation}\label{eq:registration-width-ratio}
              d_i=\mathcal W(A_i)>0,
              \qquad \rho_i=d_i/d_1.
\end{equation}
In particular, the order and relative sizes of the settings are
determined without any cross-setting component correspondence.
\end{lemma}
\begin{proof}
For $P=C+(-A_i)$ the isolated forward contribution has integral
$t\mathcal W(C)/2$. Indeed a displacement $tv$ adds a strip of area
$t\{p_C(v^\perp)+p_C(-v^\perp)\}$ to the segment sweep; integrating
the launch translation eliminates the probability density, and the
four compass directions count each coordinate width twice. This is
the slice calculation of Lemma~\ref{lem:unk-scale-flux}, which uses
only convexity. Width addition now gives
$d_i(P)=\mathcal W(P)-\mathcal W(C)=\mathcal W(A_i)$.
Translation invariance and homogeneity of width prove the ratio
formula. Nonempty interior makes the denominator positive.
\end{proof}

\subsection{An unlabeled support envelope}
Define the lower centered-support envelope at setting $i$ by
\begin{equation}\label{eq:registration-envelope}
       L_i(u)=\inf_{P\in\mathcal P_i}p_P^\circ(u),
                   \qquad u\in\mathbb S^1.
\end{equation}
This is a finite-valued Lipschitz function, although it need not be
the support function of a convex body and the infimum need not be
attained. It is defined directly from the complete collection; no
ordering, multiplicity convention, or choice of corresponding
components is involved.

\begin{theorem}[A geometric inverse without a supplied origin]
\label{thm:registration-rigidity}
Under \eqref{eq:registration-model}--\eqref{eq:registration-separation},
the exact pooled reciprocal mean pairs determine all scale factors
$\rho_i$, the centered first footprint $A_0$, and the configurations
\begin{equation}\label{eq:registration-output}
                       \mathcal O_i=\mathcal O-b_i.
\end{equation}
They therefore determine the obstacle configuration up to a common
translation and determine the relative footprint translations modulo
its full period group.

Explicitly, choose any $k$ with $d_k\ne d_1$, set
$\rho_i=d_i/d_1$, and put
\begin{equation}\label{eq:registration-support-inverse}
 q(u)=\frac{L_k(u)-L_1(u)}{\rho_k-1}.
\end{equation}
Then $q=p_{-A_0}$ is an actual centered support function. For every
$P\in\mathcal P_i$ the function
\begin{equation}\label{eq:registration-body-inverse}
                         p_P-\rho_iq
\end{equation}
is the support of the corresponding component of $\mathcal O_i$.
Writing $\Pi=\operatorname{Per}(\mathcal O)$, the relative registration
is the exactly determined coset
\begin{equation}\label{eq:registration-coset}
 \{a:\mathcal O_i+a=\mathcal O_1\}=(b_i-b_1)+\Pi.
\end{equation}
No periodicity, finite component count, smoothness, or known homothety
origin is an assumption of this theorem.
\end{theorem}
\begin{proof}
The global occupation inverse gives $\mathcal P_i$, and
Lemma~\ref{lem:registration-deficit} gives the ratios. Steiner centering
commutes with Minkowski addition. With $q_0=p_{-A_0}$ one has
\[
 p_{P_{i,C}}^\circ=p_C^\circ+\rho_iq_0,
 \qquad
 L_i(u)=\inf_Cp_C^\circ(u)+\rho_iq_0(u).
\]
The same infimum occurs at every setting, even when different
components minimize it at different directions. Subtraction proves
\eqref{eq:registration-support-inverse}. Its value is the support of
the true footprint; no smoothness or convexity of the separate
envelopes has been asserted. The identity
$P_{i,C}=(C-b_i)+\rho_i(-A_0)$ proves
\eqref{eq:registration-body-inverse} and hence
\eqref{eq:registration-output}. Formula
\eqref{eq:registration-coset} follows by translating the two recovered
configurations. Their common full period group is also determined
directly from the response by Theorem~\ref{thm:global-response-rigidity}.
\end{proof}

\begin{corollary}[The complete geometric ambiguity]
\label{cor:registration-fiber}
Suppose two admissible experiments have the same pooled reciprocal
mean pairs. Their scale ratios and centered footprints agree, and
there are a vector $h$ and periods $\pi_i\in\Pi$ such that
\begin{equation}\label{eq:registration-gauge}
            \mathcal O'=\mathcal O+h,
              \qquad A_i'=A_i+h+\pi_i.
\end{equation}
Conversely every transformation \eqref{eq:registration-gauge} is
realized with identical collision laws by translating the density at
setting $i$ by $h+\pi_i$. Thus the geometric fiber of the observed
response is exactly this orbit. In a finite nonempty configuration
$\Pi=\{0\}$, so all relative footprint translations are recovered
as vectors, and the only geometric ambiguity is a common translation.
\end{corollary}
\begin{proof}
The theorem gives $\mathcal O-b_i=\mathcal O'-b_i'$ at every setting.
Put $h=b_1'-b_1$. Then $\mathcal O'=\mathcal O+h$, and comparison
at setting $i$ gives $b_i'-b_i-h\in\Pi$. The centered footprints
already agree, proving \eqref{eq:registration-gauge}.

For the converse, couple the new launch displacement to the old one
by $Z_i'=Z_i+h+\pi_i$. Translating both the table and a flight by
$h$ preserves its collision bit, and a further period translation
$\pi_i$ preserves it again. This is a pointwise coupling for forward
and reciprocal reverse preparations, and in fact for every permitted
short direction. It proves equality of the full binary laws, not
merely of their positivity sets. A bounded nonempty configuration
has no nonzero translational period.
\end{proof}

\subsection{Finite periodic reconstruction and registration}
The global envelope is appropriate for the exact full-field inverse.
For a periodic table its role can instead be played by a finite
average. This gives the same finite statistical orders as the
registered experiment.

\begin{lemma}[Cancellation of the periodic patch defect]
\label{lem:registration-period-margin}
Let $Q$ be any fixed convex body and $b$ any vector, and assume the
expanded components $C-b+Q$ remain separated. Their full period group
and primitive component-orbit count are those of $\mathcal O$.
The component comparison distance $d_*$ and the orbitwise patch
defect are unchanged by this common expansion and translation.
Consequently the finite nonperiod margin of
\eqref{eq:patch-margin} applies to period and registration tests on
these expanded components.
\end{lemma}
\begin{proof}
A period permutes complete connected components. Cancelling the same
support function $p_Q$ from the matched pair proves the period
assertion and gives the orbit correspondence. Steiner additivity
also gives, for every pair and every vector $w$,
\[
 d_*(C-b+Q+w,D-b+Q)=d_*(C+w,D).
\]
Take the maximum over one representative of each component orbit and
the infimum over all physical components, in both translation signs.
This defect is independent of the representatives, invariant under
adding a true period to $w$, and unchanged by the displayed expansion.

For completeness, the bounded-patch margin in
\eqref{eq:patch-margin} controls every nonperiod center-difference
class. Given any two components, move their centers separately into
$Q_B$ by presentation-lattice vectors. Their difference changes by
a true period, and its new endpoints lie in the cutoffs of
\eqref{eq:patch-margin}. The defect is unchanged. Moreover $Q_B$
contains representatives of every presentation orbit, so its
maximum is precisely the orbitwise defect. This argument applies
after any common translation, and proves the claimed transfer of
the same positive margin.
\end{proof}

If $n$ is the recovered primitive orbit count, choose one complete
representative per orbit at each setting and define
\begin{equation}\label{eq:registration-orbit-average}
 H_i=\frac1n\sum_{P\bmod\Pi}p_P^\circ.
\end{equation}
The representative choices and their ordering do not matter. Indeed,
\[
 H_i=\frac1n\sum_{C\bmod\Pi}p_C^\circ+\rho_iq,
                 \qquad q=\frac{H_k-H_1}{\rho_k-1}.
\]
Thus finite orbit averages replace the global infima, without
introducing a cross-setting matching problem.

\begin{theorem}[Finite reconstruction of unregistered settings]
\label{thm:registration-finite}
Fix the periodic priors of $\mathcal B_\eta$, a finite number of
settings, uniform $C^{6,\beta}$ centered-footprint and curvature
bounds, and uniform density boundary-mass bounds with exponent
$\gamma\ge0$. Suppose the scales lie in a fixed compact subinterval
of $(0,\infty)$ and at least one satisfies $|\rho_k-1|\ge g_*>0$.
The translations $b_i$ are unknown; a fixed bound on their sizes
suffices to retain a fixed physical launch aperture. Choose a fixed
dyadic compass satisfying the strict separation
\eqref{eq:registration-separation} uniformly over the priors.

A finite controller using only the pooled collision bits recovers
the ratios, centered footprints, a periodic table in the frame
$\mathcal O-b_1$, and representatives of all registration cosets
\eqref{eq:registration-coset}. With probability at least $1-\delta$,
the ratio errors, laboratory $C^2$ errors in this translation frame,
period-basis and free-area errors, and errors of the registration
representatives modulo the true $\Pi$ are at most $C\nu$.
The primitive orbit count and rational period relations are correct.

Let $N_\nu$ count attempted bits, $J_\nu$ count center-labelled
batches, and $S_\nu$ count distinct nominal sites. A later batch at
the same site is counted again in $J_\nu$; each center sampled for
an integrated measurement is one single-attempt batch. Then
\begin{align}
 N_\nu&\le C\nu^{-((\gamma+3/2)s+1)/(s-2)}
           \log(C/(\nu\delta))
           +C\nu^{-2}\log(C/\delta),
                                    \label{eq:registration-attempts}\\
 S_\nu\le J_\nu&\le
       C\nu^{-1/(s-2)}\log(C/\nu)
                  +C\nu^{-2}\log(C/\delta).
                                    \label{eq:registration-batches}
\end{align}
The finest nominal mesh is $O(\nu^{s/(s-2)})$, and the batch list,
including setting labels and repetition counts, has binary
description length at most $C J_\nu\log(C/(\nu\delta))$.
These counts do not assign a manufacturing or metrology cost to the
physical homothety relation or to the realization of nominal controls.
\end{theorem}
\begin{proof}
At each setting run the shrinking-layer rare-collision construction
of Proposition~\ref{prop:shrinking-rare} before footprint subtraction. Its geometric and boundary-mass constants
depend on centered footprint bounds and diameter, not on the
unknown translation. With $h\asymp\nu^{1/(s-2)}$ and $e=c h^s$,
the expanded components are reconstructed with support errors
$O(e)$ in supremum norm and $O(\nu)$ in $C^2$. The fixed physical
aperture includes the unknown translations through their given
coarse bounds. Lemma~\ref{lem:registration-period-margin} permits
period locking at each setting before any calibration. It gives
the common primitive orbit count and the complete representative
lists needed in \eqref{eq:registration-orbit-average}.

Choose any complete component separately at each setting and use
Proposition~\ref{prop:unk-scale-flux-measure} to estimate its
integrated forward mean to error $c\nu$. The proof of that
proposition uses only the isolated expanded component, and hence
applies unchanged to the translated footprint. Form the deficits
\eqref{eq:registration-deficit} and their ratios. Their errors are
$O(\nu)$, since the first footprint width has a fixed positive
lower bound. The separated scale can be selected from these
estimates: a setting maximizing the measured distance from one
has true gap at least $g_*/2$ at sufficiently small error.
The orbit-average formula then estimates $q$ with $C^2$ and
supremum error $O(\nu)$. Subtracting the same $\rho_iq$ from
every component at setting $i$ gives all configurations
$\mathcal O_i$ with the same error. Positive curvature margins
ensure physical output bodies, and the retained finite smoothing
preserves these estimates.

It remains to register the reconstructed frames. Choose a complete
root component of $\mathcal O_i$ and enumerate its center differences
to the finite representative list of $\mathcal O_1$. A true
registration coset has a representative on this list. Test each
candidate by translating all component-orbit representatives and
comparing to the other table modulo its recovered lattice. Only
bounded integer translations can enter the protected test patch,
by the fixed period-basis bounds. An
incorrect candidate differs from a correct one by a nonperiod
center difference in $\mathcal O$. Its defect is at least $\eta$
by Lemma~\ref{lem:registration-period-margin}; the numerical defect
error is $O(\nu)$. Thus the same finite threshold separates the
correct registration coset. Accepted candidates differ by periods
and approximate true coset representatives to $O(\nu)$. No unique
shape, asymmetry, or a distinguished component is required. Use one
reconstructed frame and these registrations to give a common
periodic output. The original period-locking and parallel-area
estimates give the stated discrete, basis, and free-area conclusions.

There are a fixed number of geometric constructions. Each uses
$C h^{-1}\log(C/h)$ center batches, and the shrinking-layer sum of
Proposition~\ref{prop:shrinking-rare} bounds its attempted bits by
$C h^{-1}e^{-(\gamma+3/2)}\log(C/(h\delta))$. A fixed
number of integrated measurements cost
$C\nu^{-2}\log(C/\delta)$ single-attempt batches. Confidence
allocation and fresh preparations give the joint event and both
counts. Orbit averages and registration tests use the finite
reconstructions and require no further collision observations.
The dyadic mesh and the digital description bound are obtained
as in Theorem~\ref{thm:unk-scale-finite}.
\end{proof}

\begin{corollary}[Finite nonperiodic configurations]
\label{cor:registration-finite-cloud}
In place of the periodic prior in
Theorem~\ref{thm:registration-finite}, suppose that $\mathcal O$ has
$1\le n\le n_0$ components, with the fixed diameter, separation,
$C^{6,\beta}$ centered-support and positive curvature bounds.
Assume a known bounded protected aperture contains all expanded
components at every setting and all their forward-response envelopes.
The aperture includes the fixed collars needed for coarse acquisition
and reciprocal inversion. Retain the finite setting count, centered
footprint regularity and curvature bounds, boundary-mass bounds,
compact scale bounds, positive scale gap, coarse bounds on the unknown
translations, and strict compass separation of that theorem.

Without a periodicity or nonperiod-patch-margin assumption, a finite
controller recovers the number $n$, the scale ratios, the centered
footprints, the complete configuration $\mathcal O-b_1$, and the
relative translations $b_i-b_1$ as vectors. With probability at least
$1-\delta$, all scalar errors and the laboratory $C^2$ support errors
in this common translation frame are at most $C\nu$, and $n$ is
correct. The attempted-bit, center-batch and distinct-site bounds
\eqref{eq:registration-attempts}--\eqref{eq:registration-batches},
the nominal precision, and the digital description bound remain
valid, with constants depending also on the complete-aperture and
component-count bounds.
\end{corollary}
\begin{proof}
Coarse acquisition over the known aperture finds every expanded
component. The uniform inballs and separation used in that
construction ensure that each component has retained coarse nodes
and is counted once. The protective collar excludes exterior
fragments, and completeness of the aperture rules out unobserved
components. Thus every setting has the same recovered count $n$,
on the common successful coarse event. This step uses no period
test.

Replace \eqref{eq:registration-orbit-average} by the finite averages
\[
 H_i=\frac1n\sum_{P\in\mathcal P_i}p_P^\circ,
       \qquad z_i=\frac1n\sum_{P\in\mathcal P_i}s(P).
\]
There is one expanded component for each original component, so
\begin{equation}\label{eq:registration-finite-cloud-identities}
 H_i=\frac1n\sum_{C\in\mathcal C}p_C^\circ+\rho_iq,
 \qquad
 z_i=\frac1n\sum_{C\in\mathcal C}s(C)-b_i.
\end{equation}
The deficits again determine the ratios, and any separated scale
gives $q=(H_k-H_1)/(\rho_k-1)$. In addition,
\begin{equation}\label{eq:registration-centroid-registration}
                         b_i-b_1=z_1-z_i.
\end{equation}
These formulas are invariant under independent permutations of the
component lists. Thus neither shape matching nor a discrete
registration decision is required.

Apply the shrinking-layer reconstruction to all components and
the integrated measurement to one arbitrary component per setting,
with the same confidence and accuracy reserves as before. Each
expanded support has $C^2$ error $O(\nu)$ and supremum error
$O(e)$, where $e\asymp\nu^{s/(s-2)}$. Averaging cannot enlarge
these bounds. The scalar deficits have error $O(\nu)$, so the
positive width and scale-gap margins give ratio and footprint
$C^2$ errors $O(\nu)$. Subtraction from the first-setting supports
recovers $\mathcal O-b_1$ with the same error. Steiner points have
error $O(e)$; hence \eqref{eq:registration-centroid-registration}
recovers all relative translations with error $O(e)$, in particular
$O(\nu)$. Finite smoothing preserves these estimates and strict
convexity. The component count is already correct on the coarse
event.

The number of radial searches is bounded by the fixed component
and setting counts times $C\nu^{-1/(s-2)}$. Their shrinking-layer
sum and the fixed number of integrated measurements are exactly
those used to prove \eqref{eq:registration-attempts} and
\eqref{eq:registration-batches}. Complete-component averaging and
the centroid formula require no further observations. The common
dyadic mesh and finite batch encoding therefore give the same
nominal precision and digital description bounds.
\end{proof}

\section{Reciprocal designs without preferred directions}
\label{sec:isotropic-rigidity}
The exact rigidity is not tied to the four compass directions. We
first identify the property of the displacement law used in the
stopping argument, and then give a rotationally invariant collision
experiment with an intrinsic perimeter normalization. The finite
compass minimax experiment remains the one specified in
Corollary~\ref{cor:stationary-minimax}.

\subsection{A general centered displacement law}
Let $\mu$ be a known probability measure on displacements in $\R^2$
satisfying
\begin{equation}\label{eq:general-command-law}
 \int a\,\mu(\dd a)=0,\qquad
 \sigma_\mu^2=\int|a|^2\,\mu(\dd a)>0,\qquad
 \operatorname{supp}\mu\subset\ell\overline B.
\end{equation}
No symmetry, density, or nondegeneracy of the covariance matrix is
assumed. The obstacles and stationary launch density are as in
Section~\ref{sec:global-response}, with
\begin{equation}\label{eq:general-command-separation}
                              \ell+\Delta<d.
\end{equation}
Draw $Z\sim j$ independently of $a\sim\mu$. The forward attempt
is $(x+Z,a)$, and its reciprocal reverse attempt is
$(x+Z+a,-a)$ with fresh draws from that same nominal joint law.
Let $F_\mu,R_\mu$ denote the pooled means, and put
\[
 g_\mu=F_\mu-R_\mu,\qquad
 T_\mu f(x)=\int f(x+a)\,\mu(\dd a).
\]
The displacement realization is not returned by the sensor.

\begin{theorem}[Design-independent translation rigidity]
\label{thm:general-command-rigidity}
Under \eqref{eq:general-command-law}--\eqref{eq:general-command-separation},
the occupation field satisfies $g_\mu=(T_\mu-I)v$ and is uniquely
determined by $g_\mu$. With
\begin{equation}\label{eq:general-exit-budget}
                     H_\mu=\frac{(D+\Delta+\ell)^2}{\sigma_\mu^2},
\end{equation}
one has
\begin{equation}\label{eq:general-stopping-inverse}
 v(x)=\sup_{\E_x\tau\le H_\mu}
           \E_x\left[-\sum_{k<\tau}g_\mu(X_k)\right],
\end{equation}
where $X$ is the virtual walk with increment law $\mu$ and immediate
stopping is allowed. Consequently
\begin{equation}\label{eq:general-response-periods}
 \operatorname{Per}(F_\mu,R_\mu)
 =\operatorname{Per}(g_\mu)
 =\operatorname{Per}(v)
 =\operatorname{Per}(\mathcal O).
\end{equation}
The same prior-uniform Lipschitz inverse bound holds with constant
$H_\mu$. None of these assertions assumes periodicity of the table.
\end{theorem}
\begin{proof}
The endpoint balance identity applies to each displacement of length
at most $\ell<d$. Independence of $Z$ and $a$ therefore gives
$g_\mu=(T_\mu-I)v$. The positive components are still $C+(-A)$,
with diameter at most $D+\Delta$ and gaps greater than $\ell$.
An exit step from one positive component consequently lands at zero
occupation.

The mean-zero assumption and independence of increments make
$|X_n-x|^2-n\sigma_\mu^2$ a martingale. Before component exit,
and on its exit step, the displacement from $x$ is at most
$D+\Delta+\ell$. Optional stopping first at a bounded time and
then monotone convergence give expected exit time at most $H_\mu$.
This proof uses the trace of the covariance, not its invertibility;
it also applies to a law supported on a line. Markov's inequality and
the Markov property give the same geometric block tail as in
Lemma~\ref{lem:global-exit}.

Stopping the bounded occupation martingale gives a payoff at most
$v(x)$ for every integrable rule. The containing-component exit,
or immediate stopping when $v(x)=0$, attains equality and proves
\eqref{eq:general-stopping-inverse}. Comparing its payoffs for two
forcings gives the Lipschitz bound. The stopping functional commutes
with translations, so every period of $g_\mu$ is a period of $v$.
Such a period permutes the separated expanded components; cancelling
$p_{-A}$ from their laboratory support functions gives a period of
$\mathcal O$. The remaining implications in
\eqref{eq:general-response-periods} follow from translation covariance
and subtraction, exactly as in
Theorem~\ref{thm:global-response-rigidity}.
\end{proof}

For a law with finite lattice support, a sufficiently large finite
killing set gives a stopping polytope as in
Lemma~\ref{lem:stencil-polytope}. For a law with continuous support,
\eqref{eq:general-stopping-inverse} is an exact response-field theorem,
not a claim that finitely many mean samples evaluate the integral
operator exactly. In either case the displacement law is a specified
control, whereas the stationary launch density is unknown.

\subsection{Isotropic collision means and a perimeter deficit}
Now fix a length $t>0$ and take the direction uniformly on the unit
circle:
\begin{equation}\label{eq:isotropic-design}
       a=t(\cos\theta,\sin\theta),\qquad
                      \theta\sim\operatorname{Unif}[0,2\pi).
\end{equation}
Its mean is zero and $\sigma_\mu^2=t^2$. The observation operator
commutes with every orthogonal transformation; there is no distinguished
apparatus axis. Nominal centers and the reciprocal joint law are
still prescribed controls.

Use the translated homothetic supports of
\eqref{eq:registration-model}, with at least two distinct scales, and
assume $2t+\max_i\diam A_i<d_0$. Let $\mathcal L(H)$ denote the
Euclidean perimeter of a planar convex body. From the exact reciprocal
difference at setting $i$, recover the complete expanded-component
collection $\mathcal P_i$ by
Theorem~\ref{thm:general-command-rigidity}. For $P=C+(-A_i)$,
choose a bounded region $E_P$ containing $P+t\overline B$ and no
other component's forward-response support, and put
\begin{equation}\label{eq:isotropic-deficit}
 d_i^{\mathrm{iso}}(P)=\mathcal L(P)
                    -\frac{\pi}{t}\int_{E_P}F_{i,\mu}(x)\dd x.
\end{equation}

\begin{lemma}[Isotropic collision perimeter]
\label{lem:isotropic-perimeter}
For every component and setting,
\begin{equation}\label{eq:isotropic-perimeter}
 \int_{E_P}F_{i,\mu}(x)\dd x=\frac{t}{\pi}\mathcal L(C),
 \qquad d_i^{\mathrm{iso}}(P)=\mathcal L(A_i)>0.
\end{equation}
In particular the deficit is independent of the selected component,
the launch density, and the admissible region $E_P$.
\end{lemma}
\begin{proof}
At a fixed direction $e$, horizontal slicing in that direction gives
\[
 |(C+[-te,0])\setminus C|=t\,w_C(e^\perp),
 \qquad w_C(u)=p_C(u)+p_C(-u).
\]
Tonelli's theorem eliminates the independent launch translation from
the integrated mean. Averaging the strip identity in $e$ gives
\[
 \int_{E_P}F_{i,\mu}
        =\frac{t}{2\pi}\int_0^{2\pi}w_C(n(\theta))\dd\theta
        =\frac{t}{\pi}\mathcal L(C).
\]
The last equality is Cauchy's perimeter formula. To fix its
normalization, for a smooth strictly convex body it follows by
integrating $p_C+p_C''$ over the circle and using periodicity of
$p_C'$. General convex bodies follow by smooth approximation of
support functions and continuity of perimeter. The same formula
shows $\mathcal L(C+H)=\mathcal L(C)+\mathcal L(H)$; these are
classical support-function identities, as in \cite{Schneider}.
Apply them to $P=C+(-A_i)$ and note that reflection preserves
perimeter. The separation assumption justifies using one isolated
response envelope throughout.
\end{proof}

\begin{theorem}[Origin-free rigidity from isotropic pooled responses]
\label{thm:isotropic-registration}
In the isotropic experiment, the exact pooled reciprocal mean pairs
determine the scale ratios, the centered physical first footprint,
and the configurations $\mathcal O-b_i$. No component matching
between settings, preferred direction, homothety origin, periodicity,
finite species list, or boundary smoothness is supplied.

Explicitly, set
\begin{equation}\label{eq:isotropic-ratio}
                   \rho_i=d_i^{\mathrm{iso}}/d_1^{\mathrm{iso}},
\end{equation}
choose $k$ with $\rho_k\ne1$, and form the lower centered support
envelopes $L_i$ in \eqref{eq:registration-envelope}. Then
\begin{equation}\label{eq:isotropic-footprint}
 q=\frac{L_k-L_1}{\rho_k-1}=p_{-A_0},\qquad
                   p_P-\rho_iq=p_{C-b_i}\quad(P\in\mathcal P_i).
\end{equation}
The full geometric ambiguity is exactly
\begin{equation}\label{eq:isotropic-gauge}
 \mathcal O'=\mathcal O+h,\qquad
 A_i'=A_i+h+\pi_i,\qquad \pi_i\in\operatorname{Per}(\mathcal O).
\end{equation}
The relative registrations are recovered modulo the same full period
group as in \eqref{eq:registration-coset}.
\end{theorem}
\begin{proof}
The preceding lemma and homogeneity of perimeter give
\eqref{eq:isotropic-ratio}; the first deficit is positive because
the first footprint has nonempty interior. The global occupation
inverse supplies $\mathcal P_i$ without knowledge of the footprint.
Steiner additivity gives
\[
 L_i(u)=\inf_C p_C^\circ(u)+\rho_i p_{-A_0}(u).
\]
Subtracting at two distinct scales proves
\eqref{eq:isotropic-footprint}, including its physical support
interpretation. No separate convexity of the envelopes or attainment
of their infima is used. Cancelling the centered footprint component
by component gives $\mathcal O-b_i$.

Equality of observed mean pairs for two experiments therefore implies
$\mathcal O-b_i=\mathcal O'-b_i'$ at every setting. Taking
$h=b_1'-b_1$ gives $b_i'-b_i-h\in\operatorname{Per}(\mathcal O)$
and proves necessity of \eqref{eq:isotropic-gauge}. Conversely,
translate each launch density by $h+\pi_i$. Translation of the table
by $h$ and then of a flight by a period preserve each collision bit,
for every direction. Thus the gauge is realized by identical binary
laws, and is sufficient as well.
\end{proof}

The perimeter normalization uses the raw forward mean, not merely its
difference with the reverse mean. The general occupation and period
inverse uses only that difference. Exact homothety and labelled
settings remain assumptions, although their origins and ratios are
recovered. No assertion about finite isotropic boundary minimax rates
is substituted for the proved compass and short-command rates.

\section{An information cost of a fixed common launch density}
\label{sec:stationary-information}
The universal one-bit converse does not measure the information lost
through a fixed random launch spread. We now prove a stronger lower
bound in the stationary experiment itself. Both competing tables use
exactly the same known density and exact implementation. There is no
table-dependent noise map. The proof combines a uniform geometric
bound on the variation of every controlled collision probability with
a conditional binary-channel bound. It applies to expected stopping
times and to continuously selected adaptive commands.

\subsection{Uniform contraction of short-command probabilities}
For two periodically indexed tables with the same lattice and component
correspondences, write
$\epsilon=\max_i d_H(C_i,\widetilde C_i)$ in fixed laboratory frames.
Only the physical packing from Lemma~\ref{lem:packing} will be needed;
there is no matching problem in this definition. Let a fresh launch
law have density $j$ of compact support in a fixed ball and
$\|j\|_\infty\le J$. The density may be given in full to the controller.
Let $t_{\max}$ be a fixed upper bound on command length. Constants below
may depend on $J,t_{\max}$, the footprint envelope and the physical
priors, but not on the command center or the number of observations.

\begin{lemma}[Lipschitz contraction of collision means]
\label{lem:mean-contraction}
For the matched tables above and $0\le\epsilon\le1$,
\begin{equation}\label{eq:mean-contraction}
 \sup_{x,\,|a|\le t_{\max}}
 \left|\int j(z)B_a^{\mathcal O}(x+z)\dd z
       -\int j(z)B_a^{\widetilde{\mathcal O}}(x+z)\dd z\right|
                      \le C J\epsilon.
\end{equation}
The same bound holds for translated reverse commands and for any
parameter-independent mixture of these directions. It holds uniformly
for a family of available launch densities having a common compact
support envelope and a common $L^\infty$ bound.
\end{lemma}
\begin{proof}
For convex compact bodies $A,A'$ at Hausdorff distance at most
$\epsilon$, the inclusions
$A\subset A'+\epsilon\overline B$ and $A'\subset A+\epsilon\overline B$,
together with the planar parallel-area formula, give
\begin{equation}\label{eq:symmetric-difference-area}
 \area(A\mathbin\triangle A')
 \le\{\operatorname{Per}(A)+\operatorname{Per}(A')\}\epsilon
                                                   +2\pi\epsilon^2.
\end{equation}
The bodies and their sweeps $S_a(C)=C+[-a,0]$ have uniformly bounded
perimeters; for example each is at most $\pi$ times its diameter.
Minkowski addition by the same segment does not increase Hausdorff
distance. Thus \eqref{eq:symmetric-difference-area} bounds both solid
and swept symmetric differences by $C\epsilon$ for each matched pair.
These are standard parallel-body facts \cite{Schneider}.

Under the closed-solid convention,
$B_a=\one_{\cup_C S_a(C)}-\one_{\cup_C C}$.
The absolute difference of two union indicators is at most the sum
of absolute differences of the matched member indicators. A fixed
translate of the launch support can meet only a uniformly bounded
number of relevant solids or sweeps: their parent components meet a
ball of prior-bounded radius, and points selected from distinct parent
components are $d_0$-separated. This also bounds components whose matched
partners, rather than themselves, meet that ball. Integrate the indicator
bound and multiply by $J$ to prove \eqref{eq:mean-contraction}.
The estimate is uniform in $x$, so it also applies after the reverse
translation $x\mapsto x+a$. Integrating the bound over a commanded
mixture, or choosing among densities satisfying the same bounds,
preserves it. The noise realization is never part of the observed bit.
\end{proof}

\subsection{A binary channel with a narrow range of means}
All entropies and mutual informations in this section are in bits.
Let $h_2$ denote binary entropy, with its endpoint values defined by
continuity.

\begin{lemma}[Range bound for one binary observation]
\label{lem:binary-range}
Let $V$ have any distribution on a finite set. Conditional on $V=v$,
let $Y$ be Bernoulli with success probability $p_v$. Then
\begin{equation}\label{eq:binary-range}
                     I(V;Y)\le \max_v p_v-\min_v p_v.
\end{equation}
This holds without a positive lower bound on $p_v$ or $1-p_v$ and
also holds conditional on any prior observation history.
\end{lemma}
\begin{proof}
Put $a=\min_vp_v$ and $b=\max_vp_v$. Introduce $U$ uniform on $[0,1]$
and independent of $V$, and generate $Y=\one_{\{U\le p_V\}}$.
Revealing this auxiliary variable can only increase information. Hence
\[
 I(V;Y)\le I(V;Y,U)=I(V;Y\mid U)\le H(Y\mid U).
\]
For $U<a$ or $U>b$ the bit is deterministic independently of $V$.
For $a\le U\le b$ its conditional entropy is at most one.
Integration gives $H(Y\mid U)\le b-a$. A conditional distribution
of $V$ gives the identical proof after every history. This is an
elementary erasure comparison, not a new source-coding principle.
\end{proof}

\begin{proposition}[Stopped experiments with contracted responses]
\label{prop:stopped-contraction}
Consider $M$ parameters and an adaptive binary protocol with independent
controller seed $U_0$ and almost surely finite stopping time $\mathsf T$.
Stopping and commands are measurable from $U_0$ and preceding bits.
Suppose, after each nonterminal history and for every permitted command,
the possible response means lie in an interval of length at most $d$.
Under the uniform parameter distribution, the stopped transcript $Z$
satisfies
\begin{equation}\label{eq:contracted-transcript}
                         I(V;Z,U_0)\le d\,\E\mathsf T.
\end{equation}
If the parameters have disjoint acceptable output sets and the average
error probability is at most $\delta<1/2$, then
\begin{equation}\label{eq:contracted-Fano}
 d\,\frac1M\sum_{v=1}^M\E_v\mathsf T
                 \ge (1-\delta)\log_2M-h_2(\delta).
\end{equation}
An infinite expected cost satisfies the latter conclusion automatically.
\end{proposition}
\begin{proof}
Condition on the independent seed. Pad the transcript with a fixed
symbol after stopping and apply the chain rule up to time $n$.
A stopped history contributes zero, since stopping is determined by the
earlier history and seed. At an active history Lemma~\ref{lem:binary-range}
bounds the conditional information by $d$. Thus the information in the
first $n$ padded entries and the seed is at most
$d\sum_{k=1}^n\Prob(\mathsf T\ge k)=d\E\min(n,\mathsf T)$.
The increasing sigma fields generated by the finite padded records and
the seed generate the stopped record. Since $V$ has a finite alphabet,
conditional probabilities converge by the bounded martingale convergence
theorem; continuity and boundedness of their entropy give convergence
of the conditional entropies. Monotone convergence on the cost side
proves \eqref{eq:contracted-transcript}. Decoding by the unique acceptable
output set and the error-indicator entropy argument used in
Section~\ref{sec:sequential} give
$H(V\mid Z,U_0)\le h_2(\delta)+\delta\log_2M$.
This proves \eqref{eq:contracted-Fano}. Input coordinates, directions,
batch sizes and the stopping time supply no additional information:
conditional on the seed they are functions of the same observed record.
The launch randomness in Lemma~\ref{lem:binary-range} is an auxiliary
comparison variable, not a new physical observation.
\end{proof}

\begin{theorem}[A stationary fixed-spread lower bound]
\label{thm:stationary-lower}
For each $s=6+\beta$, there is a fixed nondegenerate subclass
$\mathcal P\subset\mathcal B_\eta$ and a single fixed, known uniform
disk launch density with the following property. Every adaptive protocol
using fresh draws of that same law, exact nominal centers, and collision
bits from commands of length at most a fixed $t_{\max}$, and having
centered component $C^2$ error at most $\nu$ with probability at least
$1-\delta$ uniformly on $\mathcal P$, satisfies, for fixed
$0<\delta<1/4$ and all small $\nu$,
\begin{equation}\label{eq:stationary-lower}
       \sup_{\mathcal O\in\mathcal P}\E_{\mathcal O}\mathsf T
                            \ge c\nu^{-(s+1)/(s-2)}.
\end{equation}
The lattice, primitive orbit count, component species and the
parameter-independent laboratory frame may all be supplied. The unknown
actual starts are not observed. Arbitrary input precision, adaptive
nominal centers and direction-specific short commands are permitted,
so the lower bound applies in particular to the pooled reciprocal
compass experiment.

The same conclusion holds if the controller may choose a launch scale
$\lambda$ in any fixed interval $0<\lambda_-\le\lambda\le\lambda_+<\infty$;
at each chosen scale the density is the known dilation of the same fixed
uniform disk law. It therefore also bounds the two-scale experiment even
after the footprint and density are disclosed in full.
\end{theorem}
\begin{proof}
Use the fixed lattice, fixed second obstacle, and support-bump family
of Lemmas~\ref{lem:packing} and \ref{lem:centered-packing}.
For every small $h$ it has $M=2^m$ parameters, $m\ge c_0/h$,
pairwise centered $C^2$ separation at least $c_1h^{s-2}$, and
\begin{equation}\label{eq:packing-small-Hausdorff}
                  \max_{\omega,\omega'}d_H(C_\omega,C_{\omega'})
                                         \le C_0h^s.
\end{equation}
The last assertion follows from the disjoint bump supports in
\eqref{eq:packing}; it is a bound in the fixed laboratory frame, not
a disclosure of a parameter-dependent translation. All other component
copies either agree or are translates of these matched bodies.

Choose a fixed $0<t_{\max}<d_0$. Choose once a disk of positive radius
small enough that its largest allowed dilation, together with $t_{\max}$,
obeys the strict separation margin. Its uniform density and all its
allowed dilations have one finite common $L^\infty$ bound and a bounded
support envelope. They also meet the stationary boundary-mass condition
with $\gamma=0$ and a fixed known $b_0>0$. The same law is used under
every parameter; indeed it is disclosed to the controller.

Lemma~\ref{lem:mean-contraction} and
\eqref{eq:packing-small-Hausdorff} imply that the range of every
permitted response probability across the packing has length at most
$d_h=C_2h^s$. This estimate is uniform over all nominal locations,
short directions, scale choices and adaptive histories. Conditioning
on a posterior supported on fewer parameters only shortens that range.
Take $h$ proportional to $\nu^{1/(s-2)}$, with its fixed factor
large enough that $c_1h^{s-2}>2\nu$. The accuracy sets are disjoint,
and Proposition~\ref{prop:stopped-contraction} gives
\[
 C_2h^s\frac1M\sum_\omega\E_\omega\mathsf T
       \ge (1-\delta)m-h_2(\delta)\ge c_3h^{-1}.
\]
The maximum expected cost is at least the average, proving
\eqref{eq:stationary-lower}. If a deterministic cap is imposed, it is
at least the same maximum expected cost. No adversarial implementation,
change of nuisance law between parameters or uncertainty in the period
structure was used.
\end{proof}

The range estimate remains a useful converse independently of the
Hellinger geometry of Section~\ref{sec:sharp-stationary}. For comparison,
let $N^*_{\rm stat}(\nu)$ denote the infimum of worst-case expected
attempts at fixed confidence in the known uniform-disk, fixed-scale
reciprocal experiment on a nondegenerate class containing this packing.
Choose the fixed pooled compass step so that $2t+2R<d_0$, where $R$
is the fixed launch radius; this is the stronger design condition for
the rare-query upper bound. The retained range argument gives
\begin{equation}\label{eq:stationary-bracket}
 c\nu^{-(s+1)/(s-2)}
       \le N^*_{\rm stat}(\nu)
       \le C\nu^{-(3s/2+1)/(s-2)}\log^2(C/\nu).
\end{equation}
Its upper bound is a deterministic high-confidence cap and hence also
an expected-cost bound. Its lower bound holds in the more informative
direction-controlled short-command experiment. The polynomial gap in
this earlier comparison is closed by Theorem~\ref{thm:sharp-stationary-lower}
and Corollary~\ref{cor:stationary-minimax}, which compare the same
physical class, fixed common disk law, confidence criterion and
worst-case expected cost. The shrinking-layer construction also removes
one logarithm from the sufficient attempted count.

The older upper bound of Theorem~\ref{thm:stationary-jitter}, with
exponent $(3s+1)/(s-2)$, continues to hold under its weaker fixed
condition $t+2R<d_0$. The sharp statement concerns the common uniform-disk
model; necessity of the general factor $\gamma+3/2$ for every boundary
mass exponent is not asserted. The converse permits exact nominal
positions, so arbitrary coordinate precision does not remove its
information cost. Its constants are for a fixed positive spread, or a
fixed positive compact scale interval, and are not asserted uniform as
the spread tends to zero. This distinguishes stationary statistical
loss from both shrinking-spread localization and the adversarial
calibration experiment.

\section{Retained reconstruction statements and resource bounds}
\label{sec:localized-statements}
\label{sec:retained-summaries}
We retain the localized experiment, its binary converse and the earlier
stationary bounds. Their proofs follow in the remaining appendices and
in Sections~\ref{sec:stationary}--\ref{sec:stationary-information}.
The stronger stationary bounds proved in Section~\ref{sec:rare-stationary}
use the additional prior-fixed compass design stated there.

\begin{theorem}[Adaptive scalar reconstruction]\label{thm:adaptive}
Fix the numerical bounds of $\mathcal B_\eta$ and the kernel above.
There are $C,c,\nu_0>0$ such that, for $0<\nu<\nu_0$ and
$0<\delta<1/4$, a finite adaptive experiment with the reciprocal
compass commands reconstructs a smooth strictly convex periodic
presentation with geometric loss at most $C\nu$, with probability at
least $1-\delta$. It recovers the true primitive orbit count and exact
rational relations of a true period-generating list in a true independent-pair basis. The Euclidean coordinates of these vectors are estimates.

After a prior-dependent coarse acquisition, the finest localization
scale and sufficient coupled calibration are
\begin{equation}\label{eq:precision}
 \sigma=c\nu^{s/(s-2)},\qquad \ell+\tau+t\alpha\le c\sigma.
\end{equation}
The number of nominal command centers, counted with repetitions if
necessary, and the number of attempted preparations satisfy
\begin{align}
 J_\nu&\le C\nu^{-1/(s-2)}\log(C/\nu),\label{eq:centers}\\
 N_\nu&\le C\nu^{-1/(s-2)}\log(C/\nu)
                               \log(C/(\nu\delta)).\label{eq:attempts}
\end{align}
Every solid start and free miss is charged in $N_\nu$. All commands lie
in one fixed prior-controlled aperture. The reconstruction uses finite
local computations, radial bisections and finite polynomial data, not
a search over an infinite-dimensional class of physical candidates.
Its constants include the exit-time and patch-margin bounds.
\end{theorem}

Here and below a finite smooth partition of unity, with its derivatives,
is part of the fixed numerical representation. Certified evaluations
to the reserved tolerances suffice. The theorem counts preparations and
centers, not apparatus travel or bit operations for elementary functions.
Known $s$ is used; adaptivity refers to the locations of commands, not
adaptation to unknown smoothness.

\begin{theorem}[A necessary number of binary observations]
\label{thm:lower}
For each fixed $0<\beta\le1$ there is a fixed choice of the bounds
above and a nondegenerate subclass $\mathcal P\subset\mathcal B_\eta$
with the following property. Any adaptive randomized procedure whose
only table-dependent outputs are at most $N$ bits, and which recovers
the component boundaries with $C^2$ error at most $\nu$ with probability
at least $3/4$ uniformly on $\mathcal P$, must satisfy
\begin{equation}\label{eq:lower}
                         N\ge c\nu^{-1/(s-2)}.
\end{equation}
This holds even when the lattice, its primitive orbit count, component
correspondences and laboratory pose are given, and implementation is
exact. Input choices may have arbitrary precision and depend on all
previous bits. Independent controller randomness is allowed.
\end{theorem}

The construction of $\mathcal P$ is given in Section~\ref{sec:lower}.
The assertion is not a lower bound for every subclass, in particular
not for a singleton or a finite parametric family. On any bounded class containing this construction, $\eqref{eq:attempts}$ and
$\eqref{eq:lower}$ identify the power of the attempted-bit complexity.
A logarithmic gap remains. Neither result prices input precision or
converts an active experiment into a passive billiard invariant.

\subsection{Finite controls, random stopping and physical resolution}
The next results add two resource converses to the preceding constructive
inverse. The geometric regularity remains $s=6+\beta$, and all uniform
finite global decisions still require the bounded periodic class and
known positive patch margin. No new geometric response is supplied.

\begin{theorem}[Attempted bits and calibration resolution]
\label{thm:resources}
The adaptive reciprocal experiment has a realization with dyadic centers
and scales, using $O(\log(C/\nu)+\log(C/\delta))$ binary digits per
batch command and the same attempted-bit bound
\eqref{eq:attempts}. On a fixed nondegenerate physical subclass, any
procedure whose only table-dependent data are collision bits and whose
stopping decision uses only its past bits and independent randomness
has worst-case expected attempt count at least
$c\nu^{-1/(s-2)}$ at uniform confidence $3/4$.

Moreover, under table-, command- and draw-dependent bounded-position-error
implementations, uniform centered $C^2$ accuracy $\nu$ with confidence greater
than $1/2$ for short commands requires
\[
r\le C\nu^{s/(s-2)}.
\]
Here $r$ is the allowed position error, durations and directions may be
exact, and command lengths have any fixed bound less than the component
separation. Above a constant multiple of this scale,
two distinct physical tables admit identical complete adaptive bit
experiments under two possibly different admissible implementations.
No identical common-noise device is asserted. More attempted
bits do not remove this obstruction. The sufficient combined tolerance
$\ell+\tau+t\alpha\le c\nu^{s/(s-2)}$ remains as before.
\end{theorem}

Theorem~\ref{thm:digital} proves the finite-description claim and gives
a separate conservative numerical cost. Theorem~\ref{thm:expected}
proves the expected-stopping claim by a prefix-free transcript argument,
with the centering gauge explicit. Theorem~\ref{thm:calibration-floor}
proves the physical-resolution converse through a pointwise common-response
coupling of nearby convex swept bodies. Neither converse prices
apparatus production, and neither claims sharp logarithmic factors.
The deterministic-cap lower bound is retained independently: it is
not used as an unproved substitute for the sequential result.

\subsection{A contrasting experiment with fixed stationary jitter}
The preceding resolution converse does not apply to every apparatus
whose start has a bounded displacement. The next theorem recovers the
same table from a fixed, nonzero random launch spread. A known footprint
$K$ has positive curvature, a bounded $C^{6,\beta}$ support function and
$t+\diam K<d_0$. At every nominal center the apparatus draws a fresh
translate of one unknown density $j$, supported on $K$, with
$j(z)\ge b_0\dist(z,\partial K)^\gamma$ almost everywhere in its
interior. Here $b_0>0$ and $\gamma\ge0$ are known. The density is
stationary across all commands; it need not be symmetric or mean zero.
The controller does not observe the draw or its realized start.

\begin{theorem}[Fixed launch spread and unknown density]
\label{thm:fixed-spread-main}
On $\mathcal B_\eta$, under these footprint and density priors, one
finite reciprocal-compass controller, independent of $j$, reconstructs
the geometry to $C^2$ accuracy $C\nu$ and recovers the primitive discrete
data with probability at least $1-\delta$. It uses at most
\[
 C\nu^{-((2\gamma+3)s+1)/(s-2)}
               \log(C/\nu)\log(C/(\nu\delta))
\]
attempted bits in a fixed aperture. All solid starts and misses are
charged. The launch footprint is independent of $\nu$; only the nominal
command precision and probability-estimation accuracy are refined.
For a density bounded below on a fixed disk, $\gamma=0$ and the power
is $(3s+1)/(s-2)$. No optimality of this sample exponent is asserted.
\end{theorem}

Theorem~\ref{thm:stationary-jitter} proves the full statement.
The mechanism is geometric rather than Fourier deconvolution. The
reciprocal means recover a blurred occupation whose positive components
are $C+(-K)$. Subtraction of the known support function of $-K$ recovers
$C$. Lemma~\ref{lem:cap-mass} controls the small probabilities near
this expanded boundary. Proposition~\ref{prop:jitter-modulus} also
gives a H\"older inverse modulus uniformly over different unknown
stationary densities. Theorem~\ref{thm:calibration-floor}, by contrast,
uses two possibly different table-, command- and draw-dependent bounded
implementation maps. These are different uncertainty classes for the
same binary collision sensor, not conflicting universal resolution laws.

\subsection{Unknown footprints and the cost of a common random spread}
The fixed-spread theorem still supplies an entire footprint support
function. We next remove that input in a two-scale experiment, and give
a noise-specific converse without adversarial implementation maps.

\begin{theorem}[Calibrated homothety and stationary information cost]
\label{thm:self-main}
\textup{(i)} Let two fixed positive scales $\lambda_1<\lambda_2$ of one
unknown convex launch footprint be available, with known dilation origin,
uniform inball, envelope and $C^{6,\beta}$ curvature bounds, and
$t+2\lambda_2R_K<d_0$. Suppose one unknown stationary density on the
footprint satisfies the known lower bound
$j(z)\ge b_0\dist(z,\partial K)^\gamma$. On $\mathcal B_\eta$,
finite pooled reciprocal collision bits recover both the table and the
footprint in laboratory $C^2$ error $C\nu$, and the primitive discrete
data, at confidence $1-\delta$. The attempted-bit upper bound is
\eqref{eq:self-attempts}, with the same exponent as the calibrated
one-scale bound. No footprint-function or density oracle is supplied.
Both spreads remain fixed; nominal positioning and exact homothetic
scale control are separate resources.

\textup{(ii)} For one fixed known uniform-disk launch law and exact
implementation, a nondegenerate fixed physical class requires
\[
                  \sup_{\mathcal O}\E_{\mathcal O}\mathsf T
                         \ge c\nu^{-(s+1)/(s-2)}
\]
at fixed confidence. This holds even with known lattice and component
species, arbitrary-precision adaptive nominal positions, and
individually chosen short directions. It also holds when any launch
scale in a fixed interval bounded away from zero can be selected.
The noise law is the same under all competing tables, and no actual
start is observed. This lower bound does not match the stationary
upper exponent; the exact gap is stated in \eqref{eq:stationary-bracket}.
\end{theorem}

Sections~\ref{sec:unknown-footprint} and
\ref{sec:stationary-information} prove these assertions. The first
uses the nested geometric supports at the two scales to identify
component correspondence and to separate the footprint from the body
by support addition. The second proves that the means of every allowed
bit vary by only $O(h^s)$ on a physical packing with $\log_2M$ of
order $h^{-1}$. A conditional binary-channel inequality then applies
through arbitrary stopping rules. Thus the statistical cost is not
inferred from the older worst-case calibration converse or from the
number of bits alone.

\subsection{Retained reconstruction with a fixed homothety origin}
In the stationary experiment, a forward attempt at nominal center $x$
starts at $x+Z$, with an independent compass displacement, and a reverse
attempt starts at $x+Z+a$ with displacement $-a$. The draw $Z$ has an
unknown density on an unknown convex body $K$. The second setting has
footprint $rK$, where the unknown ratio satisfies known bounds
$1+g_*\le r\le R_*$. We take the first operating footprint as the
reference unit of dilation. Both supports have fixed positive size;
they do not contract as accuracy increases. Homothety is about the
specified laboratory origin. The numerical inball, envelope, curvature
and $C^{6,\beta}$ bounds and a lower boundary-mass condition are given.
In particular, for fixed $b_0>0$ and $\gamma\ge0$ the first density obeys
$j(z)\ge b_0\dist(z,\partial K)^\gamma$ almost everywhere in $K$.
Section~\ref{sec:unknown-footprint} gives the detailed geometric priors.

Choose once a dyadic compass step so that the largest footprint
diameter plus $2t$ is less than $d_0$. This stricter design leaves a
positive collar for the translated commands used below. It is a
prior-dependent choice of the same four-direction pooled apparatus.
No direction-resolved response, launch position, occupancy bit or
obstacle area is observed.

The retained fixed-origin constructive conclusion is
Theorem~\ref{thm:unk-scale-finite}. It recovers $K$, $r$ and the original
table to laboratory $C^2$ error $C\nu$, and, on $\mathcal B_\eta$, the
primitive discrete data with probability at least $1-\delta$. Put
\begin{equation}\label{eq:new-upper-exponent}
 A_s(\gamma)=\frac{(\gamma+3/2)s+1}{s-2},\qquad
 L_\nu=\log(C/\nu),\qquad L_{\nu,\delta}=\log(C/(\nu\delta)).
\end{equation}
The sufficient attempted-bit and command-description bounds are
\begin{align}
 N_\nu&\le C\nu^{-A_s(\gamma)}L_\nu L_{\nu,\delta},
                                              \label{eq:main-new-attempts}\\
 \mathcal D_\nu&\le C\bigl\{\nu^{-1/(s-2)}L_\nu
                      +\nu^{-2}L_{\nu,\delta}\bigr\}L_{\nu,\delta}.
                                              \label{eq:main-new-control}
\end{align}
Here $\mathcal D_\nu$ is the binary length of the commanded centers,
setting indices and repetition counts. The finest dyadic nominal mesh
is $O(\nu^{s/(s-2)})$. These statements price the input descriptions
alongside the observations; physical manufacture, travel, homothety
certification and arithmetic running time remain separate resources.
All attempted preparations, including solid starts and misses, are
included in $N_\nu$. The numerical priors are known with certified
enclosures; neither $r$ nor a support-function evaluation oracle for
$K$ is supplied.

There are two new steps. First, near each recovered coarse boundary,
a finite conjunction of shifted pooled collision tests has a zero
probability outside and a uniformly positive probability at depth $e$.
Proposition~\ref{prop:rare-query} proves that its cost is
$O(e^{-(\gamma+3/2)}\log(1/\varepsilon))$. This avoids the squared
inverse-signal cost of the signed occupation query. Second, the integral of a pooled forward mean measures the sum of two
obstacle widths. A linear identity with the two positivity supports
then determines the unknown homothety ratio. Its $O(\nu^{-2})$
sampling cost is absorbed
by \eqref{eq:main-new-attempts}, since $6<s\le7$ and
$A_s(\gamma)>2$. As a second exact argument, a finite killed-walk
adjoint recovers one occupation integral, equal to the obstacle area;
a mixed-area identity determines the ratio from the two reciprocal
differences alone, under a weaker separation condition.

For $\gamma=0$, the common-noise converse in
Theorem~\ref{thm:stationary-lower} gives lower power $(s+1)/(s-2)$,
whereas \eqref{eq:main-new-attempts} has power $(3s/2+1)/(s-2)$.
Thus the difference of the powers is $s/(2(s-2))$; it was $2s/(s-2)$
for the retained signed-query construction. This earlier comparison alone does not identify the minimax exponent.
The sharp common-disk result is now Corollary~\ref{cor:stationary-minimax}. The converse permits arbitrary-precision
adaptive short commands under a single known uniform-disk law and
charges worst-case expected attempts. The finite period assertion
continues to use the known positive nonperiod-patch margin. Exact
geometric recovery and pointwise eventual period recognition are
stated separately from that uniform finite assertion.

\subsection{Location of the retained stationary arguments}
Section~\ref{sec:stationary} develops the stationary occupation inverse
and its quantitative support geometry. Section~\ref{sec:unknown-footprint}
separates two footprints when their scales are known.
Section~\ref{sec:rare-stationary} proves the faster boundary query;
Section~\ref{sec:unknown-scale} identifies an unknown ratio by integrated
collision width and, independently, by occupation area.
Section~\ref{sec:stationary-information} proves the common-noise
converse. The appendices contain the retained localized statements,
coarse reconstruction, smooth interpolation, period tests, physical
packing, finite controls and calibration converse, with their proofs.

\subsection{Exact response rigidity}
The global exact datum is the pair of pooled forward and reciprocal
reverse mean fields, denoted by $F,R$. Their difference $g=F-R$
satisfies a discrete Poisson equation for the stationary occupation
$v$, even when the law of the unobserved launch displacement is unknown.
The positivity components of $v$ are the Minkowski sums of the
obstacles and the reflected launch footprint. Our first result
recovers this unknown positive set as part of the inverse.

For arbitrary locally finite configurations of convex components with
uniformly bounded diameter and positive separation,
Theorem~\ref{thm:global-occupation} gives the constructive formula
\[
 v_0=0,\qquad v_{N+1}=\max\{0,Tv_N-g\},\qquad
                    \|v-v_N\|_\infty\le C e^{-cN}.
\]
The inverse is uniformly Lipschitz on the admissible occupation
class. The proof uses stopping at the exit from a positive component,
whose location need not be known to the reconstruction. Translation
covariance and support cancellation then give
\begin{equation}\label{eq:intro-global-periods}
 \operatorname{Per}(F,R)=\operatorname{Per}(g)
                         =\operatorname{Per}(\mathcal O).
\end{equation}
Thus the full-field scalar experiment detects whether the configuration
is crystalline without a periodicity prior. The exact result requires
only convexity, separation and an almost everywhere positive density
on the footprint interior; smoothness and quantitative boundary mass
are needed later for finite statistical rates.

The second result removes a supplied homothety origin. Let the
stationary supports at finitely many labelled settings be
\[
 A_i=b_i+\rho_i A_0,\qquad s(A_0)=0,\qquad \rho_1=1,
\]
with all $b_i,\rho_i,A_0$ unknown and at least two distinct scales.
Here $s(A_0)$ is the Steiner point. Exact homothety is a relation
between the supports; neither the densities nor the setting translations
are required to follow that relation. For each setting, choose any
complete expanded component $P$ and integrate its isolated forward
response. The deficit
\begin{equation}\label{eq:intro-footprint-deficit}
 d_i=\mathcal W(P)-\frac2t\int_{E_P}F_i
                         =\mathcal W(A_i)
\end{equation}
is independent of that choice. It gives $\rho_i=d_i/d_1$ without
cross-setting correspondence. If $L_i$ is the lower envelope of the
Steiner-centered supports of all expanded components, then
\begin{equation}\label{eq:intro-unregistered-inverse}
 p_{-A_0}=\frac{L_k-L_1}{\rho_k-1}\qquad(\rho_k\ne1).
\end{equation}
The common, possibly infinite obstacle collection cancels from this
identity. Theorem~\ref{thm:registration-rigidity} consequently recovers
each configuration $\mathcal O-b_i$ and all relative registrations
modulo its periods. Corollary~\ref{cor:registration-fiber} proves the
converse as well: the entire geometric ambiguity consists of
\[
 \mathcal O\longmapsto\mathcal O+h,\qquad
 A_i\longmapsto A_i+h+\pi_i,\qquad
                  \pi_i\in\operatorname{Per}(\mathcal O).
\]
This classifies geometric parameters; it does not assert recovery of
the nuisance densities.

\subsection{Sharp stationary powers and finite controls}
Write
\[
 q_\gamma=\frac{(\gamma+3/2)s+1}{s-2},\qquad s=6+\beta.
\]
For the finite smooth classes, assume uniform footprint and
boundary-mass bounds and choose once a dyadic step with
\begin{equation}\label{eq:intro-rare-design}
                     2t+\max_i\diam A_i<d_0.
\end{equation}
This is the stronger separation needed by the shifted rare-collision
query. The older signed-query results remain valid under their weaker
condition $t+\max_i\diam A_i<d_0$. Every rare query is used only in
the protected tubular neighborhood of an already acquired boundary
bracket.

Theorem~\ref{thm:registration-finite} reconstructs unregistered periodic
settings in a common translation frame. A known coarse bound on
$|b_i|$ keeps the physical aperture bounded while leaving the origins
unknown. Uniform finite primitive-period decisions use the original
positive patch margin $\eta$. For a finite configuration entirely
contained in a known protected aperture,
Corollary~\ref{cor:registration-finite-cloud} gives the corresponding
reconstruction without periodicity or a patch-margin assumption.

We distinguish three counts throughout the finite statements:
$N_\nu$ is the number of attempted bits; $J_\nu$ is the number of
center-labelled batches, with a later visit to the same site counted
again; and $S_\nu$ is the number of distinct spatial sites.
Each random center used for an integrated scalar measurement is one
single-attempt batch. In particular $S_\nu\le J_\nu\le N_\nu$.
The new construction gives
\begin{align}
 N_\nu&\le C\nu^{-q_\gamma}\log(C/(\nu\delta))
                          +C\nu^{-2}\log(C/\delta),
                                      \label{eq:intro-new-attempts}\\
 S_\nu\le J_\nu&\le C\nu^{-1/(s-2)}\log(C/\nu)
                          +C\nu^{-2}\log(C/\delta).
                                      \label{eq:intro-new-batches}
\end{align}
The second term is charged in both counts. The batch list, including
centers, settings and repetition counts, has binary length at most
$C J_\nu\log(C/(\nu\delta))$, and the finest nominal mesh is
$O(\nu^{s/(s-2)})$. These are observation and digital-input bounds;
the physical realization of nominal positioning and homothetic
settings remains part of the experiment.

For a fixed known uniform-disk law, the polynomial rate is necessary.
Let $N^*_{\rm pool}$ and $N^*_{\rm short}$ denote the infimal worst-case
expected attempts at a fixed confidence on the same nondegenerate
physical class. The first permits the fixed pooled reciprocal compass
design; the second permits arbitrary adaptive centers and
direction-controlled commands with bounded length, including lengths
approaching zero. The law and its fixed positive scale are the same
under both criteria. Corollary~\ref{cor:stationary-minimax} proves
\begin{equation}\label{eq:intro-sharp-bracket}
 c\nu^{-q_0}\le N^*_{\rm short}(\nu)
       \le N^*_{\rm pool}(\nu)
       \le C\nu^{-q_0}\log(C/\nu).
\end{equation}
The upper controller has a deterministic high-confidence attempt cap,
which is also an expected-cost upper bound. The lower proof works
directly with stopped transcripts and even permits selection among a
fixed positive compact interval of disclosed disk scales.

The main geometric estimate behind the lower bound is a uniform
$O(\epsilon^{3/2})$ squared-Hellinger distance between the collision
responses of nearby support-bump bodies. The two curved boundaries
of a swept collision strip cancel wherever both lie in the launch
disk. The remaining active arcs and sweep facets force a quantitative
amount of collision area and complementary area, uniformly down to
zero command length. A binary-channel capacity inequality then converts
the physical packing into the matching expected-attempt power.
Sharpness of logarithms, confidence dependence, or all other fixed
launch densities is not asserted.

\subsection{Location of these constructions}
The exact occupation and unregistered geometric inverses are proved in
Sections~\ref{sec:global-response} and \ref{sec:unregistered-laws}.
Section~\ref{sec:sharp-stationary} contains the sharp stationary power
and shrinking-layer upper construction. The stationary boundary,
known-scale footprint, rare-query and unknown-scale constructions are
proved in Sections~\ref{sec:stationary}, \ref{sec:unknown-footprint},
\ref{sec:rare-stationary} and \ref{sec:unknown-scale}, respectively.
The independent area normalization uses reciprocal differences alone.
The earlier range converse appears in
Section~\ref{sec:stationary-information}. The localized,
finite-precision, period, packing and calibration arguments are all
included. No proof is required from an external manuscript.

\subsection{Finite response inversion and isotropic commands}
There are two further consequences of the scalar formulation. First,
exact pointwise occupation is a fixed finite-dimensional functional of
the response, rather than merely the limit of an expanding computation.
Theorem~\ref{thm:finite-stencil} expresses it as a stopping linear
program whose coefficients depend only on the diameter and separation
bounds. Its feasible occupation measures have a bounded total mass.
The resulting inverse is Lipschitz even on arbitrary nonphysical noisy
data. Theorem~\ref{thm:stencil-sampling} therefore estimates occupation
to accuracy $\varepsilon$ using $O(\varepsilon^{-2}\log(1/\delta))$
attempts at a fixed finite list of nominal sites. The constants and
site count are given explicitly. The numerical output admits rational
primal and dual bounds. This is an occupation experiment, not a new
membership bit supplied by the apparatus.

Second, Theorem~\ref{thm:general-command-rigidity} replaces the compass
by any bounded centered displacement law with positive second moment.
The exact full period group remains an invariant of the scalar
reciprocal response. For uniform directions, the isolated forward
integral is $t/\pi$ times the obstacle perimeter. The resulting
perimeter deficit identifies the unknown footprint scales without a
preferred laboratory direction. Theorem~\ref{thm:isotropic-registration}
then gives the complete translation ambiguity for independently
translated homothetic supports. This exact isotropic theorem assumes
the full response fields; it does not replace the separately stated
finite compass minimax experiment.

The finite-state stopping representation and Cauchy perimeter identity
are classical. Their role here is to remove, respectively, the
iteration-dependent statistical loss and the directional normalization
from the collision inverse. The geometric issue is that the finite
killing boundary need not be a physical zero set. The containing
component supplies an attaining stopping rule while every other rule
is bounded by nonnegative terminal occupation.


\section{From reciprocal collisions to local occupation queries}
\label{sec:queries}
This section supplies the observation reduction used by every adaptive
query. The random walk is computational: its exit time is not an
observed collision record or an apparatus stopping signal.

\begin{lemma}[Reciprocal balance and calibration]\label{lem:balance}
With $\chi=\one_{\mathcal O}$, outside null boundary sets,
\begin{equation}\label{eq:balance}
 B_a(q)-B_{-a}(q+a)=\chi(q+a)-\chi(q).
\end{equation}
Consequently the difference $g_\sigma$ of the two pooled nominal means
satisfies
\begin{equation}\label{eq:forcing}
 g_\sigma=(T-I)u_\sigma,\qquad
 Tv(x)=\tfrac14\sum_{a\in\{\pm te_1,\pm te_2\}}v(x+a).
\end{equation}
For either command, coupled errors of size
$r=\ell+\tau+t\alpha\le\sigma$ change its mean by at most $Cr/\sigma$.
The constant is uniform over the protected aperture, components and
nominal directions, including tangencies.
\end{lemma}
\begin{proof}
If both endpoints are free, a segment and its reverse either both hit
or both miss. If exactly one endpoint is solid, the direction starting
there gives zero and the other direction hits. If both are solid both
bits are zero. This proves \eqref{eq:balance}. Averaging in $q,a$ and
commuting translation with convolution proves \eqref{eq:forcing}.

For a convex body put $S_a(C)=C+[-a,0]$. The hit bit is
\[
 B_a(q)=\one_{\cup_C S_a(C)}(q)-\chi(q).
\]
The displacement error is at most $\tau+t\alpha$. Support inequalities
show that membership can change only if the nominal start lies in
an $r$-tube of $\partial C$ or $\partial S_a(C)$, including the position
error. The part of a convex boundary tube in a disk of radius $\sigma$
has area at most $C(\sigma+r)r$. One proof integrates boundaries of
outer parallel bodies and inner erosions by the coarea formula; their
portions in the disk have length at most $2\pi\sigma$, by monotonicity
of the perimeter of convex sets. Degenerate levels follow by a limit.
Only a bounded number of components contribute: each has a point in
a disk of radius $\sigma+t+r$, and such points from distinct components
are $d_0$-separated. Multiplication by the density bound
$\sigma^{-2}\|\kappa\|_\infty$ proves the assertion. The tube containment
is pointwise for every allowed error, so the coupling need not have
mean zero. An exceptional coupling event adds its probability.
\end{proof}

Put $D=D_0+2\sigma_0$, $H=(D+t)^2/t^2$ and
$L_* =\max(1,\lceil2H\rceil)$. The positive set of $u_\sigma$ lies in
the $\sigma$-enlarged components. They have diameter at most $D$ and
mutual distance greater than $t$.

\begin{lemma}[Stopped inverse in one aperture]\label{lem:stopped}
For the compass walk $X_n=x+a_1+\cdots+a_n$ and
$\tau_u=\inf\{n\ge0:u_\sigma(X_n)=0\}$ one has
\begin{equation}\label{eq:survival}
 \sup_x\E_x\tau_u\le H,\qquad
 \sup_x\Prob_x(\tau_u>n)\le2^{-\lfloor n/L_*\rfloor}.
\end{equation}
Let a bounded target region $U$ contain the complete protected bodies
and their search brackets. Choose a fixed bounded $W$ containing
$U+\overline B_{D+t+1}$. On $W$ use the killed transition
$T^Wv=T(\one_Wv)$ and set iterates to zero outside $W$. For
$V_0=0$, $V_{n+1}=\max(0,T^WV_n-g_\sigma)$,
\begin{equation}\label{eq:bellman-error}
 0\le u_\sigma-V_n\le2^{-\lfloor n/L_*\rfloor}\quad\hbox{on }U.
\end{equation}
Uniform forcing error $\xi$ and numerical update error $\zeta$ add
at most $n(\xi+\zeta)$ if approximate iterates are clipped to $[0,1]$.
\end{lemma}
\begin{proof}
Until its first zero, a path stays in one enlarged component, since
one step cannot reach another. For $S=\tau_u\wedge n$,
$|X_S-x|\le D+t$. The bounded stopping identity for the martingale
$|X_n-x|^2-nt^2$ gives $t^2\E S\le(D+t)^2$. Monotone convergence
proves the mean bound. Markov's inequality at $L_*$ and the Markov
property at successive blocks prove the geometric tail.

On $W$, $T^Wu_\sigma\le Tu_\sigma$, so $0\le V_n\le u_\sigma$
and the iterates vanish on its zero set. Starting in $U$, the complete
containing component and every first-exit location lie in $W$. Along
this stopped path the killed transition has no boundary discrepancy.
For $e_n=u_\sigma-V_n$, the recursion gives $e_{n+1}\le Te_n$ before
exit, and hence
\[
 e_n(x)\le\E_x[u_\sigma(X_n)\one_{\{\tau_u>n\}}].
\]
This proves \eqref{eq:bellman-error}. Maximum, averaging and clipping
are nonexpansive in supremum norm; induction proves the error bound.
No periodic wrap or renormalization of a boundary row is used.
\end{proof}

The stopped argument also compares two physical occupations with
unrelated zero sets. Indeed, for $w=u-\widetilde u$ and
$h=(T-I)w$, bounded stopping at $\tau_u\wedge n$ gives
\[
 w(x)=\E_xw(X_{\tau_u\wedge n})
                 -\E_x\sum_{j<\tau_u\wedge n}h(X_j).
\]
At exit the terminal value is $-\widetilde u\le0$. On survival it
is at most one. Letting $n$ increase and then interchanging $u$ and
$\widetilde u$ yields $\|u-\widetilde u\|_{U,\infty}
\le H\|(T-I)(u-\widetilde u)\|_{W,\infty}$. Thus exact reciprocal
functionals at all sufficiently small scales determine the local
occupation and, by approximate-identity convergence and regular
closedness, the physical set. This retained exact conclusion requires
neither analyticity nor a supplied zero set. The martingale and killed
Green mechanisms are classical; see \cite{LawlerLimic}.

\begin{proposition}[A finite local query from pooled bits]
\label{prop:query}
For any adaptively chosen $x\in U$ and $\sigma\le\sigma_0$, an
occupation estimate with error less than $1/8$ can be computed using
only a prior-bounded number of reciprocal command centers in $W$.
At conditional confidence $1-\varepsilon$ it uses at most
$C\log(C/\varepsilon)$ attempted preparations. A sufficient calibration
is $\ell+\tau+t\alpha\le c\sigma$. Thresholding the estimate at
$1/2$ gives the correct membership label outside the $\sigma$-tube
of the component boundaries. Its label inside that tube is unrestricted.
\end{proposition}
\begin{proof}
Fix $n_*=6L_*$. To evaluate $V_{n_*}(x)$, its dependency locations
are exactly among
\begin{equation}\label{eq:diamond}
 x+t(i,j),\qquad |i|+|j|\le n_*.
\end{equation}
Intersect these with $W$ and put zero at locations outside it. Evaluate
the recursion backwards through shrinking diamonds, initializing
$V_0=0$. Induction on the layer shows equality with the full killed
iterate at the target. There are $O(n_*^2)$ centers and $O(n_*^3)$
updates, both independent of $\sigma$. No fine background spatial grid,
interpolation of the transition, or growing aperture is involved.

At each center estimate each pooled mean to error at most
$1/(1024n_*)$, and impose the same bound on its calibration bias.
The forcing error is then at most $1/(256n_*)$. Reserve update error
at most $1/(256n_*)$. Lemma~\ref{lem:stopped} gives total error at most
$1/64+1/256+1/256<1/8$. Hoeffding's inequality and a union bound over
the fixed finite number of means give the preparation bound. The
constant includes $n_*$ and may be large. The proof holds conditional
on past data because the new attempts are fresh. Inside a component
at boundary distance greater than $\sigma$ the average is one;
outside all enlarged components it is zero. Thresholding therefore
has the stated property. A membership label is a computed conclusion
from these bits, not an extra response furnished by the sensor.
\end{proof}

The full launch region may be taken to contain
$W+\overline B_{t+\sigma_0+1}$, which includes reverse positions and
kernel supports. For at most $Q$ adaptive queries, allocating conditional
failure probability $\delta/Q$ to each gives joint correctness with
probability $1-\delta$. It suffices to know a deterministic upper bound
on $Q$; repeated centers may be resampled. No independence of the
adaptively selected locations themselves is asserted.

\section{Adaptive recovery of complete component boundaries}
\label{sec:boundary}
We first acquire a coarse, fixed-resolution description. Its purpose
is to find interior centers and isolated radial brackets, not to supply
the final accuracy. This avoids both a supplied free reference point
and a bisection that might accidentally encounter a different obstacle.

\begin{lemma}[Coarse hulls and interior centers]\label{lem:coarse}
A prior-dependent finite number of queries from
Proposition~\ref{prop:query}, at one fixed scale $\sigma_c>0$, recovers
all complete components needed in the protected patch to Hausdorff
error $e_c\le3\sigma_c$. It supplies, for each such component $C$, a
polygon $H$ and a computable center $o$ for which
\begin{equation}\label{eq:inner-ball}
 \overline B_r(o)\subset C\cap H\subset\overline B_{R_1}(o)
\end{equation}
with fixed positive $r,R_1$. For every direction $v$, it supplies a
radial interval of fixed length at most $C e_c$ containing the boundary
point of $C$ on $o+\R_+v$. The interval is disjoint from the
$\sigma$-enlargements of all other components for every sufficiently
small fine scale $\sigma$.
\end{lemma}
\begin{proof}
The curvature upper bound supplies an interior rolling radius
$r_* =1/\kappa_+$. In support coordinates this follows from
$p+p''\ge r_*$: subtracting the support function of the radius-$r_*$
disk leaves a support function with nonnegative curvature measure.
In particular $C=C_{-2\sigma_c}+2\sigma_c\overline B$ when
$2\sigma_c<r_*$. Choose $\sigma_c<\min(r_*/8,d_0/40)$ and a square
grid of covering radius $\rho_c\le\sigma_c/8$ on a surrounding
square with collar larger than $2D_0+2$.

At its nodes retain the estimated occupations at least $1/2$, join
retained nodes at distance at most $4\sigma_c$, and take a convex hull
for each connected cluster. Every retained node is within $\sigma_c$
of a unique body; all nodes of $C_{-\sigma_c}$ are retained. For any
retained node assigned to $C$, choose a point in $C_{-2\sigma_c}$ at
distance at most $3\sigma_c$, and then a nearest grid node. That node
is retained and joined to the first. Convexity connects any two such
deep nodes by a segment; partitioning it into lengths at most $\rho_c$
and replacing partition points by nearest nodes yields a retained
path with successive lengths at most $3\rho_c$. Different bodies
cannot join because $d_0-2\sigma_c>4\sigma_c$. The hull thus lies in
$C+\sigma_c\overline B$, whereas every point of $C$ is within
$2\sigma_c+\rho_c$ of a retained deep node. This proves the error.
Discard clusters whose hulls are within $D_0+1$ of the outer square
boundary. Fragments cannot survive this cutoff, and the chosen collar
protects every required complete component.

Compute a largest inscribed disk of $H$ by its finite half-plane
inequalities, to any fixed reserved precision. Since
$d_H(H,C)\le e_c$, the support inequalities imply
$C_{-e_c}\subset H$ and $H_{-e_c}\subset C$. A disk of radius $r_*$
in $C$ therefore gives an inscribed disk of radius at least $r_*-e_c$
in $H$. The computed center, with the radii reduced by the error and
numerical reserve, satisfies \eqref{eq:inner-ball} for a fixed $r>0$.
Diameter bounds supply $R_1$.

Write $R_H(v),R_C(v)$ for the radial functions about $o$. Since both
bodies contain $\overline B_r(o)$,
\[
 H+e_c\overline B\subset o+(1+e_c/r)(H-o),
\]
and the analogous inclusion holds for $C$. The two Hausdorff
inclusions give $|R_H(v)-R_C(v)|\le C e_c$ uniformly. The polygonal
radial value is computable from its sides. Use a slightly enlarged
bracket about it, allowing numerical error. By decreasing the fixed
$\sigma_c$, its length and every point's distance to $C$ are less
than $d_0/4$. Separation excludes every other component from the
bracket and its sufficiently small kernel neighborhood.
\end{proof}

\begin{lemma}[Bisection with an indeterminate boundary layer]
\label{lem:bisection}
On each bracket in Lemma~\ref{lem:coarse}, the derived label of
Proposition~\ref{prop:query} may be arbitrary, and may depend on the
past, only when
\begin{equation}\label{eq:radial-ambiguity}
                       |a-R_C(v)|\le A\sigma
\end{equation}
for a fixed $A$. After $k$ midpoint tests, ordinary bisection returns
an estimate with error at most $A\sigma+2^{-k-1}w_0$, where $w_0$
is the initial bracket length. In particular $O(\log(C/\sigma))$
queries recover $R_C(v)$ to $O(\sigma)$.
\end{lemma}
\begin{proof}
Translate $o$ to zero. The inner ball implies
\[
 (1-\sigma/r)C\subset C_{-\sigma},\qquad
 C+\sigma\overline B\subset(1+\sigma/r)C.
\]
Together with $R_C\le R_1$, these show that outside
\eqref{eq:radial-ambiguity} the label is the true radial inside/outside
label. Other components were excluded from the bracket. A label one
at radius $a$ implies $a\le R_C+A\sigma$; a label zero implies
$a\ge R_C-A\sigma$. If the current endpoints are $l,u$, midpoint
updates consequently preserve the relaxed invariant
$R_C\in[l-A\sigma,u+A\sigma]$. This does not require $R_C$ to remain
in $[l,u]$, or the noisy labels to be monotone. The width halves at
every step. Its midpoint has the asserted error.
\end{proof}

\begin{lemma}[Radial regularity and stable interpolation]
\label{lem:radial}
Under \eqref{eq:inner-ball} and the physical support priors, the
$2\pi$-periodic radial function $R_C(\varphi)$ has a uniformly bounded
$C^{6,\beta}$ norm. Let $h=2\pi/m$ for an integer $m\ge9$. From
radial values at the $m$ equally spaced angles, each with error at most
$e$, one can construct a smooth function $\widehat R$ such that
\begin{equation}\label{eq:radial-interpolation}
 \|\widehat R-R_C\|_{C^j}
           \le C(h^{s-j}+e h^{-j}),\qquad 0\le j\le3.
\end{equation}
For $e\le C_0h^s$ and sufficiently small $h$, the curve
$o+\widehat R(\varphi)(\cos\varphi,\sin\varphi)$ is the boundary of
a smooth strictly convex body $\widehat C$, and
\begin{equation}\label{eq:geometric-interpolation}
 d_H(\widehat C,C)\le Ch^s,\qquad
 \|p_{\widehat C}^{\rm lab}-p_C^{\rm lab}\|_{C^2}\le Ch^{s-2}.
\end{equation}
Here uncentered support functions are compared in the same laboratory
frame; centering preserves these orders.
\end{lemma}
\begin{proof}
Use the support function $p$ relative to $o$, so $p\ge r$. In normal
angle $\theta$, the boundary is $p n+p' n^\perp$, its radial angle is
\[
 \varphi=\theta+\arctan(p'/p),\qquad
 \frac{\dd\varphi}{\dd\theta}
             =\frac{p(p+p'')}{p^2+(p')^2}>c>0.
\]
The inverse $\theta(\varphi)$ is uniformly $C^{5,\beta}$. Although
the expression $R=\sqrt{p^2+(p')^2}$ initially gives only five
ordinary derivatives, differentiation in radial angle cancels the
curvature factor:
\begin{equation}\label{eq:radial-cancellation}
                   R'=R(p'/p)\circ\theta.
\end{equation}
Its right side is uniformly $C^{5,\beta}$. Thus $R$ is uniformly
$C^{6,\beta}$, without losing a derivative from the support prior.

At each angular node interpolate the seven consecutive sampled
values by a polynomial of degree at most six in a local lift of the
angle. Combine these polynomials with a fixed translated smooth
partition of unity supported on neighborhoods of size $O(h)$.
On each such neighborhood the normalized Vandermonde inverse is
bounded. A degree-six Taylor polynomial of the true function is
reproduced exactly, and its remainder and derivatives through order
three are $O(h^{s-j})$. Errors in the values contribute $O(eh^{-j})$.
Derivatives falling on the partition have size $O(h^{-j})$ and are
absorbed by the corresponding lower-order remainders. The overlap
number is fixed. This proves \eqref{eq:radial-interpolation}, including
across the angular cut by periodic lifts.

A positive radial function is strictly convex when
\[
       F=R^2+2(R')^2-RR''>0.
\]
For the true curve this expression has a uniform positive margin,
by the inradius, curvature and diameter bounds. It persists for
$\widehat R$ by the $C^2$ estimate. The two radial graphs are uniformly
$O(h^s)$-close, which proves the Hausdorff bound.

To check the stronger support norm, set
$\theta=\varphi-\arctan(R'/R)$, now viewing the normal angle as a
function of radial angle. Its derivative is
$F/(R^2+(R')^2)>c$. At corresponding angles the support value and
radius of curvature are
\[
 p=\frac{R^2}{\sqrt{R^2+(R')^2}},\qquad
 p+p''=\frac{(R^2+(R')^2)^{3/2}}{F}.
\]
These expressions, and the tangential support derivative, depend
Lipschitzly on $R,R',R''$ on the bounded set with these margins.
The inverse normal-angle maps differ by $O(\|\widehat R-R\|_{C^1})$.
Their composition changes the true expressions by the same order,
using the uniform $C^3$ bound; the reconstructed functions also have
such a bound by \eqref{eq:radial-interpolation}. Thus the support
error through two derivatives is $O(h^{s-2})$. Translation to the
laboratory frame adds the same first harmonic to both supports. The
Steiner point is a bounded linear functional of the support, so its
error is $O(h^s)$ and centering preserves the conclusions.
\end{proof}

\begin{proposition}[Adaptive patch reconstruction]\label{prop:patch}
For a fixed protected patch, complete component boundaries can be
recovered with the errors in \eqref{eq:geometric-interpolation} using
at most $C h^{-1}\log(C/h)$ local occupation queries. The finest
localization is $\sigma=c h^s$. With probability at least $1-\delta$
the attempted-bit cost is at most
\begin{equation}\label{eq:patch-cost}
           C h^{-1}\log(C/h)\log(C/(h\delta)).
\end{equation}
The number of components, the initial coarse queries and every
local Bellman stencil size are bounded solely by the fixed prior.
\end{proposition}
\begin{proof}
On the simultaneous accuracy event, obtain the hulls and centers in
Lemma~\ref{lem:coarse}. Separation and the fixed aperture bound their
number. For each body use $m$ equally spaced radial directions and
their isolated coarse brackets. Lemma~\ref{lem:bisection} uses
$O(\log(C/h))$ midpoint queries per direction to return values with
error $e\le C\sigma$. Lemma~\ref{lem:radial} gives the claimed
geometric estimates. The total number of queries is bounded by a
deterministic $Q_h=C(1+h^{-1}\log(C/h))$. Allocate failure probability
$\delta/Q_h$ to every query, including the coarse ones, before the
experiment. Conditional confidence and the union bound in
Proposition~\ref{prop:query} give the simultaneous event. Its
preparation cost is \eqref{eq:patch-cost}. Locations chosen using
coarse or previous fine data do not alter this argument.

For a finite numerical realization, polygon intersections, midpoint
locations and radial values are enclosed with error $O(\sigma)$.
The estimate \eqref{eq:radial-interpolation} absorbs those value errors.
Represent the output by its finite polynomial coefficients and the
fixed smooth partition functions. They, their derivatives and the
normal-angle maps can be evaluated to any positive prescribed error
by finite operations on compact intervals with the stated derivative
bounds. This constructs an actual smooth convex output, without
assigning a false smoothness bound to the initial polygons. No
infinite-dimensional candidate minimization is performed.
\end{proof}

\section{Finite binary descriptions of the adaptive controls}
\label{sec:digital}
The active protocol specifies a distribution by a finite command, rather
than transmitting an exact real coordinate. We separate this digital
resource from the physical production of a localized launch distribution.
The latter remains a calibrated apparatus primitive; it is not obtained
by observing the collision bits.

\begin{theorem}[A dyadic controller and its resource bounds]
\label{thm:digital}
On the same bounded class $\mathcal B_\eta$, choose the compass length
$t>0$ and the fixed coarse numerical reserves dyadic, while retaining
all strict separation margins. The experiment of
Theorem~\ref{thm:adaptive} can be implemented with dyadic nominal
centers and dyadic localization scales. Its accuracy, deterministic
attempt bound and exact discrete outputs are unchanged up to constants.
Put
\[
 b=\lceil C_0+\log_2(1/\sigma)\rceil,
 \qquad B=1+\lceil\log_2(C/\nu)\rceil
              +\lceil\log_2(C/\delta)\rceil .
\]
The finest center mesh is $2^{-b}\le c\sigma$, with
$b=O(\log(C/\nu))$. Each batch command has a binary description of
length at most $CB$, and the total transmitted command description is
at most $CJ_\nu B$. The finite boundary representation has length
at most $C\nu^{-1/(s-2)}B$. These bounds include batch repetition
counts. They do not count the analog random numbers internal to the
prescribed launch distribution as table-dependent observations.

For rational fixed numerical priors and rational $\beta$, one may
choose the fixed smooth partition and its evaluators so that a
conservative digital bit-operation bound is
\begin{equation}\label{eq:digital-work}
              C(N_\nu+J_\nu+\nu^{-1/2})B^c
\end{equation}
for an absolute finite exponent $c$. This includes boundary output,
primitive arithmetic and free-area evaluation to error $C\nu$.
It is not a sharp computational complexity bound. For general fixed
$\beta$ the message bounds hold with certified dyadic scale enclosures;
a bit-operation bound additionally depends on their computation cost.
\end{theorem}
\begin{proof}
\emph{One exact dyadic stencil.}
Take the fine localization to be a dyadic number comparable to $c h^s$,
where $h=2\pi/m\asymp\nu^{1/(s-2)}$. Choose a dyadic rectangle $W$
with all the protective collars of Lemma~\ref{lem:stopped}. Round a
requested target $x$ once to a dyadic center $x'$, with
$|x-x'|\le C2^{-b}$. Use the exact stencil
\[
                         x'+t(i,j),\qquad |i|+|j|\le n_*.
\]
For $b$ greater than the fixed denominator exponent of $t$, all stencil
points are on the same dyadic grid. Membership in $W$ is tested by
exact rational comparisons. There is no independent rounding of
transition targets, hence no error in the killed compass identity.
The reverse nominal distribution is exactly the translated joint law:
it attempts $(q+a,-a)$ for the same specified law of $(q,a)$, not an
independently rounded reverse start. Fresh attempts are still used.

The query estimates $u_\sigma(x')$. Its derived label at the intended
point $x$ is therefore correct outside the boundary tube of width
$\sigma+C2^{-b}$. Physical position, duration and angle errors receive
an independent portion of the reserve in
Proposition~\ref{prop:query}. Thus rounding and actual calibration
together preserve its fixed accuracy bound. The means are ratios of
integer hit counts to a deterministic batch size. Taking that size to
be a common integer upper bound for all queries makes every local
average and Bellman update an exact rational computation. Clipping and
comparison with $1/2$ are also exact. The depth and stencil size are
prior constants.

\emph{Rounded adaptive geometry.}
The coarse grid and its finite polygons have a prior-bounded number
of rational vertices. A dyadic interior center with a reserved fixed
inball can be found by finite search; squared distances to rational
sides test the inball inequalities. The strict inball reserve ensures
that a sufficiently fine fixed search grid succeeds. Approximate
$\cos(2\pi j/m)$ and $\sin(2\pi j/m)$ with certified error at most
$c2^{-b}$. Polygonal radial values and brackets are computed with the
same reserve. Consequently every implemented spatial midpoint is
within $C2^{-b}$ of the intended radial point.

For completeness, rounding a radial midpoint by at most $d$ does not
accumulate $kd$ after $k$ queries. The relaxed interval invariant from
Lemma~\ref{lem:bisection} holds with a boundary layer of width
$A(\sigma+d)$. Its interval widths satisfy
\[
                     w_{k+1}\le w_k/2+d,
 \qquad w_k\le2^{-k}w_0+2d.
\]
Stop before asking for resolution finer than the mesh. With
$d=C2^{-b}\le c\sigma$, $O(\log(C/\sigma))$ queries give radial
values of error $O(\sigma)$, even when labels in the enlarged layer
are chosen adaptively. Round the final values to the same mesh.
Lemma~\ref{lem:radial} absorbs these errors and still gives actual
smooth strictly convex output bodies. The known patch margin and
bounded-denominator gaps allow the same exact period decisions.

\emph{Description lengths.}
All centers lie in a fixed rectangle. Their dyadic integer numerators
therefore use $O(b)$ bits. A batch command consists of a forward/reverse
flag, two such coordinates, a scale exponent, the fixed compass/kernel
identifier, and a repetition count of order $\log(C/(\nu\delta))$.
The last count requires only $O(\log B)$ binary digits. Fixed-width
fields give the asserted $CB$ bound without a real-valued control
field. There are at most $CJ_\nu$ batches, allowing repetitions.
The receiver interprets a fixed kernel identifier as the stipulated
localized random preparation, not as a table-dependent data stream.

Store the dyadic radial values, the dyadic centers and the integers
$j,m$. In normalized coordinates the seven-point interpolation matrix
has fixed rational entries, so its coefficients have $O(B)$-bit
rational descriptions. Local angles and partitions are generated by
a fixed program from $j,m$; no table-dependent exact trigonometric
constant is stored. There are $O(m)$ coefficient blocks, proving the
boundary-description bound.

\emph{A conservative operation count.}
For the computational assertion choose an explicit smooth bump,
for example $\exp[-1/(1-x^2)]$ on $|x|<1$, zero outside, and normalize
a finite-overlap family of its translates to make the fixed partition.
The denominator has a positive fixed lower bound. To the required
precision, the bump and its first three derivatives are evaluated by
rational series on bounded arguments; near its endpoints an explicit
exponential tail bound permits certified truncation to zero. The
trigonometric functions, $\pi$, square roots and logarithms needed
here admit rational series and bisection with a polynomial number of
bit operations in $B$. For rational $s$, dyadic powers comparable to
$h^s$ can equivalently be bracketed by raising rational quantities to
fixed integer powers. Thus all scale choices are effective in this
computational subcase.

Processing bit counts, the constant-depth rational Bellman updates,
rounded bisection and the fixed-degree coefficient blocks costs
$C(N_\nu+J_\nu)B^c$. One need not optimize global numerical work to
obtain a bound: the reconstructed support functions have uniform
$C^2$ bounds. Composite trapezoidal quadrature with
$O(\nu^{-1/2})$ nodes computes their Steiner integrals, and the area
integrals $\tfrac12\int\widehat R^2$, to error $O(\nu)$. A support
value is evaluated by inverting the uniformly monotone normal-angle
map; bisection with residual enclosures uses $O(B)$ steps. Each
partition evaluation involves only a fixed number of neighboring
coefficient blocks. The uniform derivative and convexity margins
justify all these tolerances.

The number of protected components and candidate period pairs is a
prior constant. Support-defect comparisons need only a fixed precision
smaller than the known patch and rational-separation margins, whereas
the real period coordinates and areas are evaluated to $O(\nu)$.
Their mesh tests, bounded-denominator rational search and two-dimensional
Hermite arithmetic are finite with polynomial cost in $B$. Including
the preceding quadratures proves \eqref{eq:digital-work}. The larger
$\nu^{-1/2}$ term is a disclosed conservative numerical overhead; it
is not silently identified with the attempted-bit exponent.
\end{proof}

\subsection{Separate resource axes}
Write $m_\nu=\nu^{-1/(s-2)}$ and $L_\nu=\log(C/\nu)$. The following
accounting refers to the construction, not to a measured apparatus.
\begin{center}
\small
\begin{tabular}{@{}p{0.31\textwidth}p{0.64\textwidth}@{}}
\toprule
Resource & Bound or explicit status\\
\midrule
Table-dependent outputs & $N_\nu\le C m_\nu L_\nu\log(C/(\nu\delta))$ bits;
all attempts, including zeros, are charged.\\[3pt]
Command centers/batches & $J_\nu\le C m_\nu L_\nu$, counted with repetitions.\\[3pt]
Digital controls & $O(B)$ bits per batch; $O(J_\nu B)$ in total.\\[3pt]
Spatial localization & $\sigma\asymp\nu^{s/(s-2)}$; the localized kernel
is a prescribed preparation primitive.\\[3pt]
Physical calibration & $\ell+\tau+t\alpha\le c\sigma$ is sufficient.
Section~\ref{sec:calibration-floor} gives a matching necessary position-error power.\\[3pt]
Aperture & Fixed by the geometric prior, including reverse starts and
protective collars.\\[3pt]
Numerical bit operations & The conservative bound \eqref{eq:digital-work}
under its explicit effective-prior assumptions.\\[3pt]
Motion, setup, calibration production & No physical time, energy or
cost law is assumed. No bound on these resources is inferred from $N_\nu$.\\
\bottomrule
\end{tabular}
\end{center}
Digital description length is not mechanical resolution. In particular,
a finite word specifying a narrow kernel does not prove that a device
can produce that kernel cheaply. The converse calibration theorem
below quantifies a necessary tolerance in the observation model,
without inventing a law for the cost of achieving it.

\section{Period recognition from the reconstructed patch}
\label{sec:periods}
This section retains the protected-patch argument for repeated shapes.
It acts on the acquired geometry, not on an assumed list of integer
copy labels. We include the proof to specify exactly which discrete
information the adaptive experiment recovers.

Put $B=1+L_0$ and $Q_a=[-a,a]^2$. The aperture used above is chosen
large enough to recover complete components with centers in $Q_{10B}$,
together with their protective collars. For two components define
\[
 d_*(C,D)=\|z_C-z_D\|_\infty+\|p_C-p_D\|_\infty,
\]
where the centered support functions are compared in the laboratory
directions, without an independent rotation. Define the patch defect
\begin{equation}\label{eq:defect}
 \mathcal D(w)=\max_{C:z_C\in Q_B}\ \max_{\epsilon\in\{-1,1\}}
                     \inf_D d_*(C+\epsilon w,D).
\end{equation}
The infimum ranges over physical components. Distant components cannot
minimize it. The known nonperiod margin is the requirement
\begin{equation}\label{eq:patch-margin}
 \begin{gathered}
 w=z_D-z_C,\quad z_C\in Q_{3B},\ z_D\in Q_{5B},\quad w\notin\Pi
       \quad\Longrightarrow\quad\mathcal D(w)\ge\eta>0.
 \end{gathered}
\end{equation}
This is a margin for whole-patch translations, not a separation between
all individual shapes. Every fixed table in $\mathcal B$ has some
positive such margin, since there are finitely many candidate pairs.
The margin is not claimed to be uniformly inferable from finite bits.

\begin{lemma}[Exact patch recognition]\label{lem:recognition}
One has $\mathcal D(w)=0$ if and only if $w\in\Pi$. For
$\|w\|_\infty\le8B$, this zero test uses only components centered in
$Q_B$ and $Q_{9B}$. The group $\Pi$ is a lattice with
\begin{equation}\label{eq:index-bound}
 [\Pi:L\Z^2]\le r_0,\qquad\covol\Pi\ge V_*=V_0/r_0.
\end{equation}
Choose any component centered in $Q_{2B}$ as a root. Its center
differences to components centered in $Q_{4B}$ that pass the zero
test generate $\Pi$ over $\Z$.
\end{lemma}
\begin{proof}
Each orbit modulo $L\Z^2$ has a representative centered in
$L[-1/2,1/2]^2\subset Q_B$. A zero defect gives both translates of
every such representative as physical components. Translating by
$L\Z^2$ gives $\mathcal O+w\subset\mathcal O$ and
$\mathcal O-w\subset\mathcal O$; together these imply equality.
The converse is immediate. The centers of a matching pair in the
bounded test are in $Q_{9B}$.

Map a coset of $\Pi/L\Z^2$ to the component orbit of the translated
root. This is injective: a nonempty compact body is not invariant
under a nonzero translation. Thus the quotient has at most $r$ elements.
A finite extension of a full-rank lattice is a lattice, giving
\eqref{eq:index-bound}. The two columns of $L$, and representatives
of all these period cosets, can be chosen with norm less than $B$
in the maximum norm. Translating the root by these vectors gives
targets in $Q_{3B}$, strictly within the cutoff $Q_{4B}$. They are
included in the accepted list and generate $\Pi$. Conversely each
accepted vector is a true period.
\end{proof}

\begin{lemma}[Stable discrete decisions]\label{lem:locking}
Suppose the protected complete components have matched Hausdorff
error at most $e$. On $\mathcal B_\eta$, for sufficiently small $e$
depending on the full prior, finite tests recover the true primitive
orbit count and exact rational relations of a true accepted generating
list in a true independent-pair basis. They recover a matched period
basis and primitive free area to error $O(e)$.
\end{lemma}
\begin{proof}
Hausdorff error $e$ gives Steiner-center error at most $2e$ and
centered-support error at most $3e$. Use approximate source centers
in $Q_{2B}$, a root there and root targets in $Q_{4B}$. For a measured
difference $\widehat w$, its corresponding true difference has error
at most $4e$. Test both signs on all selected sources, minimizing over
targets with approximate centers in $Q_{10B}$. The center contribution
to each comparison changes by at most $8e$ and the support contribution
by at most $6e$. True periods therefore have measured defect at most
$14e$. For a nonperiod, its true endpoints lie in $Q_{3B},Q_{5B}$,
so \eqref{eq:patch-margin} supplies a witness of defect at least $\eta$.
All witnesses centered in $Q_B$ are selected. Restricting targets cannot
lower the true infimum. Its measured defect is at least $\eta-14e$.
Reserve another $6e$ for numerical enclosures and use threshold $\eta/2$,
with $20e<\eta/4$. This accepts exactly the true periods. Strict cutoff
collars ensure that every generator in Lemma~\ref{lem:recognition}
is present. The same tests on pairs of sources give their primitive
orbit relation: a period translating the first center to the second
must translate the first body to the second body, since Steiner points
lie in the interiors and different components are disjoint. All
primitive orbits are represented, so their number is recovered.

Accepted true vectors have a prior bound $U_0$ on their Euclidean
lengths. A nonzero determinant of two such periods is an integer
multiple of $\covol\Pi\ge V_*$. Determinant perturbations are
$O(U_0e)$; hence independence is decided below this gap. Select an
independent true pair with matrix $D$. Its generated sublattice has
index
\[
       n=|\det D|/\covol\Pi\le Q:=\max(1,\lceil U_0^2/V_*\rceil).
\]
Every coordinate of $D^{-1}w_j$ has reduced denominator dividing $n$
and a prior-bounded numerator. The inverse-matrix identity bounds its
estimation error by $C_*e$. Distinct rationals of denominators at most
$Q$ are separated by at least $Q^{-2}$. Finite search locks the exact
coordinates when $C_*e<1/(3Q^2)$.

Let $q$ be their common denominator, which divides $n$. A column
Hermite basis $H_0$ of the integer span of $qI_2$ and $qD^{-1}w_j$
gives
\begin{equation}\label{eq:period-basis}
                         \Pi=(DH_0/q)\Z^2.
\end{equation}
The integer group contains $q\Z^2$, so its Hermite entries have a
prior bound. Replacing $D$ by its estimate gives a matched true basis
with error $O(e)$. We have recovered exact integer relations, not
exact real coordinates. The primitive free area is
\[
             A_* =\covol\Pi-\sum_{j=1}^{r_*}\area(C_j).
\]
Uniform perimeter bounds and the parallel-body inclusions give area
error $O(e)$ for each matched component, proving the last assertion.
Support suprema and center integrals used in the tests admit certified
finite angular meshes, because the bounded bodies have a uniform
support Lipschitz bound. Their numerical errors are charged above.
\end{proof}

\begin{proof}[Proof of Theorem~\ref{thm:adaptive}]
Choose a sufficiently large fixed target square so that every component
needed in Lemma~\ref{lem:locking} and every coarse bracket is protected.
Separation bounds their number. Choose the fixed observation and launch
apertures as in Lemma~\ref{lem:stopped}. Apply
Proposition~\ref{prop:patch} with an integer angular mesh satisfying
$h\asymp\nu^{1/(s-2)}$ and a sufficiently small constant factor.
Its reconstructed bodies have matched $C^2$ support error $O(\nu)$
and Hausdorff error $e=O(h^s)$. For small $\nu$, all thresholds in
Lemma~\ref{lem:locking} hold. It gives exact primitive discrete data
and period-basis error $O(h^s)$. Repeating one recovered body per
primitive orbit under this basis gives a physical periodic presentation.
Indeed the inverse basis remains bounded, so only a bounded number of
integer translates can enter a bounded fundamental patch. The positive
separation and curvature margins then persist. The reconstructed
bodies have area error $O(h^s)$ and the covolume has the same order;
hence the primitive free-area error is also $O(h^s)$, in particular
$O(\nu)$.

The query count of Proposition~\ref{prop:patch} and its fixed local
stencil size give \eqref{eq:centers}. Equation~\eqref{eq:patch-cost}
gives \eqref{eq:attempts}, since $s$ is fixed. The calibration bias
allocation in Proposition~\ref{prop:query} gives \eqref{eq:precision}.
Every preparation has already been charged there. All steps are finite,
and the output is a finite smooth representation with certified
numerical tolerances. This proves the theorem.
\end{proof}

Without a known $\eta$, successive refinements give the corresponding
pointwise eventual decisions for each fixed table, but not a uniformly
valid observable stopping certificate. The exact functional inverse
combined with Lemma~\ref{lem:recognition} needs no supplied positive
margin: it tests zeros exactly. This distinction, as well as the
bounded periodic presentation itself, is retained for this finite
protected-patch argument. The exact full-field criterion for
crystallinity, without a periodicity prior, is
Theorem~\ref{thm:global-response-rigidity}.

\section{A physical packing and a binary-transcript bound}
\label{sec:lower}
The lower bound counts the table-dependent binary outputs, not the
precision of the controller's choices. It therefore applies to more
powerful binary command designs as well as the reciprocal protocol.
Its geometric construction has fixed periods and a fixed positive
patch margin, so period uncertainty is not the cause of the bound.

\begin{lemma}[A uniform physical packing]\label{lem:packing}
For $s=6+\beta$ there is a fixed class $\mathcal P\subset\mathcal B_\eta$
and constants $c_0,c_1>0$ such that, for all sufficiently small $h$,
$\mathcal P$ contains $2^m$ tables, where $m\ge c_0/h$, whose variable
component support functions have pairwise $C^2$ distance at least
$c_1h^{s-2}$. All have the same known primitive lattice and two
primitive component orbits. Their correspondences and pose can be
supplied without reducing this separation.
\end{lemma}
\begin{proof}
Fix the lattice $20\Z^2$. Put one variable component near the unit
disk at the origin, and a fixed disk of radius two at $(6,0)$, and
repeat both by this lattice. Choose a nonzero smooth function
$\psi$ supported in $(-1/3,1/3)$ with $\psi''(0)\ne0$.
Take angular centers $\theta_j$ in a fixed interval with mutual
spacing at least $h$, with $c/h\le m\le C/h$. For
$\omega\in\{0,1\}^m$, define the uncentered support of the variable
body by
\begin{equation}\label{eq:packing}
 p_\omega(\theta)=1+a h^s
                   \sum_{j=1}^m\omega_j\psi((\theta-\theta_j)/h),
\end{equation}
periodically extended, where $a>0$ is a fixed sufficiently small
number. At most one summand is nonzero at each point. Its derivative
of order $k\le6$ is bounded by $Ca h^{s-k}$. The sixth derivative
has uniformly bounded $\beta$-H\"older seminorm. Within a support
this follows by rescaling; between supports it follows from their
separation and the vanishing at their edges. The same estimate follows
by comparison to the intervening zero point for two arbitrarily close
points on different sides of an edge. Thus the $C^{6,\beta}$ bound
is uniform over all $h,m,\omega$.

For small $a$, $p_\omega+p_\omega''$ lies in a fixed positive interval.
The bodies are smooth and strictly convex, with all separation,
diameter, lattice and curvature bounds independent of $h$. Subtracting
their Steiner first harmonic also preserves the uniform support prior.
The two component species cannot be translated into each other:
their diameters lie in disjoint fixed intervals. Every period must
therefore send a variable component to another variable component,
forcing that period to lie in $20\Z^2$. Conversely every such vector
is a period. This proves primitiveness and the stated orbit count.

Center differences of the same species are true periods. For any
cross-species difference, the translated source species cannot match
a target of the other species with small support error, because their
widths differ by a fixed amount. A target of the same species has
center discrepancy bounded below by a fixed positive number: the
cross-species offset $(6,0)$, perturbed by the variable Steiner point,
stays a positive distance from $20\Z^2$. There are representatives
of both species in $Q_B$. These two alternatives give a uniform
positive defect for every nonperiod candidate in
\eqref{eq:patch-margin}. Hence one fixed $\eta>0$ works throughout
the construction. Taking the union of these tables over all small
$h$ defines a single class $\mathcal P$ with fixed numerical bounds.

If $\omega\ne\omega'$, choose an index where they differ and evaluate
the second derivative at its center. All other summands vanish there,
so the absolute difference is $a|\psi''(0)|h^{s-2}$. This proves the
separation in the laboratory support norm. The variable component,
lattice and fixed component may be identified in advance; none of
these disclosures identifies its bump vector. No analytic smoothness
bound is used or claimed for this packing.
\end{proof}

\begin{lemma}[Binary decision trees with randomized commands]
\label{lem:bits}
Let $M$ parameters have pairwise disjoint acceptable output sets.
A procedure with at most $N$ binary table-dependent outputs, arbitrary
adaptive inputs and parameter-independent randomization has average
success probability at most $2^N/M$ under the uniform distribution on
these parameters. The binary response probabilities may be arbitrary
functions of the parameter and chosen input.
\end{lemma}
\begin{proof}
Represent the controller randomization by an independent seed. Each
binary response can be generated by comparing another independent
uniform random variable with its conditional response probability.
This constructs the joint adaptive experiment for every parameter on
one common seed space. Fix all these independent seeds. Commands are
then functions only of the preceding binary transcript, and the output
is a function of a binary tree with at most $2^N$ leaves, after padding
short transcripts if necessary. An output at a leaf is acceptable for
at most one parameter, by disjointness. Thus at most $2^N$ parameters
can succeed for these fixed seeds. Average over the uniform parameter
and then integrate over the seeds. The bound follows. The numerical
values of selected commands give no additional leaves, since they are
already determined by the seed and transcript.
\end{proof}

\begin{proof}[Proof of Theorem~\ref{thm:lower}]
In Lemma~\ref{lem:packing} take
$h\asymp\nu^{1/(s-2)}$ with its constant large enough that the support
separation exceeds $2\nu$. The $\nu$-accuracy output sets are disjoint.
Lemma~\ref{lem:bits}, with $M=2^m$, gives average success at most
$2^{N-m}$. Uniform success at least $3/4$ therefore requires
$N\ge m+\log_2(3/4)\ge c\nu^{-1/(s-2)}$ for small $\nu$.
This argument allows more input control than the upper-bound procedure
and assumes perfect calibration. It remains valid for the latter
procedure with all-attempt bookkeeping. A protocol that reveals a
free-start test or a continuous impact location has a different output
alphabet and is not covered by the assertion.
\end{proof}

The loss in this theorem is the supplied laboratory-frame loss used by
the active controller. It is already enough to establish a matching
power for Theorem~\ref{thm:adaptive}, in the same loss. We do not silently
replace it by an unlabelled quotient loss in the lower bound, nor extend
a deterministic cap $N$ to an expected stopping-time budget without
another argument. No lower bound for the logarithms or for confidence
dependence is asserted.

For fixed confidence, on a nondegenerate bounded class containing
$\mathcal P$, the two results give
\begin{equation}\label{eq:matched-power}
 c\nu^{-1/(4+\beta)}\le N^*(\nu)
          \le C\nu^{-1/(4+\beta)}\log^2(C/\nu),
\end{equation}
where $N^*$ is the least uniform deterministic cap on attempted bits,
and constants in the accuracy definition are absorbed by rescaling
$\nu$. The matching power comes from one-dimensional boundary
smoothness measured two derivatives below the prior. The physical
work needed for the upper bound is that these boundary queries can
be implemented from the stipulated collision bits without additional
geometric outputs or uncharged free preparations.

\section{The lower-bound gauge and sequential observations}
\label{sec:sequential}
The deterministic-cap result above is retained. We now make its gauge
explicit and prove a separate result for expected stopping times. No
real-valued stopping signal or unobserved free-start test is allowed.

\subsection{A parameter-independent laboratory frame}
Let $n(\theta)=(\cos\theta,\sin\theta)$. For a laboratory support $p$
write
\begin{equation}\label{eq:centering-map}
 z(p)=\frac1\pi\int_0^{2\pi}p(\theta)n(\theta)\dd\theta,
 \qquad p^0=p-z(p)\cdot n.
\end{equation}
The disclosed information in Lemma~\ref{lem:packing} consists of the
fixed laboratory axes, $20\Z^2$, the names of the two component species,
and the fixed radius-two disk at $(6,0)$. These data are identical for
all packing parameters. In particular, the unknown Steiner point of
the variable body is not disclosed. ``Pose given'' does not mean a
parameter-dependent translation or rotation of that body.

For the variable component define acceptable output sets by
\begin{equation}\label{eq:accuracy-sets}
 \mathcal U_\omega(\nu)=
 \{\widehat C:\ \|p^0_{\widehat C}-p^0_\omega\|_{C^2}\le\nu\}.
\end{equation}
This centered-shape loss alone is enough for the lower bound. The
stronger laboratory loss may also include
$|z_{\widehat C}-z(p_\omega)|$ and errors of the fixed component.
No independent rotation of each output is taken in either loss.

\begin{lemma}[The packing survives centering]\label{lem:centered-packing}
For the family \eqref{eq:packing}, distinct parameters satisfy
\[
 \|p^0_\omega-p^0_{\omega'}\|_{C^2}\ge c h^{s-2}
\]
for all sufficiently small $h$, uniformly over the packing. The
centered physical support bounds and the patch margin remain uniform.
\end{lemma}
\begin{proof}
The disjoint angular supports give
$\|p_\omega-p_{\omega'}\|_\infty\le C a h^s$, independently of
$m$. Equation~\eqref{eq:centering-map} implies
$|z(p_\omega)-z(p_{\omega'})|\le2Ca h^s$. The $C^2$ norm of the
subtracted first harmonic is therefore $O(a h^s)$. At a differing
bump center the uncentered second-derivative difference has magnitude
$a|\psi''(0)|h^{s-2}$. Since $h^s=o(h^{s-2})$, at least half this
difference survives centering. A first harmonic is bounded in every
fixed smoothness norm by its amplitude, so the centered prior is also
preserved. The separation and period arguments of
Lemma~\ref{lem:packing} already allow these $O(a h^s)$ center shifts.
Thus the disclosed data contain no shape-dependent centering information.
\end{proof}

\subsection{A stopped transcript carries at most its expected length}
A stopping rule is admissible when, conditional on parameter-independent
controller randomness, stopping, further commands and the eventual
output depend only on the preceding bits. Waiting times or a physical
success flag are not additional data. The number $\mathsf T$ counts
all attempted preparations, including zeros from solid starts.

\begin{lemma}[Sequential binary information bound]\label{lem:sequential}
Let $M\ge2$ parameters have disjoint acceptable output sets. If an
admissible randomized procedure succeeds with probability at least
$1-\delta$ for each parameter, $0<\delta<1/2$, then
\begin{equation}\label{eq:sequential-information}
 \frac1M\sum_{j=1}^M\E_j\mathsf T
 \ge (1-\delta)\log_2M-h_2(\delta),
\end{equation}
where $h_2(x)=-x\log_2x-(1-x)\log_2(1-x)$. There is no deterministic
cap in this assertion. An infinite expected length satisfies it
trivially.
\end{lemma}
\begin{proof}
Give the parameter $V$ the uniform distribution and denote the
independent controller seed by $U$. Suppose the average expected length
is finite, so the terminal binary word $Z$ is finite almost surely.
For each fixed $U=u$, terminal words are prefix-free: after the rule
stops at one word it cannot also stop at a proper extension. Thus
$\sum_z2^{-|z|}\le1$. If this sum is $K_u>0$, comparison with the
probability distribution $2^{-|z|}/K_u$ by nonnegativity of relative
entropy gives
\[
 H(Z\mid U=u)\le\E(|Z|\mid U=u)+\log_2K_u
                    \le\E(\mathsf T\mid U=u).
\]
The countable case follows from the same log-sum inequality by
truncation. Integrate in $u$. Since $U$ is independent of $V$,
\[
 I(V;Z,U)=I(V;Z\mid U)\le H(Z\mid U)\le\E\mathsf T.
\]
Decode the parameter from the unique acceptable output set, using an
arbitrary index on failure. Its average error $p_e$ is at most $\delta$.
For the error indicator $E$, the elementary conditional-entropy bound is
\[
 H(V\mid Z,U)\le H(E\mid Z,U)+p_e\log_2(M-1)
 \le h_2(\delta)+\delta\log_2M.
\]
Subtract from $H(V)=\log_2M$ and use the preceding information bound.
The output randomness may be included in $U$ at the outset. Neither
continuous input values nor the word length yield an extra channel:
they are already determined by the seed and stopped word. The coding
and entropy principles used here are classical \cite{Shannon}; the
argument is included to specify the stopping-time scope exactly.
\end{proof}

\begin{theorem}[Expected attempted-bit lower bound]\label{thm:expected}
On the fixed physical class of Lemma~\ref{lem:packing}, any admissible
sequential procedure with success probability at least $3/4$ uniformly
in the centered loss \eqref{eq:accuracy-sets} satisfies
\begin{equation}\label{eq:expected-lower}
                 \sup_{\mathcal O}\E_{\mathcal O}\mathsf T
                         \ge c\nu^{-1/(s-2)}.
\end{equation}
The bound holds with exact implementation and arbitrary-precision
adaptive inputs, even with the fixed lattice, species and laboratory
frame disclosed. It therefore also holds for the stronger physical
reconstruction loss.
\end{theorem}
\begin{proof}
Choose $h\asymp\nu^{1/(s-2)}$ so that Lemma~\ref{lem:centered-packing}
gives separation greater than $2\nu$. Apply Lemma~\ref{lem:sequential}
with $M=2^m$ and $m\ge c/h$. The supremum dominates the average.
\end{proof}

At fixed confidence, the least uniform expected number of attempts on
any fixed class containing this packing lies between
$c\nu^{-1/(s-2)}$ and
$C\nu^{-1/(s-2)}\log^2(C/\nu)$, since the deterministic upper bound
is also an expected upper bound. The analogous deterministic-cap
statement remains \eqref{eq:matched-power}. These are two explicitly
proved resource criteria, not an identification of random length with
a deterministic cap. No sharp logarithmic or confidence factor follows.

\section{A resolution converse for table-dependent bounded errors}
\label{sec:calibration-floor}
The sufficient calibration scale in \eqref{eq:precision} has a converse
in the same bounded-error model. The converse uses physical start
perturbations, not corrupted labels added to an abstract membership
oracle. It concerns short collision commands; arbitrarily long flights
or additional observations of the actual launch position are not included.

\subsection{A common response produced by nearby physical tables}
For a table $\mathcal O_j$ denote its attempted hit bit by $B^j_a(q)$.
An admissible position-error implementation is a measurable map
$q\mapsto F_j(q,a)$ with $|F_j(q,a)-q|\le r$; time and direction
can remain exact. Such maps may depend on the table and the nominal
command. The controller sees only the resulting bit, not the actual
start. This is a submodel of the conditionally bounded, not necessarily
mean-zero couplings in Section~\ref{sec:setting}.

\begin{lemma}[Pointwise common response under two implementations]\label{lem:common-response}
Suppose $\mathcal O_0,\mathcal O_1$ consist of matched strictly convex
components of diameter at most $D_0$, with inradius at least $r_*>0$,
and that matched bodies have Hausdorff distance at most $e>0$.
In each table distinct bodies are separated by at least $d_0$.
Fix $t_{\max}<d_0$. If $e$ is sufficiently small compared with
$r_*$ and $d_0-t_{\max}$, there are measurable position perturbations
of size at most $2e$ such that, for every $|a|\le t_{\max}$ and
with closed-solid and swept-set conventions,
\begin{equation}\label{eq:common-response}
 B^0_a(F_0(q,a))=B^1_a(F_1(q,a))
                    =\min\{B^0_a(q),B^1_a(q)\}.
\end{equation}
The closed convention assigns zero on a solid boundary and one on a
swept boundary outside the solid. In every disagreement case the
forced-zero position is moved strictly into a solid or strictly outside
all sweeps before its response is evaluated. Agreement cases keep the
nominal point and use the same closed convention. The conclusion
applies to locally finite periodic arrays and to both the forward and
translated reverse commands.
\end{lemma}
\begin{proof}
For one fixed $a$ put $S_j(C)=C+[-a,0]$. Matched swept bodies have
Hausdorff distance at most $e$, because their support functions differ
by the same amount as those of the matched bodies. In either table,
\begin{equation}\label{eq:swept-separation}
 \dist(S_j(C),S_j(D))\ge d_0-|a|\quad(C\ne D).
\end{equation}
Indeed two points in the sweeps are $c-ua$ and $d-va$, with
$u,v\in[0,1]$. Their distance is at least $|c-d|-|u-v||a|$.
Thus a start belonging to a swept body determines that body uniquely.
Moreover no solid other than that body meets this sweep. On it the
hit bit is the indicator of the sweep minus the indicator of its body.

Leave the start unchanged when the two bits agree, and also on the
side whose bit is already zero. Consider $B^0_a(q)=1$, $B^1_a(q)=0$,
and let $C_0$ be the unique body whose sweep contains $q$. Write $C_1$
for its matched body. The separation \eqref{eq:swept-separation} and
Hausdorff matching imply the following exhaustive alternatives:
\[
                   q\in C_1,\qquad\hbox{or}\qquad q\notin\operatorname{int}S_1(C_1).
\]
To check this, every other sweep in table 1 is at distance at least
$d_0-|a|-e$ from $S_0(C_0)$, so cannot supply either a competing hit
or a solid start there. If $q\in S_1(C_1)\setminus C_1$, its bit
would be one.

In the first case $\dist(q,C_0)\le e$. Project $q$ onto $C_0$, and
move the projection a further distance less than $e/2$ toward a fixed
interior point of $C_0$. This produces an interior solid start within
$3e/2$ of $q$, whose bit is zero. Uniform inballs permit a fixed small
convex-combination factor. The construction is measurable.

In the second case, when $q$ is outside the closed sweep, let $q_1$
be its metric projection onto $S_1(C_1)$ and set
$v=(q-q_1)/|q-q_1|$. At a boundary point choose an outward unit normal
from its compact normal cone, measurably, for example by the least
angle in a fixed angular convention. In both cases
$q\cdot v\ge p_{S_1(C_1)}(v)$, and support matching gives
$p_{S_0(C_0)}(v)\le p_{S_1(C_1)}(v)+e$. Hence
\[
                       q'=q+\tfrac32 e v
\]
lies strictly outside $S_0(C_0)$. It cannot enter another swept body
of table 0, by \eqref{eq:swept-separation} and smallness of $e$.
It is therefore a free miss, with bit zero. The displacement is $3e/2$.
Projection onto a closed convex set is continuous, so this map is
measurable on its case domain. Reverse the indices when the differing
bit is on side 1. Local finiteness and the separation give measurable
component selection. For $a=0$ every bit is zero and no change is needed.
Boundary cases can be assigned by the same inward or outward moves;
in particular the final forced-zero positions are not on boundaries.

The proof is pointwise in the nominal start and displacement. Apply it
to $(q,a)$ for a forward command and to $(q+a,-a)$ for its reverse.
Fresh nominal draws and deterministic perturbation maps preserve
conditional independence. They do not supply another observable.
\end{proof}

\begin{theorem}[Minimax resolution under table-dependent bounded errors]
\label{thm:calibration-floor}
For each fixed $s=6+\beta$ there is a fixed physical subclass of
$\mathcal B_\eta$, with known primitive lattice, two known species
and the fixed laboratory frame, and constants $c,r_0>0$ such that
for every $0<r<r_0$ two tables in this class have centered-support
$C^2$ distance at least
\begin{equation}\label{eq:calibration-separation}
                            c r^{(s-2)/s}
\end{equation}
yet admit position-error implementations bounded by $r$ with identical
laws for the entire adaptive bit experiment. This holds for every
nominal command design with $|a|\le t_{\max}<d_0$, arbitrary input
precision and arbitrary numbers of attempts. A sequential stopping
rule based on the bits cannot distinguish them either.

To specify the uncertainty quantifiers, let $\mathfrak I_r(\mathcal O)$
be the measurable implementation maps permitted above, and let
$\mathcal A_{t_{\max}}$ be the admissible bit-based controllers. The
construction has the stronger, controller-uniform form
\begin{gather}\label{eq:calibration-quantifiers}
 \forall\,0<r<r_0\quad
 \exists\,(\mathcal O_0,F_0),(\mathcal O_1,F_1),\quad
             F_j\in\mathfrak I_r(\mathcal O_j),\notag\\
 \forall\,\mathcal A\in\mathcal A_{t_{\max}}:\qquad
 \mathsf{Law}_{\mathcal O_0,F_0}^{\mathcal A}(Z)
       =\mathsf{Law}_{\mathcal O_1,F_1}^{\mathcal A}(Z).
\end{gather}
Here $Z$ includes the returned bits and the controller's recorded
commands and stopping decisions, not the actual launch positions.
The maps $F_0,F_1$ need not agree and may use the nominal draw and
command. They are fixed measurable maps for the two worlds, not an
extra observation or a common known jitter distribution.

Consequently, uniform success probability greater than $1/2$ at
centered $C^2$ accuracy $\nu$ under all such position errors requires
\begin{equation}\label{eq:calibration-necessary}
                            r\le C\nu^{s/(s-2)}.
\end{equation}
Constants, not an equality threshold, are asserted. The result is a
worst-case bounded-calibration statement; it does not assert this
obstruction for one fixed common offset or for a known mean-zero law.
\end{theorem}
\begin{proof}
Use the two-species periodic construction in Lemma~\ref{lem:packing},
but compare the unit-disk support $p_0=1$ with a single nonzero bump
\[
 p_1(\theta)=1+a h^s\psi((\theta-\theta_0)/h).
\]
The amplitude $a$ is fixed small enough that all curvature, separation
and patch margins remain uniform. Choose
$h=(r/(8a\|\psi\|_\infty))^{1/s}$. Then the Hausdorff distance of
the two variable bodies is at most $e=r/8$. The fixed disk and all
periodic copies are matched without alteration. At the bump center
the second derivatives differ by $a|\psi''(0)|h^{s-2}$; the
first-harmonic estimate in Lemma~\ref{lem:centered-packing} shows that
centering removes only $O(h^s)$. This gives
\eqref{eq:calibration-separation}. The lattice and disclosed data are
identical for both tables; the bump does not change primitiveness.

Lemma~\ref{lem:common-response} produces two physical implementations
with position errors at most $2e<r$. Couple the parameter-independent
controller randomness and every nominal draw across the two worlds.
Their first bits are equal by \eqref{eq:common-response}. Inductively,
identical histories generate identical next commands and identical
bits. Thus the complete bit streams, nominal controller records and
all bit-measurable stopping decisions have identical laws. No finite
or infinite sequence of those observations distinguishes the worlds.
An output law common to two disjoint accuracy sets cannot place
probability greater than $1/2$ in each. If $2\nu$ is smaller than
\eqref{eq:calibration-separation}, uniform success greater than $1/2$
is impossible. Rearranging proves \eqref{eq:calibration-necessary}.
\end{proof}

\subsection{The two resource requirements}
Together with Theorems~\ref{thm:adaptive} and \ref{thm:expected}, this
gives two necessary resources on one fixed nondegenerate subclass:
\[
 \sup_{\mathcal O}\E_{\mathcal O}\mathsf T
       \ge c\nu^{-1/(s-2)},\qquad
 r\le C\nu^{s/(s-2)}.
\]
Conversely the reciprocal compass construction attains the geometric
accuracy with the stated logarithmic overhead in attempts when
$\ell+\tau+t\alpha\le c\nu^{s/(s-2)}$. The converse calibration
necessity already holds with duration and angle exact, so it concerns
position error alone, not separate optimal exponents for all three
actuators. The lower bound allows any directions of length at most
$t_{\max}$, a larger design class than the four-atom upper construction.
The finite digital realization in Section~\ref{sec:digital} does not
alter either bound.

Thus the shrinking physical tolerance is not merely an omitted term in
a sample count: under the stipulated worst-case coupling it cannot be
replaced by additional samples. This is a statement within the active
collision model. Periodicity, the bounded presentation and the known
patch margin are still priors for the uniform global conclusion.
Neither crystallinity, apparatus calibration nor passive spectral
rigidity is inferred from the returned bits.

\subsection*{Source history and assisted analysis}
The article contains the proofs required for its statements. All
preceding localized, finite-control, stationary, calibration, global
rigidity and statistical arguments remain active in the present source.
Their historical manuscript and review directories are unchanged; they
are not independently published articles. AI-assisted analysis
contributed to the collision-query, footprint, global-response,
registration, statistical, finite-stopping, isotropic and single-law
directional-germ arguments and
the source checks. The named author is responsible for the proofs and
references. Finite diagnostics and a successful build are reproducibility
evidence, not certificates of continuum proofs or physical apparatus.
\begin{thebibliography}{99}
\raggedright
\bibitem{AFP}
L. Ambrosio, N. Fusco, and D. Pallara, \emph{Functions of Bounded
Variation and Free Discontinuity Problems}, Oxford Mathematical
Monographs, Clarendon Press, Oxford, 2000.
\bibitem{Altman}
E. Altman, \emph{Constrained Markov Decision Processes},
Chapman \& Hall/CRC, Boca Raton, 1999.
\bibitem{AverkovBianchi}
G. Averkov and G. Bianchi,
\emph{Confirmation of Matheron's conjecture on the covariogram of a planar convex body},
J. Eur. Math. Soc. \textbf{11} (2009), 1187--1202.
\href{https://doi.org/10.4171/JEMS/179}{doi:10.4171/JEMS/179}.
\bibitem{BianchiCross}
G. Bianchi, \emph{The cross covariogram of a pair of polygons determines both
polygons, with a few exceptions}, Adv. Appl. Math. \textbf{42} (2009), 519--544.
\href{https://doi.org/10.1016/j.aam.2008.10.002}{doi:10.1016/j.aam.2008.10.002}.
\bibitem{BianchiLaplace}
G. Bianchi, \emph{The covariogram and Fourier--Laplace transform in $\mathbb C^n$},
Proc. Lond. Math. Soc. (3) \textbf{113} (2016), 1--23.
\href{https://doi.org/10.1112/plms/pdw020}{doi:10.1112/plms/pdw020}.
\bibitem{BGK}
G. Bianchi, R. J. Gardner, and M. Kiderlen,
\emph{Phase retrieval for characteristic functions of convex bodies and
reconstruction from covariograms}, J. Amer. Math. Soc. \textbf{24} (2011), 293--343.
\href{https://doi.org/10.1090/S0894-0347-2010-00683-2}
{doi:10.1090/S0894-0347-2010-00683-2}.
\bibitem{BrunelCaps}
V.-E. Brunel, \emph{Uniform behaviors of random polytopes under the
Hausdorff metric}, Bernoulli \textbf{25} (2019), 1770--1793.
\href{https://doi.org/10.3150/18-BEJ1035}{doi:10.3150/18-BEJ1035}.
\bibitem{BKY}
V.-E. Brunel, J. M. Klusowski, and D. Yang,
\emph{Estimation of convex supports from noisy measurements},
Bernoulli \textbf{27} (2021), 772--793.
\href{https://doi.org/10.3150/20-BEJ1229}{doi:10.3150/20-BEJ1229}.
\bibitem{CGLSupport}
J. Capitao-Miniconi, \'E. Gassiat, and L. Leh\'ericy,
\emph{Support and distribution inference from noisy data}, Bernoulli, forthcoming.
\href{https://arxiv.org/abs/2304.09452}{arXiv:2304.09452}.
Theorem numbering follows the
\href{https://www.e-publications.org/ims/submission/BEJ/user/submissionFile/68730?confirm=a078ae36}
{accepted manuscript}, accessed October 4, 2026.
\bibitem{CN}
R. M. Castro and R. D. Nowak, \emph{Minimax bounds for active learning},
IEEE Trans. Inform. Theory \textbf{54} (2008), 2339--2353.
\href{https://doi.org/10.1109/TIT.2008.920189}{doi:10.1109/TIT.2008.920189}.
\bibitem{DelaigleMeister}
A. Delaigle and A. Meister,
\emph{Nonparametric function estimation under Fourier-oscillating noise},
Statist. Sinica \textbf{21} (2011), 1065--1092.
\href{https://doi.org/10.5705/ss.2009.082}{doi:10.5705/ss.2009.082}.
\bibitem{Galerne}
B. Galerne, \emph{Computation of the perimeter of measurable sets via
their covariogram. Applications to random sets},
Image Anal. Stereol. \textbf{30} (2011), 39--51.
\href{https://doi.org/10.5566/ias.v30.p39-51}
{doi:10.5566/ias.v30.p39-51}.
\bibitem{GLL}
\'E. Gassiat, S. Le Corff, and L. Leh\'ericy,
\emph{Deconvolution with unknown noise distribution is possible for multivariate
signals}, Ann. Statist. \textbf{50} (2022), 303--323.
\href{https://doi.org/10.1214/21-AOS2106}{doi:10.1214/21-AOS2106}.
\bibitem{GK}
P. W. Goldberg and S. Kwek,
\emph{The precision of query points as a resource for learning convex
polytopes with membership queries},
Proceedings of the Thirteenth Annual Conference on Computational
Learning Theory, Morgan Kaufmann, 2000, 225--235.
\bibitem{KiderlenRataj}
M. Kiderlen and J. Rataj, \emph{Dilation volumes of sets of finite perimeter},
Adv. Appl. Probab. \textbf{50} (2018), 1095--1118.
\href{https://doi.org/10.1017/apr.2018.52}{doi:10.1017/apr.2018.52}.
\bibitem{LawlerLimic}
G. F. Lawler and V. Limic, \emph{Random Walk: A Modern Introduction},
Cambridge Studies in Advanced Mathematics \textbf{123}, Cambridge
University Press, 2010.
\bibitem{LCK}
A. Locatelli, A. Carpentier, and S. Kpotufe,
\emph{An adaptive strategy for active learning with smooth decision boundary},
Proceedings of Algorithmic Learning Theory, Proceedings of Machine
Learning Research \textbf{83} (2018), 547--571.
\bibitem{LRSS}
A. K. Louis, M. Riplinger, M. Spiess, and E. Spodarev,
\emph{Inversion algorithms for the spherical Radon and cosine transform},
Inverse Problems \textbf{27} (2011), 035015.
\href{https://doi.org/10.1088/0266-5611/27/3/035015}
{doi:10.1088/0266-5611/27/3/035015}.
\bibitem{MartinezMaure}
Y. Martinez-Maure, \emph{Geometric study of Minkowski differences of
plane convex bodies}, Canad. J. Math. \textbf{58} (2006), 600--624.
\href{https://doi.org/10.4153/CJM-2006-025-x}{doi:10.4153/CJM-2006-025-x}.
\bibitem{MeisterCompact}
A. Meister, \emph{Deconvolving compactly supported densities},
Math. Methods Statist. \textbf{16} (2007), 63--76.
\href{https://doi.org/10.3103/S106653070701005X}{doi:10.3103/S106653070701005X}.
\bibitem{Retained}
Q. Qi, \emph{Scalar collision laws and recognition of periodic dispersing
billiards}, A2 v31, October 4, 2026. Complete reviewed manuscript retained in the repository source archive;
not an independently published article or a required supplementary volume.
\bibitem{RigolletWeed}
P. Rigollet and J. Weed, \emph{Uncoupled isotonic regression via minimum
Wasserstein deconvolution}, Inform. Inference \textbf{8} (2019), 691--717.
\href{https://doi.org/10.1093/imaiai/iaz006}{doi:10.1093/imaiai/iaz006}.
\bibitem{Schneider}
R. Schneider, \emph{Convex Bodies: The Brunn--Minkowski Theory},
second expanded edition, Encyclopedia of Mathematics and its Applications
\textbf{151}, Cambridge University Press, 2014.
\bibitem{Shannon}
C. E. Shannon, \emph{A mathematical theory of communication},
Bell System Tech. J. \textbf{27} (1948), 379--423, 623--656.
\bibitem{Villarrubia}
J. S. Villarrubia, \emph{Algorithms for scanned probe microscope image
simulation, surface reconstruction, and tip estimation},
J. Res. Natl. Inst. Stand. Technol. \textbf{102} (1997), 425--454.
\href{https://doi.org/10.6028/jres.102.030}{doi:10.6028/jres.102.030}.
\bibitem{YCJ}
S. Yan, K. Chaudhuri, and T. Javidi,
\emph{Active learning from imperfect labelers},
Advances in Neural Information Processing Systems \textbf{29} (2016).
\end{thebibliography}

\end{document}
